% !TeX program = xelatex
% 数学 LaTeX 排版 Skill v4.0.0；版式保持 v3 金标准. XeLaTeX 编译至少两次. 不分发字体.
% WZ-LOCKED-CLASS-BEGIN
\begin{filecontents*}[overwrite]{elegantbook-original-adapter.cls}
\NeedsTeXFormat{LaTeX2e}
\ProvidesClass{elegantbook-original-adapter}[2026/09/23 v3.0 fixed mathematical book environments]
\LoadClass[10pt,letterpaper,oneside]{book}
\RequirePackage[UTF8,scheme=plain,fontset=fandol]{ctex}
\RequirePackage{fontspec}
\setmainfont{cmunrm.otf}[BoldFont=cmunbx.otf,ItalicFont=cmunti.otf,BoldItalicFont=cmunbi.otf]
\setCJKfamilyfont{sourceheader}[ItalicFont=FandolKai-Regular.otf]{FandolKai-Regular.otf}
\RequirePackage{amsmath,amssymb,mathrsfs}
\RequirePackage{amsthm,etoolbox}
\RequirePackage{xcolor,geometry,setspace,fancyhdr,enumitem,needspace,xurl}
\definecolor{structurecolor}{HTML}{880E4F}
\definecolor{main}{HTML}{9C27B0}
\geometry{paperwidth=612bp,paperheight=792bp,left=20mm,right=20mm,top=95.04bp,bottom=20mm,headheight=12pt,headsep=25pt,footskip=30pt}
\setstretch{1.6}
\setlength{\parindent}{2em}
\setlength{\parskip}{0pt}
\setlist{nosep,leftmargin=2em}
\AtBeginDocument{\setlength{\abovedisplayskip}{4pt plus 1pt minus 1pt}\setlength{\belowdisplayskip}{4pt plus 1pt minus 1pt}\setlength{\abovedisplayshortskip}{2pt}\setlength{\belowdisplayshortskip}{3pt}}
\fancyhf{}
\fancyhead[L]{{\itshape\CJKfamily{sourceheader}\rightmark}}
\fancyhead[R]{{\itshape\CJKfamily{sourceheader}\leftmark}}
\fancyfoot[C]{\thepage}
\renewcommand{\headrulewidth}{0.4pt}
\renewcommand{\headrule}{{\color{structurecolor}\hrule height\headrulewidth width\headwidth\vskip-\headrulewidth}}
\pagestyle{fancy}
\fancypagestyle{cover}{\fancyhf{}\fancyfoot[C]{\thepage}\renewcommand{\headrulewidth}{0pt}\renewcommand{\headrule}{}}
\RequirePackage{hyperref}
\hypersetup{unicode=true,colorlinks=true,linkcolor=structurecolor,urlcolor=structurecolor,citecolor=structurecolor,linktoc=section,bookmarksnumbered=true,bookmarksopen=true,pdfborder={0 0 0},pdfstartview=FitH}
\setcounter{tocdepth}{1}
\renewcommand*\l@section{\@dottedtocline{1}{1.5em}{2.4em}}
\renewcommand{\thesection}{\thechapter.\arabic{section}}
\newcommand{\originalchapter}[1]{\par\addvspace{14pt}\Needspace{5\baselineskip}\refstepcounter{chapter}\setcounter{section}{0}\addcontentsline{toc}{chapter}{\protect\numberline{\thechapter}#1}\markboth{\thechapter\quad #1}{}\begingroup\centering\color{structurecolor}\bfseries\fontsize{14.4}{20}\selectfont\thechapter\quad #1\par\endgroup\nobreak\vspace{8pt}}
\newcommand{\originalsection}[1]{\par\addvspace{9pt}\Needspace{4\baselineskip}\refstepcounter{section}\addcontentsline{toc}{section}{\protect\hspace*{1.5em}\protect\numberline{\protect\textcolor{structurecolor}{\thesection}}#1}\markright{\thesection\quad #1}\noindent{\fontsize{12}{17}\selectfont\color{structurecolor}\thesection\quad}{\bfseries\fontsize{12}{17}\selectfont #1}\par\nobreak\vspace{4pt}}

% Semantic environments are registered once here, not re-designed in manuscripts.
\newtheoremstyle{wzstatement}{6pt}{5pt}
  {\normalfont\kaishu}{0pt}{\normalfont\bfseries\color{structurecolor}}
  {}{.7em}{\thmname{#1}\thmnumber{ #2}\thmnote{\normalfont\ (#3)}}
\newtheoremstyle{wzdefinition}{6pt}{5pt}
  {\normalfont\songti}{0pt}{\normalfont\bfseries\color{structurecolor}}
  {}{.7em}{\thmname{#1}\thmnumber{ #2}\thmnote{\normalfont\ (#3)}}
\newtheoremstyle{wzproblem}{7pt}{5pt}
  {\normalfont\songti}{2em}{\normalfont\bfseries\color{main}}
  {}{.7em}{\thmname{#1}\thmnumber{ #2}}
\newtheoremstyle{wzremark}{4pt}{4pt}
  {\normalfont\songti}{0pt}{\normalfont\bfseries}
  {}{.7em}{\thmname{#1}}
\theoremstyle{wzstatement}
\newtheorem{theorem}{定理}[chapter]
\newtheorem{lemma}[theorem]{引理}
\newtheorem{proposition}[theorem]{命题}
\newtheorem{corollary}[theorem]{推论}
\theoremstyle{wzdefinition}
\newtheorem{definition}[theorem]{定义}
\theoremstyle{wzproblem}
\newtheorem{example}{例题}[chapter]
\newtheorem{exercise}{习题}[chapter]
\theoremstyle{wzremark}
\newtheorem*{remark}{注}
\renewcommand{\proofname}{证明}
% Retain the native amsthm QED stack; only adjust the fixed typography.
\expandafter\patchcmd\csname\string\proof\endcsname{\normalfont \topsep6\p@\@plus6\p@\relax}
  {\normalfont\songti \topsep5\p@\@plus1\p@\relax}
  {}{\ClassError{elegantbook-original-adapter}{Cannot patch proof spacing}{Check the amsthm version.}}
\expandafter\patchcmd\csname\string\proof\endcsname{\itshape}{\normalfont\bfseries}
  {}{\ClassError{elegantbook-original-adapter}{Cannot patch proof heading}{Check the amsthm version.}}
\expandafter\patchcmd\csname\string\proof\endcsname{\ignorespaces}
  {\setlength{\parindent}{2em}\ignorespaces}
  {}{\ClassError{elegantbook-original-adapter}{Cannot patch proof paragraphs}{Check the amsthm version.}}
\AtBeginEnvironment{proof}{\par\Needspace{3\baselineskip}}
\renewcommand{\qedsymbol}{\ensuremath{\square}}
\newenvironment{solution}{\begin{proof}[解答]}{\end{proof}}
\newcommand{\wzcheckchapter}{\ifnum\value{chapter}<1
  \ClassError{elegantbook-original-adapter}{Numbered mathematics before chapter 1}{Move it after the first originalchapter.}\fi}

\newcommand{\wzregisterguard}[1]{\AtBeginEnvironment{#1}{\wzcheckchapter\par\Needspace{3\baselineskip}}}
\forcsvlist{\wzregisterguard}{theorem,lemma,proposition,corollary,definition,example,exercise}
% Front matter and bibliography are not numbered mathematical chapters.
\newcommand{\originalpreface}{\par\addvspace{10pt}\Needspace{5\baselineskip}
  \noindent{\bfseries\fontsize{12}{17}\selectfont 前言}\par\nobreak\vspace{4pt}}
\newcommand{\originalcontents}{\par\addvspace{14pt}
  \begin{center}{\color{structurecolor}\bfseries\fontsize{14.4}{20}\selectfont 目录}\end{center}
  \vspace{-10pt}\@starttoc{toc}}
\renewcommand{\bibname}{参考文献}
\renewenvironment{thebibliography}[1]{
  \par\addvspace{12pt}\Needspace{3\baselineskip}\phantomsection
  \addcontentsline{toc}{chapter}{\bibname}
  \markboth{\bibname}{}
  \noindent{\color{structurecolor}\bfseries\fontsize{14.4}{20}\selectfont\bibname}\par\nobreak\vspace{6pt}
  \list{\@biblabel{\@arabic\c@enumiv}}
    {\settowidth\labelwidth{\@biblabel{#1}}\leftmargin\labelwidth
     \advance\leftmargin\labelsep\usecounter{enumiv}
     \let\p@enumiv\@empty\renewcommand\theenumiv{\@arabic\c@enumiv}
     \setlength{\itemsep}{3pt}\setlength{\parsep}{0pt}}
  \sloppy\clubpenalty4000\widowpenalty4000\sfcode`\.=1000\relax
}{\def\@noitemerr{\@latex@warning{Empty bibliography}}\endlist}

\numberwithin{equation}{chapter}
\renewcommand{\theequation}{\thechapter.\arabic{equation}}
\allowdisplaybreaks[2]
\emergencystretch=2em
\clubpenalty=6000
\widowpenalty=6000
\raggedbottom
\endinput
\end{filecontents*}
% WZ-LOCKED-CLASS-END
\documentclass{elegantbook-original-adapter}
% ===== 元数据内容槽 =====
\newcommand{\BookTitle}{数值线性代数}
\newcommand{\BookAuthor}{}
\newcommand{\BookDate}{}
\newcommand{\PDFTitle}{数值线性代数：MIT 18.335J 中文整理与补充}
\newcommand{\PDFAuthor}{}
\hypersetup{pdftitle={\PDFTitle},pdfauthor={\PDFAuthor}}
\usepackage{tikz}
\usetikzlibrary{arrows.meta,positioning,calc}
\usepackage{longtable,booktabs}
\newcommand{\dd}{\,\mathrm d}
\newcommand{\R}{\mathbb R}
\newcommand{\C}{\mathbb C}
\newcommand{\norm}[1]{\lVert #1\rVert}
\newcommand{\inner}[2]{\langle #1,#2\rangle}
\newcommand{\eps}{\varepsilon}
\DeclareMathOperator{\diag}{diag}
\DeclareMathOperator{\rank}{rank}
\DeclareMathOperator{\range}{range}
\DeclareMathOperator{\Span}{span}
\begin{document}
\frontmatter
\begin{center}
{\color{structurecolor}\bfseries\fontsize{18}{25}\selectfont 数值线性代数\par}
\vspace{5pt}
{MIT 18.335J Spring 2019 相关课程资料中文整理与自编补充\par}
\end{center}
\originalpreface
数值线性代数研究如何在有限精度和有限计算资源下求解线性方程组、最小二乘问题及特征值问题. 本册面向已学过矩阵运算、线性空间、内积、特征值和微积分的读者, 目标是把数学结构、误差分析和实际算法连在一起. 从浮点运算出发, 依次讨论条件数、SVD、QR、LU、Cholesky、特征值算法和大型稀疏问题. 读完后应能解释为什么某个算法适合某类矩阵, 能检查其条件, 也能用残差判断计算结果的可信程度.

主要来源是 MIT OpenCourseWare 的 Steven G. Johnson 教授 18.335J Introduction to Numerical Methods, Spring 2019. 本册选取第 2--5、7--24 讲的数值线性代数主题; 原课程还含 ODE、优化、积分和 FFT, 本册不重复这些专题. 官方资源栏共有 27 份 PDF 手册, 共 261 页; 其中直接相关的核心手册为 12 份、59 页, 性能补充 3 份、16 页. 课程另有逐周 HTML 课堂说明, 部分重要讲次明确没有独立手册. 因此本册是按学科重组的中文讲义, 不是声称完整逐字翻译了原课程 39 讲的官方讲稿.

有独立手册的部分按其数学内容整理, 保留原作者和来源; 原手册简略、仅列算法或只有课堂主题的部分, 补充了连续的中文解释、证明、例题与习题. 这些扩展是本册自编补充, 不冒称 MIT 原文. 每章末说明对应课次和资料, 附录给出来源覆盖表. Per-Olof Persson 的浮点、正交化、QR 和稀疏算法手册另行署名. 本册没有复制或下载 Trefethen 与 Bau 的版权教材全文.

MIT OCW 资料一般按 Creative Commons Attribution--NonCommercial--ShareAlike 4.0 许可使用, 具体资源的特殊署名和权利说明仍适用. 本册作为非商业学习用改编材料保留相同许可, 不代表 MIT 官方出版或认可. 原件、逐文件来源和校验值另随项目保存. 阅读时应把数学结论、精确算术中的算法性质、有限精度中的经验现象分开; 数值实验用于检验实现和理解现象, 不能替代证明.
\originalcontents
\clearpage
\mainmatter
% SOURCE: OCW L2--4, Persson IEEE handout; Johnson myths, summation, recursive summation, norm equivalence.
% ADAPTATION: original connected Chinese exposition, proofs and worked examples added.
\originalchapter{浮点运算与误差分析}\label{ch:float}
\originalsection{为什么精确公式仍会算错}
线性代数中的等式是关于实数或复数的精确关系. 计算机只存有限个数, 每次运算还要把结果舍入到可表示集合. 例如在精确算术中, $((a+b)-a)-b=0$; 若 $a$ 很大、$b$ 很小, $a+b$ 可能舍入为 $a$, 随后的结果就成为 $-b$. 因此代数上相同的公式, 可能具有不同的数值行为. 我们要分析的是实际运算路径, 不能只看最终表达式.

\begin{definition}
基数为 $\beta\ge2$、有效位数为 $t$ 的正规浮点数写为 $\pm m\beta^e$, 其中有效数字满足 $1\le m<\beta$, 且 $m$ 由 $t$ 位基数 $\beta$ 数字表示, 指数 $e$ 限于固定范围. 零、次正规数、无穷和 NaN 需要单独编码. 用 $\operatorname{fl}(x)$ 表示把实数 $x$ 按给定舍入规则映射到浮点集合.
\end{definition}
IEEE 754 的 binary64 使用二进制, 含 53 位有效精度, 正规数的有效指数范围为 $-1022$ 到 $1023$. 这里的 53 包括正规数隐含的首位 $1$. 机器中从 $1$ 到下一个更大的浮点数的间距为 $2^{-52}$; 在舍入到最近数的模型下, 单位舍入误差为 $u=2^{-53}$. 有些软件把前一个间距称为 machine epsilon, 有些文献用同一名称指 $u$, 所以必须先核对定义.

\begin{proposition}\label{prop:rounding}
舍入到最近浮点数时, 若 $x$ 和舍入结果处于正规范围, 则存在 $|\delta|\le u=\tfrac12\beta^{1-t}$ 使 $\operatorname{fl}(x)=x(1+\delta)$. 对基本运算, 在未发生上溢、下溢且精确结果可按该模型处理时,
\[
\operatorname{fl}(a\mathbin{\circ}b)=(a\mathbin{\circ}b)(1+\delta),\qquad |\delta|\le u,
\quad \circ\in\{+,-,\times,/\}.
\]
\end{proposition}
\begin{proof}
在区间 $[\beta^e,\beta^{e+1})$ 内, 相邻浮点数间距为 $\beta^{e-t+1}$. 舍入到最近数的绝对误差不超过此间距的一半. 除以 $|x|\ge\beta^e$, 就得相对误差上界 $\tfrac12\beta^{1-t}$. 基本运算先取数学上的精确结果, 再作一次正确舍入, 便得到同一结论.
\end{proof}
次正规数的目的是让接近零的数仍可表示, 从而渐进下溢. 但它们没有统一的相对误差界; 接近零时宜用绝对误差界. 上溢、除零和 NaN 也不满足上述实数相对误差模型. 后面的推导会明确假定这些异常没有发生.

\begin{example}
解释 binary64 中为什么按通常的舍入到最近、遇中点取偶数的规则, $(1+2^{-53})-1$ 可以得到 $0$, 而不是 $2^{-53}$.
\end{example}
\begin{solution}
$1$ 附近向上的间距为 $2^{-52}$. $1+2^{-53}$ 恰好处在 $1$ 与 $1+2^{-52}$ 的中点, 中点规则选择尾数末位为偶数的 $1$. 第一步结果因此是 $1$, 第二步为零. 这里不是减法本身有很大舍入误差, 而是加法已经丢掉了所需信息.
\end{solution}
浮点加法一般不满足结合律. 编译器、并行归约和不同的矩阵乘法实现可能改变加法顺序, 从而改变最后几位. 这并不自动说明某个实现错误; 需要判断误差规模是否符合问题条件数和算法稳定性.

\originalsection{IEEE 编码与常见误解}
在 binary32 中, 一个数含 1 位符号、8 位指数和 23 位显式尾数字段; 在 binary64 中分别为 1、11 和 52 位. 正规数用隐含首位补成 24 或 53 位精度. 若指数编码为 $E$, 小数字段为 $0.M$, 正规值分别为 $(-1)^S(1.M)2^{E-127}$ 和 $(-1)^S(1.M)2^{E-1023}$. 指数全零用于零和次正规数, 全一用于无穷和 NaN. 次正规数不含隐含首位 $1$.
\begin{center}
\begin{tabular}{lcc}
\toprule
 & binary32 & binary64\\
\midrule
有效二进制位 & 24 & 53\\
最小正规正数 & $2^{-126}$ & $2^{-1022}$\\
最小次正规正数 & $2^{-149}$ & $2^{-1074}$\\
最大有限数 & $(2-2^{-23})2^{127}$ & $(2-2^{-52})2^{1023}$\\
单位舍入误差 & $2^{-24}$ & $2^{-53}$\\
\bottomrule
\end{tabular}
\end{center}
例如 binary32 的 $E=129$、$1.M=1.101_2=13/8$ 表示 $6.5$. 正规最小指数编码 $E=1$ 对应 $2^{-126}$; 同一最低尺度下允许首位为零, 最小非零尾数给出 $2^{-149}$. 渐进下溢能使相邻有限浮点数的非零差仍被表示; 若硬件选择把次正规数直接刷成零, 这种性质可能不再成立, 实验必须记录该行为.

浮点误差不是每次自动加入的独立随机数. 可精确表示的结果, 例如 $1+1=2$, 在正确舍入下仍精确; 乘零在不含 NaN 或无穷异常时仍得零. 在有效整数范围内, 浮点可以精确表示整数并精确执行结果仍可表示的整数加减. “存了多少位”与“结果准了多少位”也不同: 相消、累计误差或病态问题可以使准确位数远少于存储位数; 某些整数结果则完全准确.

binary64 中 $0.1$ 通常不可精确表示, 因其分母含因子 $5$, 二进制有限小数的分母只含 $2$. 这是一种确定性的舍入, 不是运行时随机挑选尾数. 十进制浮点可以准确表示这个十进制数, 但它同样只有有限精度, 并不对任意实数精确. 有符号零 $+0,-0$ 的区分也有用, 例如有符号无穷和某些复函数分支需要知道趋零方向. 所以“浮点数不可能恰好为零”和“零符号总无意义”都不正确.

若为平均误差建立随机模型, 必须另加假设. 假设各次舍入近似零均值、弱相关, 方差约为 $u^2$, 长为 $k$ 的误差和的均方根可为 $O(\sqrt{k}u)$ 而非最坏的 $O(ku)$. 这解释了顺序求和常比最坏界好, 成对求和可把相应树深降到 $O(\log n)$. 但输入最后几位、运算关联和舍入方向可能高度相关, 平均模型不能替代对所有输入成立的稳定性证明.

\originalsection{误差怎样沿运算链传播}
\begin{lemma}\label{lem:gamma}
若 $|\delta_i|\le u$, $\rho_i\in\{-1,1\}$, 且 $ku<1$, 则
\[
\prod_{i=1}^{k}(1+\delta_i)^{\rho_i}=1+\theta_k,
\qquad |\theta_k|\le\gamma_k:=\frac{ku}{1-ku}.
\]
\end{lemma}
\begin{proof}
先注意 $|\log(1+\delta_i)|$ 的控制也可导出相近估计, 但这里给出适用于正负指数的直接论证. 每个因子及其倒数均位于 $[1-u,(1-u)^{-1}]$, 故乘积位于 $[(1-u)^k,(1-u)^{-k}]$. Bernoulli 不等式给出 $(1-u)^k\ge1-ku$, 从而乘积不超过 $(1-ku)^{-1}=1+\gamma_k$, 也不小于 $1-ku\ge1-\gamma_k$.
\end{proof}
$\gamma_k$ 的意义是把多次小舍入合成一个可控制的扰动. 当 $ku\ll1$ 时, $\gamma_k\approx ku$; 不能在 $ku$ 已很大时仍随意用这个一阶近似.

\begin{theorem}[顺序求和]\label{thm:sum}
对输入 $x_1,\ldots,x_n$, 按 $s_1=x_1$, $s_k=\operatorname{fl}(s_{k-1}+x_k)$ 计算. 若上述浮点模型成立且 $(n-1)u<1$, 则
\[
\widehat s=\sum_{i=1}^n x_i(1+\eta_i),\quad |\eta_i|\le\gamma_{n-1},
\qquad
|\widehat s-s|\le\gamma_{n-1}\sum_{i=1}^n|x_i|,
\]
其中 $s=\sum_i x_i$.
\end{theorem}
\begin{proof}
记第 $k$ 次加法误差为 $\delta_k$. 把递推式展开,
\[
\widehat s=x_1\prod_{k=2}^n(1+\delta_k)
+\sum_{i=2}^n x_i\prod_{k=i}^n(1+\delta_k).
\]
每个输入沿运算树最多经过 $n-1$ 次舍入. 由引理 \ref{lem:gamma}, 相应乘积可写成 $1+\eta_i$, 且 $|\eta_i|\le\gamma_{n-1}$. 减去精确和并用三角不等式, 得到误差界.
\end{proof}
这一定理一方面给出了绝对误差界, 另一方面说实际结果是略微改变输入后的精确和. 但若正负数严重相消, $|s|$ 可以远小于 $\sum|x_i|$, 相对误差便很大. 这需要把算法稳定性与问题本身的敏感性分开.

\begin{definition}
当 $s\ne0$ 时, 求和问题在逐分量相对扰动下的条件数为 $\kappa_{\rm sum}=\sum_i|x_i|/|s|$. 若 $s=0$, 相对条件数可能无穷, 此时应改用绝对误差尺度.
\end{definition}
\begin{proof}[条件数公式的推导]
若 $|\Delta x_i|\le\eps|x_i|$, 则 $|\Delta s|\le\eps\sum|x_i|$. 对实输入, 令 $\Delta x_i=\eps|x_i|\operatorname{sign}(s)$ 即能取到这个绝对上界; 因而最坏相对放大倍数正是所给比值.
\end{proof}
成对求和把输入分为两半, 各自求和后再相加. 若 $n=2^m$, 每个输入仅经过 $m=\log_2 n$ 层, 同样的展开给出绝对误差不超过 $\gamma_m\sum|x_i|$. 对一般 $n$ 可用深度不超过 $\lceil\log_2n\rceil$ 的近似平衡二叉树. 顺序与成对求和都对输入是后向稳定的, 但误差常数随 $n$ 的增长不同. 递归求和的好处来自浅运算树, 不是来自递归这个程序语法本身.

\begin{proposition}[点积与矩阵乘法]\label{prop:dot}
对实数向量 $x,y$, 用通常乘法加顺序累加计算 $x^Ty$, 在 $nu<1$ 和浮点模型成立的条件下,
\[
|\operatorname{fl}(x^Ty)-x^Ty|\le\gamma_n\sum_{i=1}^n|x_i||y_i|.
\]
因此普通实矩阵乘法的结果可写成 $\widehat C=AB+E$, 满足逐分量界 $|E|\le\gamma_n|A||B|$.
\end{proposition}
\begin{proof}
先把第 $i$ 次乘法的舍入因子写入对应输入项, 再沿求和树展开. 第一项最多经历一次乘法和 $n-1$ 次加法, 其余项不会更多. 引理 \ref{lem:gamma} 控制每项总扰动为 $\gamma_n$, 再求绝对值之和. 对矩阵乘法的每个元素应用点积界即得结论. 复数实现需把实乘法、加法的额外次数吸收入常数, 不能原封不动沿用实数操作计数.
\end{proof}

\originalsection{相消与等价表达式}
相近数相减有时被笼统称为“灾难性相消”. 更精确的说法是: 减法可能把此前输入中的误差暴露并放大. 在可表示数的适当范围内, 一次减法甚至可以精确完成, 但这不恢复已经丢失的信息.

\begin{example}
在 $x$ 接近零时, 如何稳定计算 $f(x)=\sqrt{1+x}-1$?
\end{example}
\begin{solution}
直接公式先把 $1+x$ 与平方根舍入到接近 $1$ 的数, 再相减, 相对误差可能被约 $1/|x|$ 放大. 由有理化可改写为
\[
f(x)=\frac{x}{\sqrt{1+x}+1},\qquad x\ge-1.
\]
当 $x$ 很小时, 分母接近 $2$, 没有两个相近大数相减. 若输入 $x$ 已可靠且没有异常范围问题, 分母的相对小误差只引起结果的相对小误差. 在 $x=-1$ 附近应另看输入敏感性, 不能把这个改写解释为所有区间均无条件良态.
\end{solution}
线性方程组中常见的相消来自消元更新 $a_{ij}-l_{ik}u_{kj}$. 我们不能避免一切相消, 而应控制参与相减的数不被算法放大到极大, 这就是主元选择和增长因子的作用. 正交变换则保持向量的 $2$-范数, 更容易建立可靠的全局误差界.

\originalsection{范数与有限维等价性}
\begin{definition}
向量范数 $\norm{\cdot}$ 满足正定性、绝对齐次性和三角不等式. 常用范数为 $\norm{x}_1=\sum|x_i|$, $\norm{x}_2=(\sum|x_i|^2)^{1/2}$, $\norm{x}_\infty=\max|x_i|$. 由向量范数诱导的矩阵范数为
\[
\norm{A}=\sup_{x\ne0}\frac{\norm{Ax}}{\norm{x}}.
\]
矩阵的 Frobenius 范数为 $\norm{A}_F=(\sum_{ij}|a_{ij}|^2)^{1/2}$.
\end{definition}
诱导范数衡量线性映射对向量长度的最坏放大. 它直接给出 $\norm{Ax}\le\norm{A}\norm{x}$ 和 $\norm{AB}\le\norm{A}\norm{B}$. Frobenius 范数便于逐元素误差分析, 但它不是由同维向量 $2$-范数诱导的算子范数.

\begin{proposition}
对 $x\in\C^n$, 有
\[
\norm{x}_\infty\le\norm{x}_2\le\norm{x}_1,\qquad
\norm{x}_1\le\sqrt n\norm{x}_2,\qquad
\norm{x}_2\le\sqrt n\norm{x}_\infty.
\]
并且 $\norm{A}_1=\max_j\sum_i|a_{ij}|$, $\norm{A}_\infty=\max_i\sum_j|a_{ij}|$.
\end{proposition}
\begin{proof}
第一组不等式分别来自最大项不超过平方和、展开平方和后的非负交叉项, 以及 Cauchy--Schwarz 不等式; 最后一条来自各项不超过最大项. 对矩阵 $1$-范数,
$\norm{Ax}_1\le\sum_j(\sum_i|a_{ij}|)|x_j|$, 故上界是最大列和, 取对应标准基向量可达到. 对无穷范数, 先用最大行和给上界; 取各分量的相位使最大行的各项同相, 可达到该界.
\end{proof}
\begin{theorem}[有限维范数等价]
固定有限维空间上的任意两个范数 $\norm{\cdot}_a,\norm{\cdot}_b$ 都存在常数 $c,C>0$ 使 $c\norm{x}_a\le\norm{x}_b\le C\norm{x}_a$ 对所有 $x$ 成立.
\end{theorem}
\begin{proof}
先与标准 $1$-范数比较. 若 $x=\sum x_ie_i$, 三角不等式给 $\norm{x}_b\le M\norm{x}_1$, 其中 $M=\max_i\norm{e_i}_b$. 因而 $\norm{\cdot}_b$ 关于标准坐标连续: $|\norm{x}_b-\norm{y}_b|\le M\norm{x-y}_1$. 标准 $1$-范数单位球面是紧集, 连续函数 $\norm{x}_b$ 在其上取得正的最小值 $m$; 最小值不能为零, 因球面不含零向量. 缩放得 $m\norm{x}_1\le\norm{x}_b$. 对范数 $a$ 同样比较, 再合并两个界即得结论.
\end{proof}
这里的常数可以依赖维数. 因此“有限维等价”不等于大型问题中不同范数的误差界一样好; 也不能据此推断无穷维空间上所有范数都等价.

\originalsection{稳定性与一个可执行的检查原则}
\begin{definition}
设问题为 $y=f(x)$, 算法输出 $\widehat y$. 前向误差比较 $\widehat y$ 与 $f(x)$. 后向误差问是否存在小扰动 $\Delta x$ 使 $\widehat y=f(x+\Delta x)$. 若后向误差在选定范数和尺度下为 $O(u)$, 其常数可合理依赖规模而不依赖病态的具体输入, 称算法后向稳定.
\end{definition}
后向稳定不承诺结果有很多正确数字. 它承诺算法相当于把原问题轻微扰动后精确求解; 若问题本身对小扰动高度敏感, 前向误差仍会很大. 反过来, 在某个例子上前向误差很小也不证明算法对所有输入稳定.

实际检查先分三件事. 第一, 输入能否准确表达、是否存在溢出下溢. 第二, 算法是否具有合适的后向误差界. 第三, 问题条件数是否会放大这个后向误差. 后文以线性系统的残差、最小二乘的正交条件和特征向量残差具体实施这个原则.

\begin{exercise}
对 $x=(1,-1+\eps)$, 求精确和的条件数, 并解释 $\eps\to0$ 时为什么后向稳定求和也可能丢掉相对有效数字.
\end{exercise}
\begin{solution}
当 $0<\eps<1$ 时, 和为 $\eps$, 条件数为 $(2-\eps)/\eps$. 输入各自只有 $O(u)$ 相对误差时, 和的绝对扰动可以达 $O(u)$, 相对扰动可达 $O(u/\eps)$. 当 $\eps$ 与 $u$ 同阶时, 相对精度无法保证; 这源于问题相消的敏感性.
\end{solution}
\begin{exercise}
证明 $\norm{A}_F\le\sqrt{\min(m,n)}\norm{A}_2$ 对 $A\in\C^{m\times n}$ 成立. 可在读完第 \ref{ch:svd} 章后用奇异值证明, 并说明何时取等.
\end{exercise}
\begin{solution}
若非零奇异值为 $\sigma_1,\ldots,\sigma_r$, 则 $\norm{A}_F^2=\sum_{i=1}^r\sigma_i^2\le r\sigma_1^2\le\min(m,n)\norm{A}_2^2$. 非零矩阵取等要求 $r=\min(m,n)$ 且全部奇异值相同.
\end{solution}
\begin{remark}
本章对应原课 L2--4. 浮点表示、CGS/MGS 之外的误差背景来自 Per-Olof Persson 的 IEEE 手册及 Johnson 的浮点误解、顺序求和、递归求和与范数等价性手册. 误差模型的证明、例题和章内教学连接为自编补充; 不把现实硬件的所有指令和次正规数处理纳入同一无条件相对误差模型.
\end{remark}

% SOURCE: OCW L4--5, L7; no complete handout for L5. Self-authored proofs and examples.
\originalchapter{条件数与线性系统扰动}\label{ch:condition}
\originalsection{问题的敏感性}
同一套方程的精确解可能对输入极敏感. 这发生在算法开始之前. 若把两个几乎平行的直线相交, 很小的斜率变化就会让交点移动很远. 线性系统的病态性正是这一几何现象在高维中的表现. 条件数应先描述问题, 稳定性再描述解题方法.

\begin{definition}
设 $f$ 是有限维赋范空间之间的映射, $x\ne0$, $f(x)\ne0$. 相对条件数定义为
\[
\kappa_f(x)=\lim_{\eps\downarrow0}\sup_{0<\norm{\Delta x}\le\eps\norm{x}}
\frac{\norm{f(x+\Delta x)-f(x)}/\norm{f(x)}}{\norm{\Delta x}/\norm{x}}.
\]
若 $f$ 在 $x$ 处 Fr\'echet 可微, 则 $\kappa_f(x)=\norm{Df(x)}\norm{x}/\norm{f(x)}$, 其中导数范数是相应诱导范数.
\end{definition}
\begin{proof}
可微性给出 $f(x+h)-f(x)=Df(x)h+o(\norm{h})$, 余项对方向一致趋于零, 因而最坏放大上界趋于 $\norm{Df(x)}$. 有限维单位球面的紧致性保证线性映射范数可达到; 沿达到最大值的方向取趋零的 $h$, 得到相同下界.
\end{proof}
条件数依赖输入、输出尺度和范数. 若输出为零, 相对误差本身无定义, 此时使用绝对条件数或混合尺度比硬套相对条件数更恰当. 参数化方式也重要: 对矩阵条目的相对扰动与对模型参数的扰动可能给出不同问题.

\begin{example}
对 $f(x)=\sqrt{x}$ ($x>0$) 与 $g(x)=1/(1-x)$ ($x\ne1$), 分别求相对条件数.
\end{example}
\begin{solution}
由导数公式,
$\kappa_f=|x f'(x)/f(x)|=1/2$,
$\kappa_g=|x g'(x)/g(x)|=|x/(1-x)|$.
后者在 $x\to1$ 时无界, 因分母趋零; 在 $x=0$ 处相对输入尺度失效, 但绝对导数为 $1$, 问题并不因此病态.
\end{solution}

\originalsection{可逆矩阵的条件数}
\begin{definition}
对可逆方阵 $A$, 在给定诱导范数下定义 $\kappa(A)=\norm{A}\norm{A^{-1}}$. 对奇异矩阵记 $\kappa(A)=\infty$. $\kappa(A)\ge1$, 因为 $1=\norm{I}\le\norm{A}\norm{A^{-1}}$.
\end{definition}
把 $A$ 看成线性映射, $\norm{A}$ 控制最大伸长, $1/\norm{A^{-1}}$ 控制最小伸长. 二者之比衡量一个方向比另一个方向更容易被压扁多少. 第 \ref{ch:svd} 章会把它写成最大与最小奇异值之比.

\begin{proposition}\label{prop:bcondition}
固定可逆 $A$, 解 $x=A^{-1}b$ 对非零右端 $b$ 的相对条件数为
\[
\kappa_b(A,b)=\frac{\norm{A^{-1}}\norm{b}}{\norm{x}}\le\kappa(A).
\]
在 $2$-范数下对某些 $b$ 可以取到 $\kappa(A)$; 但某个特定 $b$ 的条件数可能明显更小.
\end{proposition}
\begin{proof}
映射 $b\mapsto A^{-1}b$ 是线性的, 套用定义即得精确公式. 由 $b=Ax$ 得 $\norm{b}\le\norm{A}\norm{x}$, 给出上界. 在 $2$-范数下, 取 $x$ 为最大奇异值的右奇异向量, 就有 $\norm{b}/\norm{x}=\norm{A}_2$, 故达到上界. 这里最后一步使用第 \ref{ch:svd} 章的 SVD 结论.
\end{proof}
\begin{lemma}[小扰动可逆性]\label{lem:neumann}
若 $\norm{E}<1$, 则 $I+E$ 可逆, 且 $(I+E)^{-1}=\sum_{j=0}^{\infty}(-E)^j$, $\norm{(I+E)^{-1}}\le1/(1-\norm{E})$.
\end{lemma}
\begin{proof}
矩阵级数在范数下绝对收敛, 因 $\sum\norm{E}^j<\infty$. 有限部分和乘以 $I+E$ 得 $I+(-1)^mE^{m+1}$, 其余项趋零. 极限因此是双侧逆; 三角不等式给出范数界.
\end{proof}
这个引理是扰动估计的基础. 它明确给出可逆性条件, 而不仅是“扰动很小”这种没有尺度的说法.

\begin{theorem}[线性系统扰动界]\label{thm:perturb}
设 $Ax=b$, $(A+\Delta A)(x+\Delta x)=b+\Delta b$, $A$ 可逆, $x\ne0$. 若 $\eta=\norm{A^{-1}}\norm{\Delta A}<1$, 则 $A+\Delta A$ 可逆, 且
\[
\frac{\norm{\Delta x}}{\norm{x}}
\le\frac{\kappa(A)}{1-\kappa(A)\norm{\Delta A}/\norm{A}}
\left(\frac{\norm{\Delta A}}{\norm{A}}+\frac{\norm{\Delta b}}{\norm{b}}\right),
\]
这里假定 $b\ne0$, 以使用右端相对误差.
\end{theorem}
\begin{proof}
由 $A+\Delta A=A(I+A^{-1}\Delta A)$ 和引理 \ref{lem:neumann} 可知可逆. 两个系统相减, 得
\[
\Delta x=(I+A^{-1}\Delta A)^{-1}A^{-1}(\Delta b-\Delta A\,x).
\]
因此 $\norm{\Delta x}\le\norm{A^{-1}}(\norm{\Delta b}+\norm{\Delta A}\norm{x})/(1-\eta)$. 除以 $\norm{x}$, 再用 $\norm{b}\le\norm{A}\norm{x}$, 就得到所述界.
\end{proof}
若两个输入相对误差都约为 $u$, 且 $\kappa(A)u\ll1$, 则前向误差约为 $O(\kappa(A)u)$. 但当分母接近零时, 一阶近似失效. 矩阵可能被扰动成奇异, 此时不存在有意义的统一相对解误差保证.

\begin{example}
设 $A=\diag(1,\eps)$, $b=(1,\eps)^T$, $0<\eps<1$. 比较右端扰动 $\Delta b=(0,u)^T$ 与 $\Delta b=(u,0)^T$ 的影响.
\end{example}
\begin{solution}
精确解为 $x=(1,1)^T$, $\kappa_2(A)=1/\eps$. 第一种扰动给 $\Delta x=(0,u/\eps)^T$, 第二种给 $\Delta x=(u,0)^T$. 同样大小的右端误差, 落在不同方向上会有不同影响. 条件数是最坏方向的保证, 不预测每一次误差一定被放大到上界.
\end{solution}

\originalsection{残差、后向误差与前向误差}
\begin{definition}
对近似解 $\widehat x$, 定义残差 $r=b-A\widehat x$. 范数后向误差使用允许扰动的尺度: 寻找尽量小的 $\eta$ 使
\[
(A+\Delta A)\widehat x=b+\Delta b,\quad
\norm{\Delta A}_2\le\eta\norm{A}_2,\quad
\norm{\Delta b}_2\le\eta\norm{b}_2.
\]
\end{definition}
\begin{theorem}\label{thm:backward}
若 $\widehat x\ne0$ 且 $\norm{A}_2\norm{\widehat x}_2+\norm{b}_2>0$, 上述最小后向误差为
\[
\eta_2(\widehat x)=\frac{\norm{r}_2}{\norm{A}_2\norm{\widehat x}_2+\norm{b}_2}.
\]
若仅允许扰动 $A$, 相应最小值为 $\norm{r}_2/(\norm{A}_2\norm{\widehat x}_2)$.
\end{theorem}
\begin{proof}
扰动方程等价于 $\Delta A\widehat x-\Delta b=r$, 所以三角不等式给出
$\norm{r}_2\le\eta(\norm{A}_2\norm{\widehat x}_2+\norm{b}_2)$.
若 $r=0$, 零扰动达到下界. 若 $r\ne0$, 令 $d=\norm{A}_2\norm{\widehat x}_2+\norm{b}_2$, 取
\[
\Delta A=\frac{\norm{A}_2}{d\norm{\widehat x}_2}\,r\widehat x^*,
\qquad \Delta b=-\frac{\norm{b}_2}{d}\,r.
\]
秩一矩阵的 $2$-范数为两个向量范数的乘积, 因此两项满足恰好所需的相对界, 并且 $\Delta A\widehat x-\Delta b=r$. 只扰动 $A$ 时直接取 $\Delta A=r\widehat x^*/\norm{\widehat x}_2^2$ 即可.
\end{proof}
这说明缩放残差是有效的后向准确度指标. 单看 $\norm{r}$ 不够, 因把方程同乘一个常数会改变它; 分母消除了这种整体缩放影响.

\begin{corollary}\label{cor:resforward}
若 $Ax=b\ne0$, 则
\[
\frac{\norm{x-\widehat x}}{\norm{x}}
\le\kappa(A)\frac{\norm{b-A\widehat x}}{\norm{b}}.
\]
\end{corollary}
\begin{proof}
由 $x-\widehat x=A^{-1}r$ 得分子不超过 $\norm{A^{-1}}\norm{r}$, 再用 $\norm{b}\le\norm{A}\norm{x}$ 即可.
\end{proof}
残差很小并不保证解很准, 因还要乘条件数. 一个残差相对为 $10^{-12}$ 的计算, 若条件数约为 $10^{10}$, 最坏解相对误差界仍只有 $10^{-2}$ 量级. 反过来, 大条件数的上界不说明每个实际解都这么差.

\begin{example}
令 $A=\diag(1,10^{-12})$, $b=(1,0)^T$, 近似解 $\widehat x=(1,1)^T$. 判断残差与解误差.
\end{example}
\begin{solution}
精确解 $x=(1,0)^T$. 残差为 $(0,-10^{-12})^T$, 相对残差仅 $10^{-12}$, 但相对解误差为 $1$. 它恰好位于被 $A$ 压缩的方向上. 因此停止迭代时除了残差阈值, 还应考虑所需解精度和可获得的条件数估计.
\end{solution}

\originalsection{分量尺度与缩放}
某些系统的不同方程、变量采用不同物理单位. 范数尺度可能掩盖小分量的巨大相对误差. 因而还要考虑逐分量扰动.

\begin{proposition}\label{prop:componentwise}
对实矩阵与实向量, 允许 $|\Delta A|\le\eta|A|$, $|\Delta b|\le\eta|b|$, 则近似解的最小分量后向误差为
\[
\eta_c=\max_i\frac{|r_i|}{(|A||\widehat x|+|b|)_i}.
\]
若分母、分子同为零, 该项记为零; 若分母为零而分子非零, 记为无穷.
\end{proposition}
\begin{proof}
逐行取绝对值给出必要的下界. 记第 $i$ 行分母为 $d_i>0$, 令 $t_i=r_i/d_i$, 并取
$\Delta a_{ij}=t_i|a_{ij}|\operatorname{sign}(\widehat x_j)$,
$\Delta b_i=-t_i|b_i|$.
则该行 $\sum_j\Delta a_{ij}\widehat x_j-\Delta b_i=t_i d_i=r_i$, 且所有扰动不超过 $\eta_c$ 的分量界. 分母为零的行按约定处理即可.
\end{proof}
选择对角矩阵 $D_r,D_c$ 后, 可解缩放系统 $(D_rAD_c)y=D_rb$, 最后取 $x=D_cy$. 行缩放调整方程量级, 列缩放调整变量量级. 合理缩放可改善算法行为和范数条件数, 但不能消除真正的线性相关性. 缩放也必须与原变量的误差目标一致, 不能只报告缩放空间中的好看数字.

\originalsection{误差改进与可靠停止}
若已有近似解 $\widehat x$, 可计算残差 $r=b-A\widehat x$, 求校正方程 $Ad=r$, 再更新 $\widehat x\leftarrow\widehat x+d$. 这称为迭代改进. 其数学依据是精确校正 $d=x-\widehat x$; 实际中校正也有误差, 但较高精度残差往往能补回求解阶段丢掉的数字.

\begin{proposition}
假设每次校正实际解的是 $(A+E)d=r$, 且残差与更新在本命题中视为精确. 若 $\norm{A^{-1}E}\le q<1/2$, 则改进后的误差满足
\[
\norm{e_{\rm new}}\le\frac{q}{1-q}\norm{e},\qquad e=x-\widehat x.
\]
\end{proposition}
\begin{proof}
由 $r=Ae$ 得 $d=(I+A^{-1}E)^{-1}e$. 因而
$e_{\rm new}=e-d=(I+A^{-1}E)^{-1}A^{-1}Ee$.
应用引理 \ref{lem:neumann} 即得系数 $q/(1-q)<1$.
\end{proof}
这个简单模型解释了改进的收缩机制, 但真正浮点实现还要加残差计算和更新误差的项. 当病态性使 $q$ 不小, 或残差被低精度相消污染, 改进可能停滞. 合理停止要同时检查缩放残差、相邻更新量和可用精度; 不能要求浮点算法无限降低残差.

\begin{exercise}
设 $A$ 可逆, 证明对任何非零标量 $\alpha$, 有 $\kappa(\alpha A)=\kappa(A)$. 说明为什么方程整体倍乘不改善条件数.
\end{exercise}
\begin{solution}
$\norm{\alpha A}=|\alpha|\norm{A}$, $\norm{(\alpha A)^{-1}}=|\alpha|^{-1}\norm{A^{-1}}$, 两因子相消. 它改变绝对数值尺度, 但最大伸长与最小伸长的比例不变.
\end{solution}
\begin{exercise}
若 $A$ 与 $B$ 均可逆, 证明 $\kappa(AB)\le\kappa(A)\kappa(B)$. 对 $2$-范数给出严格不等式的例子.
\end{exercise}
\begin{solution}
由次乘性和 $(AB)^{-1}=B^{-1}A^{-1}$ 即得. 取 $A=\diag(2,1)$, $B=\diag(1,2)$, 两个条件数都为 $2$, 而 $AB=2I$ 的条件数为 $1$.
\end{solution}
\begin{remark}
本章对应原课 L5 及 L7 的条件数主题, 官方资源索引在这些讲次没有独立手册. 定义、线性系统扰动、后向误差的可达构造与迭代改进分析均为自编补充, 用来补足课堂说明不能代替的推理环节.
\end{remark}

% SOURCE: OCW L7--8 HTML, no independent complete handout; self-authored mathematical development.
\originalchapter{奇异值分解与最小二乘}\label{ch:svd}
\originalsection{从伸缩方向认识矩阵}
特征值描述方阵在某些方向上的不变性, 却不能直接处理长方矩阵, 对非正规方阵也不一定准确描述伸缩. 奇异值分解把任何矩阵分为输入空间的正交旋转、坐标轴上的非负伸缩和输出空间的正交旋转. 它因此同时解释秩、条件数、最小二乘与低秩逼近.

以下对复空间使用 $x^*y$ 内积; 实矩阵只需把共轭转置改为转置. 复矩阵 $A$ 称为 Hermitian, 若 $A=A^*$; 酉矩阵满足 $U^*U=I$.

\begin{theorem}[Hermitian 谱分解]\label{thm:hermitian}
若 $H=H^*\in\C^{n\times n}$, 则存在酉矩阵 $V$ 和实对角矩阵 $\Lambda$ 使 $H=V\Lambda V^*$. 若 $x^*Hx\ge0$ 对所有 $x$ 成立, 全部特征值非负.
\end{theorem}
\begin{proof}
$ x^*Hx$ 为实数, 并在单位球面上连续, 因而取得最大值 $\lambda$ 于单位向量 $v$. 对任意 $w\perp v$, 考察归一化后的 $v+tw$ 的 Rayleigh 商; 在实参数 $t=0$ 处的一阶导数为 $2\operatorname{Re}(w^*Hv)$, 必须为零. 再以 $iw$ 替代 $w$, 得虚部也为零. 因而 $Hv$ 与所有 $w\perp v$ 正交, 所以 $Hv=\lambda v$. 对任意 $w\perp v$, 有 $v^*Hw=(Hv)^*w=0$, 说明 $v^\perp$ 在 $H$ 下不变. 在这个维数少一的空间上重复, 用归纳法得到正交归一特征基. 半正定情形对单位特征向量有 $\lambda=v^*Hv\ge0$.
\end{proof}
\begin{theorem}[奇异值分解]\label{thm:svd}
对 $A\in\C^{m\times n}$, 存在酉矩阵 $U\in\C^{m\times m}$、$V\in\C^{n\times n}$ 使
\[
A=U\Sigma V^*,\qquad
\Sigma=\diag(\sigma_1,\ldots,\sigma_p)\ \text{的长方矩阵形式},\quad
\sigma_1\ge\cdots\ge\sigma_p\ge0,\ p=\min(m,n).
\]
若 $r=\rank(A)$, 则 $A=\sum_{i=1}^r\sigma_i u_iv_i^*$, 且 $Av_i=\sigma_i u_i$, $A^*u_i=\sigma_i v_i$ 对 $i\le r$ 成立.
\end{theorem}
\begin{proof}
$A^*A$ 是 Hermitian 半正定矩阵, 由定理 \ref{thm:hermitian} 有正交特征基 $v_i$, 特征值写为 $\sigma_i^2\ge0$. 对正特征值定义 $u_i=Av_i/\sigma_i$. 于是
$u_i^*u_j=v_i^*A^*Av_j/(\sigma_i\sigma_j)=\delta_{ij}$.
这些向量可扩充为输出空间的正交归一基. 零特征值对应 $\norm{Av_i}_2^2=0$, 故 $Av_i=0$. 对全部输入基向量核对即可得 $A=U\Sigma V^*$. 最后 $A^*u_i=A^*Av_i/\sigma_i=\sigma_i v_i$, 而正奇异值的个数正是像空间维数.
\end{proof}
这个证明说明 SVD 一定存在, 但并不建议在计算中先形成 $A^*A$ 再求特征值. 因 $A^*A$ 的条件数会平方, 还会丢掉小奇异值的信息. 稳定的数值 SVD 通常先用正交变换约化为双对角矩阵, 再用专用迭代处理.

\begin{proposition}\label{prop:svdnorm}
$\norm{A}_2=\sigma_1$, $\norm{A}_F^2=\sum_i\sigma_i^2$. 对可逆方阵, $\kappa_2(A)=\sigma_1/\sigma_n$. 酉左右乘保持这两种范数和全部奇异值.
\end{proposition}
\begin{proof}
写 $x=Vy$, 则 $\norm{Ax}_2^2=\sum_i\sigma_i^2|y_i|^2\le\sigma_1^2\norm{x}_2^2$, 取 $x=v_1$ 达到. Frobenius 范数的平方等于 $\operatorname{tr}(A^*A)$, 等于特征值之和. 可逆时 $A^{-1}=V\Sigma^{-1}U^*$, 因而逆范数为 $1/\sigma_n$. 酉变换保持内积, 可从范数定义或 SVD 表示直接验证不变性.
\end{proof}
\begin{theorem}[到奇异矩阵的距离]\label{thm:singulardistance}
对可逆 $n\times n$ 矩阵,
\[
\min_{A+E\text{ 奇异}}\norm{E}_2=\sigma_n(A),\qquad
\frac{\min\norm{E}_2}{\norm{A}_2}=\frac1{\kappa_2(A)}.
\]
\end{theorem}
\begin{proof}
若 $A+E$ 奇异, 可取单位向量 $z$ 使 $(A+E)z=0$, 从而 $\norm{E}_2\ge\norm{Ez}_2=\norm{Az}_2\ge\sigma_n$. 取 $E=-\sigma_nu_nv_n^*$, 则 $(A+E)v_n=0$ 且 $\norm{E}_2=\sigma_n$, 达到下界.
\end{proof}
因此大条件数也可解释为相对靠近奇异集合. 一个只需极小相对扰动就变成奇异的矩阵, 不可能保证解对所有输入误差都稳健.

\originalsection{正交投影与伪逆}
\begin{definition}
若 $Q$ 的列正交归一, $P=QQ^*$ 称为到 $\range(Q)$ 的正交投影. 它满足 $P^*=P$, $P^2=P$. 对 $A$ 的 SVD, 定义 Moore--Penrose 伪逆
\[
A^\dagger=V\Sigma^\dagger U^*=
\sum_{i=1}^r\sigma_i^{-1}v_iu_i^*,
\]
其中正奇异值取倒数, 零奇异值保持零, 并转置长方矩阵形状.
\end{definition}
\begin{proposition}
对任意向量 $b$, $Pb$ 是子空间 $\range(Q)$ 中离 $b$ 最近的唯一向量. 最小距离为 $\norm{(I-P)b}_2$.
\end{proposition}
\begin{proof}
对 $z\in\range(Q)$, 分解 $b-z=(I-P)b+(Pb-z)$. 两项正交, 因而
$\norm{b-z}_2^2=\norm{(I-P)b}_2^2+\norm{Pb-z}_2^2$.
第二项仅在 $z=Pb$ 时为零, 得到唯一最优点.
\end{proof}
\begin{theorem}[最小二乘的全部解]\label{thm:ls}
问题 $\min_x\norm{Ax-b}_2$ 的全部解为 $x=A^\dagger b+z$, $z\in\ker A$. 其中唯一的最小范数解为 $x_\star=A^\dagger b$, 且拟合值为 $AA^\dagger b$, 残差为 $(I-AA^\dagger)b$. 最小二乘解等价于满足正规方程 $A^*(Ax-b)=0$.
\end{theorem}
\begin{proof}
令 $y=V^*x$, $c=U^*b$. 酉不变性使目标的平方变成
\[
\sum_{i=1}^r|\sigma_i y_i-c_i|^2+\sum_{i=r+1}^{m}|c_i|^2.
\]
前 $r$ 个坐标必须取 $y_i=c_i/\sigma_i$; 其余输入坐标不影响目标, 正是 $\ker A$ 的分量. 为使解范数最小, 这些自由坐标必须全为零. 这也证明投影和残差公式. 最优时残差与 $\range(A)$ 正交, 等价于 $A^*(Ax-b)=0$; 反过来, 该正交性使所有目标增量的交叉项为零, 故是全局最小点.
\end{proof}
正规方程是数学刻画, 不一定是最佳计算方法. 满列秩时, $A^*A$ 的特征值为 $\sigma_i^2$, 因而 $\kappa_2(A^*A)=\kappa_2(A)^2$. 若 $\kappa_2(A)\approx10^8$, binary64 中形成正规矩阵便可能丢掉最小谱方向. QR 直接使用 $A$ 的正交结构, 不必把条件数先平方.

\begin{example}
令 $A=\begin{pmatrix}1&0\\0&2\\0&0\end{pmatrix}$, $b=(3,4,5)^T$. 求最小二乘解与残差.
\end{example}
\begin{solution}
目标平方为 $(x_1-3)^2+(2x_2-4)^2+25$, 故 $x_\star=(3,2)^T$, 残差 $b-Ax_\star=(0,0,5)^T$. 第三个输出方向不在 $A$ 的像中, 任何输入都无法拟合这一分量. SVD 形式把这个事实写为 $AA^\dagger=\diag(1,1,0)$.
\end{solution}

\originalsection{回归、主成分与不同的观察角度}
在线性回归中, 矩阵 $A$ 的第 $i$ 行表示第 $i$ 个样本的特征, $x$ 为待拟合系数, $b_i$ 为观测输出. 线性代数角度是把 $b$ 投影到列空间; 优化角度是极小化凸二次函数; 若另假设观测噪声独立、同方差、正态, 统计角度的最大似然估计也给出同一最小二乘目标. 最后一解释依赖概率假设, 不是只由 QR 分解便能推导出数据真的正态. 若噪声有已知不同方差, 可按权重先缩放各行, 再解加权最小二乘.

主成分分析用于找数据中方差最大的正交方向. 设已中心化的数据矩阵 $D\in\R^{m\times n}$ 每行是一个样本, $m>1$. 沿单位方向 $v$ 的样本投影为 $Dv$, 方差为 $\norm{Dv}_2^2/(m-1)$. 因而最大方差方向是 $D$ 的第一右奇异向量. 在其正交补中继续最大化, 依次得到其余右奇异向量, 对应方差为 $\sigma_i(D)^2/(m-1)$.
\begin{proof}
对 $D=U\Sigma V^T$, 写单位 $v=Vy$, 则 $\norm{Dv}_2^2=\sum_i\sigma_i^2 y_i^2\le\sigma_1^2$, 第一坐标向量达到上界. 限制 $v\perp v_1$ 即使 $y_1=0$, 相同论证选出第二个坐标; 递推即可. 这也表明样本协方差 $D^TD/(m-1)$ 的特征方向与右奇异向量相同.
\end{proof}
数据先中心化是必要步骤, 否则最大方向可能主要描述均值而非波动. 主成分对变量单位和缩放敏感, 原始量纲的方差最大不等同于标准化后最大. 计算上直接对 $D$ 做 SVD, 可避免形成协方差矩阵导致的小奇异值信息损失. 第一个方向方差最大并不自动意味着它对预测任务最有用, 因这里的目标没有使用响应变量 $b$.

\originalsection{最小二乘的敏感性}
当残差不为零时, 最小二乘的输入扰动比方阵系统多一个几何因素. 设 $A$ 满列秩, $r=b-Ax$, 则 $A^*r=0$. 对微小 $\Delta A,\Delta b$, 把扰动后的正规方程保留一阶项, 得
\[
A^*A\,\Delta x
=A^*(\Delta b-\Delta A\,x)+(\Delta A)^*r+O(\norm{\Delta}^2).
\]
因此一阶解变化为
\[
\Delta x=A^\dagger(\Delta b-\Delta A\,x)
+(A^*A)^{-1}(\Delta A)^*r+O(\norm{\Delta}^2).
\]
这里的 $O(\norm{\Delta}^2)$ 是在固定满秩 $A$ 邻域中的局部记号, 常数依赖 $A,b$. 第一项类似线性方程组, 第二项却含 $\norm{(A^*A)^{-1}}=1/\sigma_n^2$ 与残差. 当拟合不好时, 改变列空间方向会使解更敏感. 所以即便使用后向稳定 QR, 原问题本身仍可能具有 $\kappa(A)^2$ 量级的矩阵扰动敏感性. 稳定算法不能消除这个问题性质.

\begin{proof}[一阶公式的推导]
扰动后的残差为 $r+\Delta b-\Delta A\,x-A\Delta x-\Delta A\Delta x$. 与 $A+\Delta A$ 的列正交, 展开 $(A+\Delta A)^*$ 与这个残差的乘积. 原项 $A^*r$ 为零, 留下一阶项 $A^*\Delta b-A^*\Delta A\,x-A^*A\Delta x+(\Delta A)^*r$. 把二阶乘积并入余项, 再左乘可逆 $A^*A$ 的逆即得结论. 满列秩保证局部解映射光滑, 因此这个线性化在小扰动极限中有效.
\end{proof}

\originalsection{最佳低秩逼近}
\begin{theorem}[截断 SVD]\label{thm:ey}
令 $A_k=\sum_{i=1}^k\sigma_i u_iv_i^*$, $0\le k<r$. 对所有秩不超过 $k$ 的矩阵 $B$, 有
\[
\norm{A-B}_2\ge\sigma_{k+1}=\norm{A-A_k}_2,
\qquad
\norm{A-B}_F^2\ge\sum_{i>k}\sigma_i^2=\norm{A-A_k}_F^2.
\]
\end{theorem}
\begin{proof}
先看算子范数. 维数为 $k+1$ 的空间 $\Span(v_1,\ldots,v_{k+1})$ 与 $\ker B$ 有非零交, 因 $\rank B\le k$. 取交中的单位向量 $z$, 有 $Bz=0$ 及 $\norm{Az}_2\ge\sigma_{k+1}$, 给出下界. 截断矩阵的差有最大奇异值 $\sigma_{k+1}$, 故达到.

再看 Frobenius 范数. 设 $P$ 为到 $\range(B)$ 的正交投影. 因 $(I-P)B=0$, 且正交投影不增 Frobenius 范数,
\[
\norm{A-B}_F^2\ge\norm{(I-P)A}_F^2
=\sum_{i=1}^r\sigma_i^2(1-\norm{Pu_i}_2^2).
\]
记 $w_i=\norm{Pu_i}_2^2$, 则 $0\le w_i\le1$, 且 $\sum_iw_i\le\rank P\le k$. 由于 $\sigma_i^2$ 降序, 把权重从较小系数处移到较大系数处只会增大和, 所以 $\sum_i\sigma_i^2w_i\le\sum_{i=1}^k\sigma_i^2$. 得所需尾和下界. $A_k$ 的误差恰好有该尾和, 故达到.
\end{proof}
若 $\sigma_k=\sigma_{k+1}$, 最优低秩子空间可能不唯一. 即使存在谱间隙, 算子范数意义下的所有最优矩阵也不能简单宣称唯一. 本定理只给出最优值与一个最优构造.

\originalsection{正则化与广义奇异值}
小奇异值方向中的数据噪声经 $1/\sigma_i$ 放大. 对近似低秩问题, 把所有非零奇异值都倒数并不一定合适. 截断伪逆取
$x_k=\sum_{i=1}^k(u_i^*b/\sigma_i)v_i$, 丢弃小奇异值方向; 它降低噪声放大, 同时引入截断偏差. 截断阈值应反映噪声和模型, 不能单由机器精度决定.

\begin{proposition}[Tikhonov 正则化]
对 $\lambda>0$, 问题 $\min_x(\norm{Ax-b}_2^2+\lambda\norm{x}_2^2)$ 有唯一解
\[
x_\lambda=(A^*A+\lambda I)^{-1}A^*b
=\sum_{i=1}^r\frac{\sigma_i}{\sigma_i^2+\lambda}(u_i^*b)v_i.
\]
\end{proposition}
\begin{proof}
目标的 Hessian 为 $2(A^*A+\lambda I)$, 对任何非零 $z$ 的二次型是 $2(\norm{Az}_2^2+\lambda\norm{z}_2^2)>0$, 因而严格凸. 令梯度为零即得第一式. 在 SVD 坐标中逐个解 $(\sigma_i^2+\lambda)y_i=\sigma_i u_i^*b$, 得第二式和唯一性.
\end{proof}
滤波因子从 $1/\sigma_i$ 改为 $\sigma_i/(\sigma_i^2+\lambda)$. 很小的 $\sigma_i$ 不再被无限放大; 但真实信号的这些方向也会被压小. 正则化是在噪声与偏差之间作数学选择, 不等同于“把不稳定算法修好”.

若惩罚为 $\norm{Lx}_2^2$, 则应同时考虑 $A$ 与 $L$. 以下给出一个条件明确、便于证明的广义 SVD 版本.
\begin{theorem}[矩阵对的广义分解]
设 $A\in\C^{m\times n}$, $L\in\C^{p\times n}$, $m,p\ge n$, 且叠放矩阵 $\begin{pmatrix}A\\L\end{pmatrix}$ 满列秩. 则存在可逆 $X$、列正交归一的 $U\in\C^{m\times n}$ 和 $V\in\C^{p\times n}$, 以及非负对角矩阵 $C,S$, 使
\[
A=UCX^{-1},\qquad L=VSX^{-1},\qquad C^2+S^2=I.
\]
\end{theorem}
\begin{proof}
先作薄 QR: $\begin{pmatrix}A\\L\end{pmatrix}=\begin{pmatrix}Q_A\\Q_L\end{pmatrix}R$, 其中 $R$ 可逆且 $Q_A^*Q_A+Q_L^*Q_L=I$. 酉对角化 $Q_A^*Q_A=W C^2W^*$, 特征值位于 $[0,1]$, 再令 $S^2=I-C^2$. $Q_AW$ 的列彼此正交, 第 $i$ 列长度为 $c_i$; 把非零列归一化, 零列补成正交归一列, 得 $Q_AW=UC$. 同理 $Q_LW=VS$. 因 $m,p\ge n$, 所需补列均可完成. 取 $X^{-1}=W^*R$ 即得分解.
\end{proof}
一般尺寸的 GSVD 需要长方对角块, 本册不把上述方块写法无条件套到所有矩阵对. 若令 $y=X^{-1}x$, 则带 $L$ 的正则化在这个坐标中按 $c_i^2+\lambda s_i^2$ 分离. 对 $\lambda>0$ 全部该分母为正, 解为
\[
x=Xy,\qquad y_i=\frac{c_i\,u_i^*b}{c_i^2+\lambda s_i^2}.
\]
$X$ 一般不是酉矩阵, 所以不能据此断言原 $x$ 的欧氏范数得到保持. 数值求解也可直接对叠放矩阵 $\begin{pmatrix}A\\\sqrt\lambda L\end{pmatrix}$ 作 QR, 避免显式形成正规矩阵.

\begin{exercise}
证明 $AA^\dagger$ 和 $A^\dagger A$ 分别是到 $\range(A)$ 与 $\range(A^*)$ 的正交投影, 并推导 $AA^\dagger A=A$.
\end{exercise}
\begin{solution}
在 SVD 中, $AA^\dagger=\sum_{i=1}^ru_iu_i^*$, $A^\dagger A=\sum_{i=1}^rv_iv_i^*$. 它们正是相应正交基组成的投影. $A$ 的每个列向量都属于 $\range(A)$, 对其投影不变, 所以 $AA^\dagger A=A$.
\end{solution}
\begin{exercise}
对 $A=\diag(1,10^{-6})$, $b=(1,10^{-6}+10^{-4})^T$, 比较精确逆解与取 $\lambda=10^{-8}$ 的 Tikhonov 解. 说明为什么“更接近数据”不一定等于“更接近真实信号”.
\end{exercise}
\begin{solution}
若无噪声真实数据为 $(1,10^{-6})^T$, 真实解是 $(1,1)^T$. 精确逆解是 $(1,101)^T$. 正则化解第一分量为 $1/(1+10^{-8})$, 第二分量为 $10^{-6}(10^{-6}+10^{-4})/(10^{-12}+10^{-8})\approx0.0101$. 它仍对第二个真实分量有偏差, 但避免了 $101$ 的噪声放大. 正则化强度应由噪声和先验选择, 这个数值比较不证明给定 $\lambda$ 最优.
\end{solution}
\begin{remark}
本章对应原课 L7--8 的 SVD、回归、广义 SVD 主题, 两讲由 Alan Edelman 客座讲授. 这些讲次主要给课堂说明与外部示例, 没有一份覆盖本章证明的官方手册. 本章数学证明、GSVD 的限定版本、正则化例题均为自编补充.
\end{remark}

% SOURCE: OCW L9--10; Persson Gram-Schmidt and Householder/Givens handouts.
\originalchapter{正交化与 QR 分解}\label{ch:qr}
\originalsection{把列空间换成正交坐标}
给定满列秩矩阵 $A=(a_1,\ldots,a_n)\in\C^{m\times n}$, $m\ge n$, 希望寻找同一列空间的正交归一基 $q_1,\ldots,q_n$. 每个 $a_j$ 仅用前 $j$ 个基向量表示时, 系数矩阵为上三角, 即薄 QR 分解 $A=QR$, $Q^*Q=I_n$, $R$ 上三角. 这个分解把几何部分放在 $Q$, 把尺度和线性相关程度放在 $R$.

\begin{theorem}
满列秩矩阵存在薄 QR 分解. 若要求 $R$ 的对角元为正实数, 该分解唯一.
\end{theorem}
\begin{proof}
对各列依次去掉已有基方向的投影, 得
\[
v_j=a_j-\sum_{i<j}(q_i^*a_j)q_i,\quad
r_{jj}=\norm{v_j}_2>0,\quad q_j=v_j/r_{jj}.
\]
满列秩保证 $v_j\ne0$. 设 $r_{ij}=q_i^*a_j$ ($i<j$), 就得存在性. 若 $A=QR=\widetilde Q\widetilde R$, 两组 $Q$ 的列空间相同, $S=Q^*\widetilde Q=R\widetilde R^{-1}$ 既酉又上三角. 上三角酉矩阵必为对角酉矩阵: 从第一列长度与正交性开始逐列归纳即可. 对角正实的条件使 $S$ 全部对角元为 $1$, 因而两分解相同.
\end{proof}
秩亏时可以做带列主元 QR 或允许零对角的分解, 但上面“归一化非零余量”和唯一性的条件不能直接保留.

\originalsection{经典与改进 Gram--Schmidt}
经典 Gram--Schmidt (CGS) 先用原列 $a_j$ 计算全部系数 $r_{ij}=q_i^*a_j$, 再作一次整体减法. 改进 Gram--Schmidt (MGS) 则让余量不断更新:
\[
v\leftarrow a_j;\qquad
r_{ij}\leftarrow q_i^*v,\quad v\leftarrow v-q_i r_{ij}
\quad (i=1,\ldots,j-1).
\]
最后 $r_{jj}=\norm{v}_2$, $q_j=v/r_{jj}$. 在精确算术中, 因已有 $q_i$ 彼此正交, 两者完全等价; 在浮点算术中, CGS 与 MGS 使用的点积输入不同, 误差因此不同.

\begin{proposition}
设 $q_i$ 正交归一, 则
\[
I-QQ^*=(I-q_kq_k^*)\cdots(I-q_1q_1^*),\qquad
Q=(q_1,\ldots,q_k).
\]
\end{proposition}
\begin{proof}
展开乘积时, 因 $q_i^*q_j=0$ ($i\ne j$), 含两个不同秩一投影的交叉项全为零. 留下 $I-\sum q_iq_i^*=I-QQ^*$.
\end{proof}
从矩阵角度看, Gram--Schmidt 是在 $A$ 右边依次乘上三角矩阵, 用线性组合把列正交化, 最后有 $Q=AR^{-1}$. 若这些三角变换病态, 正交方向容易受到误差污染. Householder 则在 $A$ 左边乘酉矩阵把它三角化, 把病态三角因子留在结果 $R$ 而不用于生成 $Q$.

CGS 实现左边的整体投影, MGS 实现右边的顺序投影. 当已有正交性被舍入破坏时, 这个等式不再精确, MGS 能较及时地删除仍残留的旧方向. 但单次 MGS 也不是在任意病态输入上都保持机器精度正交.

\begin{example}
考虑三列 $a_1=(1,\eps,0,0)^T$, $a_2=(1,0,\eps,0)^T$, $a_3=(1,0,0,\eps)^T$, 其中 $0<\eps\ll1$, 并假设在演示舍入模型下 $1+\eps^2$ 被舍入为 $1$. 比较 CGS 与 MGS 的第三列.
\end{example}
\begin{solution}
先得 $q_1\approx(1,\eps,0,0)^T$. 第二列减去 $q_1$ 后约为 $(0,-\eps,\eps,0)^T$, 故 $q_2\approx(0,-1,1,0)^T/\sqrt2$.

CGS 用原 $a_3$ 计算 $q_2^Ta_3\approx0$, 所以第三列余量约为 $(0,-\eps,0,\eps)^T$. 归一化后 $q_3\approx(0,-1,0,1)^T/\sqrt2$, 与 $q_2$ 的内积约为 $1/2$, 正交性严重破坏.

MGS 先从 $a_3$ 去掉 $q_1$ 得 $v=(0,-\eps,0,\eps)^T$, 再算 $q_2^Tv=\eps/\sqrt2$. 去掉这个投影后得 $(0,-\eps/2,-\eps/2,\eps)^T$, 归一化为 $(0,-1,-1,2)^T/\sqrt6$, 此时与 $q_2$ 正交到演示模型的精度. 这个例子说明运算次序的差异, 不是对所有误差情形的严格 MGS 稳定性证明.
\end{solution}
再正交化是在第一遍得到 $v$ 后再对已有 $Q$ 作一遍投影删除, 并累加到 $R$ 的系数. 它常能显著恢复正交性, 但若输入实际已到数值秩亏、余量被噪声淹没, 不能凭第二遍凭空恢复有效方向. 此时应报告秩判定或使用带主元 QR/SVD.

\originalsection{Householder 反射}
\begin{definition}
对非零向量 $v$, Householder 反射为 $H=I-2vv^*/(v^*v)$. 它沿 $v$ 方向变号, 在 $v^\perp$ 上不变.
\end{definition}
\begin{proposition}
$H^*=H$, $H^2=I$, 所以 $H$ 酉. 对非零 $x$, 令 $\alpha=-\operatorname{phase}(x_1)\norm{x}_2$, $v=x-\alpha e_1$, 其中 $x_1=0$ 时相位取 $1$. 则 $Hx=\alpha e_1$.
\end{proposition}
\begin{proof}
$P=vv^*/(v^*v)$ 满足 $P^*=P$, $P^2=P$, 因而 $(I-2P)^2=I$. 又 $\overline\alpha x_1=-|x_1|\norm{x}_2$, 所以
\[
v^*x=\norm{x}_2^2+|x_1|\norm{x}_2,
\quad v^*v=2(\norm{x}_2^2+|x_1|\norm{x}_2).
\]
故 $2v(v^*x)/(v^*v)=v$, $Hx=x-v=\alpha e_1$.
\end{proof}
选择负相位使 $v_1=x_1-\alpha$ 的两项同方向相加, 避免相近数相减. 若反过来取同相位的 $\alpha$, 当 $x$ 已接近坐标轴时 $v$ 可能几乎为零, 其方向受到舍入噪声支配. 这是从数学构造走向稳定实现的关键细节.

QR 算法依次对剩余子矩阵的第一列作反射. 第 $j$ 个反射只作用于行 $j,\ldots,m$, 不破坏此前列的零元. 得到
\[
H_n\cdots H_1 A=\begin{pmatrix}R\\0\end{pmatrix},\qquad
Q_{\rm full}=H_1\cdots H_n.
\]
只取 $Q_{\rm full}$ 的前 $n$ 列便是薄 $Q$. 实际通常保存每个反射向量与标量, 不显式构造所有 $m\times m$ 反射矩阵. 要算 $Q^*b$, 按消元顺序把反射作用于 $b$ 即可.

\begin{example}
求把 $x=(3,4)^T$ 变到第一坐标轴的 Householder 反射.
\end{example}
\begin{solution}
$\norm{x}_2=5$, 取 $\alpha=-5$, $v=(8,4)^T$, $v^Tv=80$. 因此
\[
H=I-\frac1{40}\begin{pmatrix}64&32\\32&16\end{pmatrix}
=\begin{pmatrix}-3/5&-4/5\\-4/5&3/5\end{pmatrix},\quad
Hx=(-5,0)^T.
\]
$H$ 两行的长度均为 $1$、内积为零, 可直接核查正交性.
\end{solution}

\originalsection{Givens 旋转与局部消元}
对实数 $a,b$ 不同时为零, 令 $r=\sqrt{a^2+b^2}$, $c=a/r$, $s=b/r$. 平面旋转
\[
G=\begin{pmatrix}c&s\\-s&c\end{pmatrix}
\]
满足 $G^TG=I$, $G(a,b)^T=(r,0)^T$. 把这个 $2\times2$ 块嵌入单位矩阵, 就能只混合两行, 消去某个指定元素. 当 $a=b=0$ 时取恒等变换.

直接计算 $a^2+b^2$ 可能上溢, 即使最终 $r$ 可表示. 可先取 $t=\max(|a|,|b|)$, 再算 $r=t\sqrt{(a/t)^2+(b/t)^2}$; 实现中的 hypot 函数采用类似尺度保护. 复数 Givens 需要共轭相位的专门公式, 不应把实数公式原样套入.

Householder 一次消去一列的一整段, 适合稠密 QR; Givens 仅动两行, 适合已有结构的矩阵、稀疏局部更新和迭代算法中的小型 Hessenberg 最小二乘. 对一般稠密 $m\times n$ 矩阵, 逐项 Givens QR 的主导操作数约为 $3mn^2-n^3$, 约比 Householder 高 $50\%$. 两种变换都靠正交性控制范数, 选择主要取决于待保留的结构和数据访问方式.

\originalsection{误差与计算成本}
正交变换本身条件数为 $1$, 不会把误差的 $2$-范数放大. 为看到这点怎样进入算法分析, 给出一个局部扰动的累积结论.
\begin{proposition}\label{prop:orthaccum}
设第 $j$ 步实际结果满足 $A_j=H_jA_{j-1}+E_j$, 其中 $H_j$ 是精确酉矩阵. 则
\[
A_s=(H_s\cdots H_1)(A_0+\Delta A),\qquad
\norm{\Delta A}_{2,F}\le\sum_{j=1}^s\norm{E_j}_{2,F}.
\]
\end{proposition}
\begin{proof}
逐步展开得
$A_s=H_s\cdots H_1A_0+\sum_jH_s\cdots H_{j+1}E_j$.
左乘 $(H_s\cdots H_1)^*$, 把误差项定义为 $\Delta A$. 每项仅被酉矩阵左右作用, 范数保持不变, 再用三角不等式即得.
\end{proof}
标准尺度保护的 Householder QR 在正常浮点模型下有 $A+\Delta A=\widetilde Q\widehat R$, $\widetilde Q$ 精确列正交, 且 $\norm{\Delta A}_F\le C(m,n)u\norm{A}_F+O(u^2)$. 存储的 $\widehat Q$ 通常只有近似正交, 它与分析中的精确 $\widetilde Q$ 不应混为一谈. 上述累积命题解释了为什么误差不经过病态三角因子的放大; 完整逐指令常数还需计入反射向量生成、点积与舍入, 本册不把这个实现层面的标准结果伪装成只由命题便已无条件证明.

对实 $m\times n$ ($m\ge n$) 稠密矩阵, Householder QR 的主导操作数为 $2mn^2-\tfrac23n^3$. 其来源是第 $j$ 步对约 $(m-j+1)\times(n-j)$ 尾块作一次秩一更新, 每元素约四次实数加乘, 累加三次多项式的主导项即可. CGS/MGS 的主导操作数约 $2mn^2$, 但正交性保证不同. 最快的实现还会把多个反射合并成块以调用高效矩阵乘法, 第 \ref{ch:performance} 章讨论数据移动.

\originalsection{用 QR 解最小二乘}
对满列秩 $A$, 若 $A=Q_{\rm full}\begin{pmatrix}R\\0\end{pmatrix}$, 把 $Q_{\rm full}^*b=(c_1,c_2)^T$ 按前 $n$ 个坐标分块, 则
\[
\norm{Ax-b}_2^2=\norm{Rx-c_1}_2^2+\norm{c_2}_2^2.
\]
所以先用反射求 $c_1$, 再回代解 $Rx=c_1$. 残差范数就是 $\norm{c_2}_2$, 不必先计算一个易相消的 $b-Ax$ 才知道理论拟合误差. 实际仍应回到原输入检查残差和最优性.

\begin{example}
用薄 QR 求 $A=\begin{pmatrix}1&1\\1&0\\0&1\end{pmatrix}$, $b=(1,2,0)^T$ 的最小二乘解.
\end{example}
\begin{solution}
由第一列得到 $q_1=(1,1,0)^T/\sqrt2$, $r_{11}=\sqrt2$. 第二列投影系数 $r_{12}=1/\sqrt2$, 余量为 $(1/2,-1/2,1)^T$, 故 $q_2=(1,-1,2)^T/\sqrt6$, $r_{22}=\sqrt{3/2}$. 于是
\[
R=\begin{pmatrix}\sqrt2&1/\sqrt2\\0&\sqrt{3/2}\end{pmatrix},
\qquad Q^Tb=(3/\sqrt2,-1/\sqrt6)^T.
\]
第二行给 $x_2=-1/3$, 第一行给 $x_1=5/3$. 残差 $b-Ax=(-1/3,1/3,1/3)^T$, 与两列均正交, 范数为 $1/\sqrt3$.
\end{solution}
若矩阵近秩亏, 回代会放大小对角造成的误差. 带列主元 QR 每步在剩余列中选择较大的余量范数, 作 $AP=QR$. 它有助于揭示数值秩和选择较独立列, 但不能无条件代替 SVD 的最佳秩判定. 数值秩阈值必须考虑数据噪声、矩阵尺度和目标误差.

\begin{exercise}
设满列秩 $A=QR$, 证明 $\kappa_2(R)=\kappa_2(A)$, 并据此解释 QR 为什么不改变问题本身的条件数.
\end{exercise}
\begin{solution}
$A^*A=R^*Q^*QR=R^*R$, 所以二者全部奇异值相同. QR 避免额外引入条件数平方的计算路线, 但原 $A$ 的小奇异值仍由 $R$ 保留.
\end{solution}
\begin{exercise}
对非零 $v$, 证明 Householder 反射的行列式为 $-1$, 并确定其全部特征值.
\end{exercise}
\begin{solution}
$v$ 是特征值 $-1$ 的特征向量, $v^\perp$ 上特征值都是 $1$. 在这组基中反射为 $\diag(-1,1,\ldots,1)$, 因而行列式为 $-1$.
\end{solution}
\begin{remark}
本章整理 Persson 的 Gram--Schmidt 与 Householder/Givens 手册, 对应 L9--10. 投影乘积、CGS/MGS 的演示例、反射构造与计算结构来自这些手册; 唯一性证明、误差累积命题、完整最小二乘计算和秩判定说明为自编补充.
\end{remark}

% SOURCE: OCW L13--14 HTML, no independent full handout. Self-authored derivations and examples.
\originalchapter{LU 与 Cholesky 直接法}\label{ch:direct}
\originalsection{消元就是分解}
解 $Ax=b$ 时, 高斯消元逐列消去对角线以下的元素. 如果把每一步减去的倍数保存下来, 原矩阵就被分解为下三角与上三角的乘积. 因此一次分解可以服务多个右端, 不必每次重新消元.

考虑第一步分块
\[
A=\begin{pmatrix}a&v^*\\w&B\end{pmatrix},\qquad a\ne0.
\]
减去第一行的倍数后剩余块为 $S=B-wv^*/a$, 称为 Schur 补. 恒等式
\[
A=\begin{pmatrix}1&0\\w/a&I\end{pmatrix}
\begin{pmatrix}a&v^*\\0&S\end{pmatrix}
\]
正是消元的矩阵表达. 接着对 $S$ 作同一操作, 直到得到 $A=LU$, $L$ 为单位下三角, $U$ 为上三角.

\begin{theorem}\label{thm:lu}
若 $A\in\C^{n\times n}$ 的全部顺序主子式 $\det A_{1:k,1:k}$ ($k=1,\ldots,n$) 非零, 则存在唯一的 $A=LU$, $L$ 单位下三角, $U$ 上三角且对角非零. 反之, 若存在这样的分解, 全部顺序主子式非零.
\end{theorem}
\begin{proof}
$k=1$ 保证首主元 $a\ne0$. 分块消元恒等式说明 $S$ 的前 $k-1$ 阶顺序主子式为 $\det A_{1:k,1:k}/a$, 因而全非零. 对维数归纳得到存在性. 若 $LU=\widetilde L\widetilde U$, 则 $\widetilde L^{-1}L=\widetilde U U^{-1}$, 左边单位下三角、右边上三角, 故二者都是 $I$, 得唯一性. 反向对分解取前 $k$ 行列, 该块行列式为 $\prod_{i=1}^k u_{ii}\ne0$.
\end{proof}
无主元 LU 不是对任意可逆矩阵都可执行. 例如 $\begin{pmatrix}0&1\\1&0\end{pmatrix}$ 可逆, 但首主元为零. 即使主元只是很小而不为零, 乘子也会很大, 造成增长与舍入损失.

\originalsection{部分主元选择}
在第 $k$ 列, 部分主元策略从当前行 $k,\ldots,n$ 中选择绝对值最大的条目, 把该行与第 $k$ 行交换. 精确算术中可逆矩阵的这个剩余列不会全为零. 最终得到 $PA=LU$, $P$ 为置换矩阵, 且消元乘子满足 $|l_{ik}|\le1$.

一次分解后的求解流程是 $Ly=Pb$, 再解 $Ux=y$. 下三角系统用前代,
$y_i=(c_i-\sum_{j<i}l_{ij}y_j)/l_{ii}$;
上三角系统用回代,
$x_i=(y_i-\sum_{j>i}u_{ij}x_j)/u_{ii}$.
单位下三角的 $l_{ii}=1$, 无需除法. 每个三角求解约需 $n^2$ 次实数加乘.

\begin{remark}
实现行交换时, 不仅交换尚未消元的矩阵行, 也要交换已经保存的前 $k-1$ 列乘子. 否则程序中的 $L$ 与累计置换 $P$ 不匹配, 即便最终 $U$ 看起来上三角, 也未必满足 $PA=LU$.
\end{remark}
\begin{example}
对 $A=\begin{pmatrix}\delta&1\\1&1\end{pmatrix}$, $0<\delta\ll1$, 比较无主元与部分主元 LU.
\end{example}
\begin{solution}
无主元分解为
\[
L=\begin{pmatrix}1&0\\1/\delta&1\end{pmatrix},\qquad
U=\begin{pmatrix}\delta&1\\0&1-1/\delta\end{pmatrix}.
\]
原矩阵元素不超过 $1$, 中间却出现 $1/\delta$. 若舍入误差落在这些大数上, 后续相减可能留下远大于输入精度的误差.

部分主元先交换两行, 得
\[
PA=\begin{pmatrix}1&1\\\delta&1\end{pmatrix},\qquad
L=\begin{pmatrix}1&0\\\delta&1\end{pmatrix},\quad
U=\begin{pmatrix}1&1\\0&1-\delta\end{pmatrix}.
\]
乘子和中间元素保持温和. 此处矩阵本身在 $\delta\to0$ 时并不趋奇异; 无主元算法的失败风险不能归咎于问题必然病态.
\end{solution}
部分主元只交换行, 开销约为 $O(n^2)$ 的比较, 不改变分解的 $\tfrac23n^3$ 主导实数操作数. 完全主元在剩余子矩阵中选最大条目并交换行列, 得 $PAQ=LU$, 更强地控制增长但增加搜索和数据移动. 实际选法需兼顾可靠性、稀疏结构及计算成本.

\originalsection{增长因子与后向误差}
\begin{definition}
设消元过程中当前矩阵依次为 $A^{(k)}$, 定义元素增长因子
\[
\rho=\frac{\max_{i,j,k}|a_{ij}^{(k)}|}{\max_{i,j}|a_{ij}|}.
\]
采用只计 $U$ 元素的版本时, 需明确与这里的定义有细微区别.
\end{definition}
乘子小并不自动保证所有中间元素小. 每次 Schur 补更新仍可能把原来同号或异号的数相减为更大的数. 部分主元的可靠性要通过 $\rho$ 进入误差分析.

\begin{proposition}
部分主元高斯消元在精确算术下满足 $\rho\le2^{n-1}$. 对适当矩阵, 若相等主元优先保留当前行, 这个指数上界可以达到.
\end{proposition}
\begin{proof}
因 $|l_{ik}|\le1$, 一次更新的每个元素满足
$|a_{ij}^{(k+1)}|\le|a_{ij}^{(k)}|+|l_{ik}a_{kj}^{(k)}|\le2M_k$,
其中 $M_k$ 是当前最大绝对值. 最多 $n-1$ 次更新, 得上界. 达到的构造是主对角为 $1$、主对角以下为 $-1$、最后一列全为 $1$、其余上三角元素为零的矩阵. 消元时乘子均为 $-1$, 最后一列沿主元行依次成为 $1,2,4,\ldots,2^{n-1}$, 且相等主元不触发行交换.
\end{proof}
例如 $n=4$ 时的构造为
\[
\begin{pmatrix}1&0&0&1\\-1&1&0&1\\-1&-1&1&1\\-1&-1&-1&1\end{pmatrix},
\]
最后一个主元为 $8$. 因此部分主元 LU 不能无条件称为对所有输入具有温和常数的后向稳定算法. 它通常有效, 但要认识增长例外, 必要时改用更强主元或正交方法.

\begin{proposition}[LU 因子残差的标准尺度]
在常规实数浮点模型下, 已完成的高斯消元因子可解释为
\[
PA+\Delta A=\widehat L\widehat U,\qquad
|\Delta A|\le\gamma_{cn}|\widehat L||\widehat U|,
\]
这里 $c$ 是由固定实现的每项操作链长度确定的常数, $cnu<1$, 并假定没有上溢、下溢及异常主元. 本命题给出误差量级, 不声称所有实现都具有同一个精确 $\gamma_n$ 常数.
\end{proposition}
\begin{proof}[误差尺度的推导]
每个因子关系对应一次已计算乘子和消元点积. 将计算的关系反向写成原条目加误差,
\[
(PA)_{ij}+\Delta a_{ij}=\sum_{k\le\min(i,j)}\widehat l_{ik}\widehat u_{kj}.
\]
若最后一个未知量通过除以主元生成, 把除法的因子移回等式左边; 引理 \ref{lem:gamma} 的正负指数版本同样控制它. 点积每项累积的舍入因子数为至多一个与 $n$ 同阶的常数. 因此逐条误差不超过对应绝对值乘积之和乘 $\gamma_{cn}$ (常数 $c$ 由相应操作组织确定). 把逐条关系组成矩阵, 即得到所述逐分量尺度. 这一步没有说 $|L||U|$ 与 $|A|$ 一样大; 两者的差距正是增长风险.
\end{proof}
再加前代、回代的误差, 可得到最终计算解是一个邻近系统的精确解, 扰动仍受 $|L||U|$ 类尺度控制. 对实际浮点因子, 应用实际计算增长 $\widehat\rho=\max|\widehat u_{ij}|/\max|a_{ij}|$, 而不是直接代入精确消元的 $\rho$. 若计算乘子仍满足 $|\widehat l_{ij}|\le1+O(u)$, 则 $\norm{\widehat L}_\infty\le n(1+O(u))$, $\norm{\widehat U}_\infty\le n\widehat\rho\max|a_{ij}|$. 一个粗略界因此为
$\norm{\Delta A}_\infty/\norm{A}_\infty\le\gamma_{cn} n^2\widehat\rho(1+O(u))$.
真实界往往更好, 但这个公式清楚说明: 只有增长温和, 后向误差才与机器精度同阶. 前向误差还要按第 \ref{ch:condition} 章乘问题条件数.

\originalsection{正定结构与 Cholesky 分解}
\begin{definition}
Hermitian 矩阵 $A$ 若对所有非零 $x$ 满足 $x^*Ax>0$, 称为 Hermitian 正定. 实情形称对称正定 (SPD). 它的全部特征值为正.
\end{definition}
\begin{theorem}[Cholesky 分解]\label{thm:chol}
Hermitian 正定矩阵存在唯一分解 $A=LL^*$, 其中 $L$ 为下三角且对角元为正实数.
\end{theorem}
\begin{proof}
按维数归纳. 写 $A=\begin{pmatrix}a&w^*\\w&B\end{pmatrix}$, 正定性给 $a>0$. 令 $\ell_{11}=\sqrt a$, $\ell_{21}=w/\sqrt a$. 对非零 $z$, 取向量 $(-w^*z/a,z)^T$, 可算得
\[
\begin{pmatrix}-w^*z/a\\z\end{pmatrix}^{\!*}
A\begin{pmatrix}-w^*z/a\\z\end{pmatrix}
=z^*(B-ww^*/a)z>0.
\]
所以 Schur 补也是正定, 可继续归纳. 拼接下三角因子得存在性. 首对角必须为正平方根, 第一列随之唯一, 剩余块再由归纳唯一, 得全部唯一性.
\end{proof}
Cholesky 不需要主元, 因每一步正定 Schur 补的对角主元均为正. 实数递推公式为
\[
\ell_{kk}=\sqrt{a_{kk}-\sum_{j<k}\ell_{kj}^2},\qquad
\ell_{ik}=\frac{a_{ik}-\sum_{j<k}\ell_{ij}\ell_{kj}}{\ell_{kk}},\quad i>k.
\]
复数版本的平方与乘积应改为绝对值平方及带共轭的内积. 只处理一个三角半部, 主导操作数为 $\tfrac13n^3$, 约为 LU 的一半. 解系统时作 $Ly=b$, $L^*x=y$.

\begin{example}
对 $A=\begin{pmatrix}4&2&0\\2&5&1\\0&1&3\end{pmatrix}$ 求 Cholesky 分解.
\end{example}
\begin{solution}
第一列为 $\ell_{11}=2$, $\ell_{21}=1$, $\ell_{31}=0$. 第二列给 $\ell_{22}=\sqrt{5-1}=2$, $\ell_{32}=1/2$. 最后 $\ell_{33}=\sqrt{3-1/4}=\sqrt{11}/2$. 因而
\[
L=\begin{pmatrix}2&0&0\\1&2&0\\0&1/2&\sqrt{11}/2\end{pmatrix}.
\]
直接相乘恢复 $A$, 同时正对角保证其正定, 因 $x^TAx=\norm{L^Tx}_2^2>0$ 对非零 $x$ 成立.
\end{solution}
对已完成的浮点 Cholesky, 标准分量误差界为 $|\Delta A|\le\gamma_{n+1}|\widehat L||\widehat L|^*$ 类形式. 与 LU 的区别是 $A$ 正定时因子尺度受到二次型结构限制, 不会发生部分主元 LU 那种任意指数增长. 对角主元计算若变成非正, 可能表示输入并非正定, 也可能表示矩阵离半正定边界太近, 舍入已破坏正定性. 此时不能简单把负数改成小正数继续并声称解了原问题; 应检查对称性、数据与尺度, 或明确采用正则化的新问题.

\originalsection{带状、三对角与不定问题}
带宽为 $w$ 的矩阵若消元保持带状, 可以把成本从 $O(n^3)$ 降到约 $O(nw^2)$, 存储降到 $O(nw)$. 对 SPD 带状矩阵, Cholesky 保持同一下带宽: 更新只涉及此前有公共非零列的位置, 不会越出带宽. 三对角 SPD 情形因 $w=1$ 可用 $O(n)$ 成本分解和求解.

对称但不正定的矩阵不能直接用实 Cholesky. 常用 $PAP^T=LDL^T$, $D$ 为 $1\times1$ 或 $2\times2$ 块对角, 配合对称主元策略. 为什么需要 $2\times2$ 块? 矩阵 $\begin{pmatrix}0&1\\1&0\end{pmatrix}$ 可逆且对称, 两个对角元都为零, 单个非零对角主元无从选择; 取整块即可继续. 这里的分解结构与正定性保证不同, 不应混称为 Cholesky.

多右端问题应复用分解. 通常也不应为解 $Ax=b$ 先显式计算 $A^{-1}$: 这增加成本和舍入步骤, 并丢掉结构化三角求解的直接后向误差解释. 若任务确实需要逆矩阵条目或多次线性映射, 仍应根据用途决定, 而不是把“逆”作为默认求解步骤.

\begin{exercise}
证明 SPD 矩阵的每个顺序主子式为正, 并解释为什么这保证无主元 LU 存在.
\end{exercise}
\begin{solution}
对前 $k$ 阶主子矩阵 $A_k$, 任意非零 $z$ 扩展为 $(z,0)^T$, 有 $z^*A_kz>0$. 因而 $A_k$ 正定, 所有特征值为正, 其行列式也为正. 定理 \ref{thm:lu} 的条件随即满足. Cholesky 进一步利用了对称性, 不只是无主元 LU 的另一个名字.
\end{solution}
\begin{exercise}
设 $A=\begin{pmatrix}2&-1&0\\-1&2&-1\\0&-1&2\end{pmatrix}$. 写出其 Cholesky 对角和次对角, 并说明所有更低位置为何保持零.
\end{exercise}
\begin{solution}
$\ell_{11}=\sqrt2$, $\ell_{21}=-1/\sqrt2$, $\ell_{22}=\sqrt{3/2}$, $\ell_{32}=-\sqrt{2/3}$, $\ell_{33}=\sqrt{4/3}$, 而 $\ell_{31}=0$. 三对角结构在第一步余块中仍保持三对角, 继续递推也不会产生跨越一个下对角的非零.
\end{solution}
\begin{remark}
本章对应原课 L13--14 的 LU、部分主元、Cholesky 与结构化求解器. 这两讲的官方索引明确没有手册. 本章分解存在性、增长构造、Cholesky 证明和例题为自编补充; 浮点误差常数依具体实现的操作组织而变, 不据一个抽象公式保证任意程序实现均稳定.
\end{remark}

% SOURCE: OCW L11--12; Johnson 2008 matrix multiplication performance and cache-oblivious experiments, ideal cache terminology.
\originalchapter{计算成本与数据移动}\label{ch:performance}
\originalsection{操作数为什么不能预测运行时间}
两个矩阵乘法程序都按 $c_{ij}=\sum_k a_{ik}b_{kj}$ 计算, 都有约 $2n^3$ 次实数加乘, 实际速度却可能相差很大. 原课程的实验是在 2.66GHz Intel Core 2 Duo、gcc 4.1.2 与 ATLAS 3.6.0 的历史平台上进行, 展示缓存容量、访问顺序和寄存器复用的影响. 这些曲线不能被当作今日硬件的测试结果, 但它们提出的问题仍适用于现代算法设计: 数据从哪里来, 会被使用几次, 下一次使用前是否已经被替换?

算 $q(q^*a)$ 时, 先算长度 $m$ 的点积再作向量缩放, 成本为 $O(m)$. 若先形成 $qq^*$ 再乘 $a$, 要存储 $m^2$ 个元素并作 $O(m^2)$ 运算. 矩阵乘法在精确算术中结合, 但括号的位置既改变成本, 也可能改变舍入路径. 数学表达式必须被转化为恰当的计算图.

\begin{definition}
在理想二级存储模型中, 快速存储可容纳 $Z$ 个数, 主存容量充足. 需要的数据已在快速存储中称命中, 否则称缺失, 需要把数据搬入并可能替换旧数据. 完全相联表示任意数据都可放到任意位置; 最优替换使用未来访问信息, 淘汰最久不会再用的数据. 缓存复杂度 $Q(n;Z)$ 计数据移动量或缺失次数, 必须先说明一个移动单位是一个数还是一个缓存行.
\end{definition}
现实缓存有有限相联度、多个层级和预取, 不能与理想模型逐一等同. 理想模型用于看清主要数量级, 不保证常数和每个矩阵尺寸下的性能. 时间局部性指一个值在短时间内反复用; 空间局部性指相邻地址在短时间内用到, 使一整条缓存行有效.

\originalsection{行列存储与访存方向}
行主序按每一行连续存储, 列主序按每一列连续存储. 设缓存行含 $B$ 个数. 顺序读 $n$ 个数约需 $n/B$ 次搬运; 若访问步长很大, 每次只用一条缓存行中的一个数, 可能接近 $n$ 次搬运. 运算次数相同, 带宽成本却差约 $B$ 倍.

Julia 的通常稠密数组采用列主序. 所以对矩阵逐元素遍历时, 内层循环宜让行指标改变, 沿一列连续走. C 中常见二维数组采用行主序, 适宜方向恰好相反. 语言和库的存储约定要实查, 不能凭矩阵的数学行列想象内存布局.

\begin{example}
一个 $n\times n$ 矩阵按列主序存储, 连续缓存行有 $B$ 个数. 比较逐列与逐行扫描全部元素的缺失次数, 假设步长和缓存容量使逐行访问不能保留有用的缓存行.
\end{example}
\begin{solution}
逐列扫描沿连续地址, 每条缓存行中的 $B$ 个数都被使用, 约为 $n^2/B$ 次. 逐行扫描每次跳过 $n$ 个数, 若缓存相联和容量使行间不能复用, 每个元素都可能触发一次搬入, 约为 $n^2$ 次. 这里的第二个数量级带有明确模型假设; 若整个矩阵能放入缓存, 或预取和跨行复用有效, 实测会不同.
\end{solution}
真实缓存的地址映射还可能造成某些特殊尺寸的冲突缺失. 因此时间曲线中的尖峰不一定是随机噪声. 测性能时应该记录尺寸、布局、线程、库与重复测量, 而不是只给一个最快时间.

\originalsection{分块矩阵乘法}
对 $n\times n$ 稠密矩阵, 把各矩阵分为 $b\times b$ 块. 以块指标 $I,J,K$ 写
\[
C_{IJ}\leftarrow C_{IJ}+A_{IK}B_{KJ}.
\]
若 $3b^2\le Z$, 三个活动块可同时进入快速存储. 这时每个读入的 $A$、$B$ 元素能服务一个块内的多个乘加, 不再只用一次就被替换.

\begin{proposition}\label{prop:blockcache}
在以一个数为移动单位的理想模型中, 取 $b$ 与 $\sqrt Z$ 同阶、并使三个活动块能装入快速存储. 经典分块矩阵乘法的数据移动量为
\[
O\left(n^2+\frac{n^3}{\sqrt Z}\right).
\]
若按含 $B$ 个数的连续缓存行计数, 且块布局与 tall-cache 等条件支持有效连续搬运, 相应量级为 $O(n^2/B+n^3/(B\sqrt Z))$.
\end{proposition}
\begin{proof}
忽略边缘不整块的常数, 共有 $(n/b)^3$ 次块乘法. 每次搬运至多常数个 $b^2$ 规模块, 总量为 $O(n^3/b)$. 初始读入与最终写出至少涉及 $O(n^2)$ 个元素. 取 $b=\Theta(\sqrt Z)$ 即得. 缓存行版本要求活动块的有效访问沿连续地址发生, 每次移动的一行能利用 $\Theta(B)$ 个数, 所以相同元素搬运界除以 $B$; 若这个布局条件失败, 不能直接这么除.
\end{proof}
这个算法需要知道适当块大小. 多级缓存需要多层分块; 处理器寄存器又构成更小的快速存储层. 高性能 BLAS 往往采用缓存块、数据打包和寄存器小核的结合, 不仅是给三重循环外面再加三层循环.

\originalsection{为什么平方根尺度自然出现}
$\sqrt Z$ 不是随意选的参数. 经典矩阵乘法的每个标量乘积 $a_{ik}b_{kj}$ 需要一对输入并更新输出位置 $(i,j)$. 把一次阶段中完成的乘积索引集合记为 $T\subset\{(i,k,j)\}$, 对应三个投影集合是用过的 $A$、$B$、$C$ 位置.

\begin{lemma}
若这些三个投影集合分别有 $a,b,c$ 个位置, 则 $|T|\le\sqrt{abc}$.
\end{lemma}
\begin{proof}
用 $0$--$1$ 矩阵 $M,N,P$ 分别指示三个位置集合. $T$ 的每个三元组都被这三个指标包含, 所以
\[
|T|\le\sum_{i,k,j}M_{ik}N_{kj}P_{ij}
=\inner{P}{MN}_F
\le\norm{P}_F\norm{MN}_F
\le\norm{P}_F\norm{M}_F\norm{N}_F=\sqrt{abc}.
\]
这里矩阵 Frobenius 内积的 Cauchy--Schwarz 不等式给第一界; $\norm{MN}_F\le\norm{M}_F\norm{N}_F$ 可逐列应用算子范数界与 $\norm{M}_2\le\norm{M}_F$ 得到.
\end{proof}
在经典不重复计算、显式维护各输出部分和的模型中, 把执行划成每段至多 $Z$ 次元素搬运. 一段初始快存有 $Z$ 个数, 搬运再引入至多 $Z$ 个, 所以输入位置的总规模为 $O(Z)$. 输出部分和可能从零新建, 不能只数搬入: 每个被更新的输出位置或者在段初已存在, 或在段中被搬入/搬出, 或在段末仍驻留快存. 由于不能丢弃待保留的部分和且不允许重计算, 三类数量均为 $O(Z)$, 输出位置也只有 $O(Z)$ 个. 因此三个投影位置集合的总规模为 $O(Z)$. 引理给该段最多 $O(Z^{3/2})$ 个不同乘积. 要完成 $n^3$ 个乘积, 需要 $\Omega(n^3/Z^{3/2})$ 段, 每完整段搬运 $Z$ 个数, 给出 $\Omega(n^3/\sqrt Z)$ 主项, 再加输入输出的 $\Omega(n^2)$.

这是经典矩阵乘法存储下界的主要结构解释. 若使用快速矩阵乘法、允许大量重计算, 或换成不同数据移动模型, 下界形式必须重新论证. 不能由此宣称任何可能的矩阵乘法算法都至少要 $n^3$ 次算术.

\originalsection{缓存无关的递归算法}
缓存无关算法不把 $Z$ 或 $B$ 写进程序. 对矩阵乘法, 把三个维度分成两半, 用分块恒等式递归计算, 直到足够小的基础问题. 对方阵, 一层产生八个半尺寸乘法, 总算术仍是 $\Theta(n^3)$.

当递归子问题的三个矩阵块达到能同时放入快存的尺寸 $b=\Theta(\sqrt Z)$ 时, 以下层级的访问都能在快存中完成. 该算法虽然不知道 $Z$, 递归树却自动经过这一尺度. 共有约 $(n/b)^3$ 个这种叶块, 每个搬运 $O(b^2)$ 个数, 因而也得到 $O(n^3/\sqrt Z+n^2)$ 元素搬运. 缓存行版本仍需恰当布局与容量条件.

对很不方的矩阵, 通常优先切最长维度, 避免产生极细长、不利于复用的块. 实际递归不应一直到 $1\times1$: 函数调用开销、寄存器利用与向量化可能压过理论优势. 原课程手册用较小块作为基例, 并指出即便消除 L1/L2 的明显容量断崖, 未优化的递归代码仍可能明显慢于专门 BLAS. 因而“缓存无关的渐近最优”不等于“任何尺寸的实际速度最快”.

\begin{center}
\begin{tikzpicture}[scale=.85,>=Stealth]
\draw (0,0) rectangle (2,2);\draw (1,0)--(1,2);\draw (0,1)--(2,1);
\node at (1,-.35) {$A$};
\draw (3,0) rectangle (5,2);\draw (4,0)--(4,2);\draw (3,1)--(5,1);
\node at (4,-.35) {$B$};
\draw (6,0) rectangle (8,2);\draw (7,0)--(7,2);\draw (6,1)--(8,1);
\node at (7,-.35) {$C$};
\node at (2.5,1) {$\times$};\node at (5.5,1) {$=$};
\node at (4,2.45) {$C_{11}=A_{11}B_{11}+A_{12}B_{21}$};
\end{tikzpicture}
\end{center}
图中的分块等式是精确矩阵恒等式. 不同求和树在浮点中可能给不同末位, 因此性能优化也应搭配误差与残差检查.

\originalsection{BLAS 层级与分解的块实现}
BLAS 一级操作是点积、向量加法等, 每个数据通常只参与少量运算; 二级操作是矩阵向量乘法与秩一更新; 三级操作是矩阵乘法和多右端三角求解, 能用大量数据复用提高算术强度. 算术强度是每搬运一个字节完成的操作数. 当强度低时, 运算往往受内存带宽限制; 当强度高且小核优化合理时, 更可能接近计算吞吐极限.

分块 LU 先分解一个窄面板, 再用三角求解和矩阵乘法更新尾块. 分块 Cholesky 对对角块分解, 解下方块, 再作对称秩更新. 分块 QR 把多个反射表示为 $I-VTV^*$, 利用矩阵乘法成批作用. 它们保留原有分解的数学目标, 但把大部分工作移到三级 BLAS. 若为重排计算改变了舍入次序, 应核查原算法的稳定性分析是否仍适用.

可比较的性能实验至少要先预热编译和库初始化, 固定线程设置, 用足够多次重复区分波动, 把输入生成和文件操作从计时中分离, 并检查输出正确性. 所谓某个程序“快十倍”应附具体环境和规模, 不能从历史课件曲线直接推给本次实现.

\begin{exercise}
若 $Z=3b^2$, 比较 $n\gg b$ 时逐元素矩阵乘法的 $O(n^3)$ 元素搬运与分块算法的 $O(n^3/b)$ 搬运. 当快存容量增加四倍时, 分块算法的主导搬运项如何改变?
\end{exercise}
\begin{solution}
$b$ 与 $\sqrt Z$ 同阶, 因而 $Z$ 增加四倍使合适的 $b$ 增加两倍, 主导 $n^3/b$ 项约减半. 这只是在同一理想模型中的主项比较, 不保证实际墙钟时间也恰好减半.
\end{solution}
\begin{remark}
本章整理 L11--12 的三份性能手册及逐周说明. 历史硬件、实验现象、寄存器复用和缓存无关观点来自 Johnson 手册; 分块上下界的中文推导和图为自编补充. 本册没有把原手册的旧性能曲线改标为新实测.
\end{remark}

% SOURCE: OCW L14--17 HTML; Persson Hessenberg (2p), QR I (4p), QR II (9p). Original proofs/supplements distinguished at chapter end.
\originalchapter{特征值问题与 QR 迭代}\label{ch:eigen}
\originalsection{Schur 形式与迭代的必要性}
求全部特征值不应通常先展开特征多项式再求根. 系数的形成可能丢掉矩阵结构, 从根到系数的敏感性也可能很差. 对首一多项式
$p(z)=z^n+c_{n-1}z^{n-1}+\cdots+c_0$, 伴随矩阵
\[
C=\begin{pmatrix}
0&0&\cdots&0&-c_0\\
1&0&\cdots&0&-c_1\\
0&1&\cdots&0&-c_2\\
\vdots&\vdots&\ddots&\vdots&\vdots\\
0&0&\cdots&1&-c_{n-1}
\end{pmatrix}
\]
满足 $\det(zI-C)=p(z)$, 可由行列式递推或特征向量递推核查. 这把一般多项式求根嵌入特征值问题. 一般五次以上多项式不存在统一的有限次根式公式, 所以不能期待像 LU 那样, 对所有矩阵以有限次四则运算与根式精确给出特征值. 数值算法采用迭代逼近; 这个代数限制不妨碍达到给定精度, 也不表示每个特殊矩阵都需要无限迭代.

\begin{theorem}[复 Schur 分解]\label{thm:schur}
对任意 $A\in\C^{n\times n}$, 存在酉矩阵 $Q$ 和上三角矩阵 $T$ 使 $A=QTQ^*$. $T$ 的对角元是 $A$ 的全部特征值. 若 $A$ Hermitian, $T$ 为实对角矩阵.
\end{theorem}
\begin{proof}
由代数基本定理, $A$ 有特征值 $\lambda$ 和单位特征向量 $q_1$. 扩充 $q_1$ 为酉矩阵 $Q_1$, 则 $Q_1^*AQ_1=\begin{pmatrix}\lambda&*\\0&B\end{pmatrix}$. 对 $B$ 按维数归纳, 在后 $n-1$ 个坐标内再作酉变换, 得上三角 $T$. 相似保持特征多项式, 上三角矩阵的特征值是对角元. 若 $A=A^*$, 则 $T=T^*$, 而同时上三角与 Hermitian 的矩阵必为实对角.
\end{proof}
Schur 分解对缺陷矩阵也存在, 比特征向量对角化更稳健. 例如 $\begin{pmatrix}1&1\\0&1\end{pmatrix}$ 只有一维特征空间, 不可对角化, 但本身就是 Schur 形式. 实矩阵若坚持实正交变换, 目标通常是实准上三角形式, 对角含 $1\times1$ 实根块与 $2\times2$ 复共轭根块.

特征多项式路线的另一个风险是系数扰动. 原课程提到的 Wilkinson 多项式为 $p(z)=\prod_{j=1}^{20}(z-j)$, 所有根本来都是简单整数. 若只对 $z^{19}$ 的系数加 $\eta$, 简单根 $r$ 的一阶变化满足
\[
\Delta r\approx-\frac{\eta r^{19}}{p'(r)},\qquad
p'(r)=(-1)^{20-r}(r-1)!(20-r)!\quad(r=1,\ldots,20).
\]
这个公式来自把 $(p+\eta z^{19})(r+\Delta r)=0$ 保留一阶项. 即使原根很整齐, $r^{19}/p'(r)$ 也可非常大. 这不是说所有多项式表示都同样病态, 而是说明把矩阵先转成某种系数表示, 可以引入不同于原矩阵问题的敏感性. 正交相似方法尽可能直接保持矩阵结构, 不经过这条额外路线.

\originalsection{特征值敏感性与残差}
\begin{theorem}[Bauer--Fike 型界]\label{thm:bauer}
若 $A=X\Lambda X^{-1}$ 可对角化, $\mu$ 是 $A+E$ 的任一特征值, 则
\[
\min_i|\mu-\lambda_i|\le\kappa_2(X)\norm{E}_2.
\]
若 $A$ 正规, 可取酉 $X$, 因而界为 $\norm{E}_2$.
\end{theorem}
\begin{proof}
若 $\mu$ 已是 $A$ 的特征值, 结论成立. 否则取 $(A+E)z=\mu z$, $z\ne0$, 得
$z=(\mu I-A)^{-1}Ez$.
所以 $1\le\norm{(\mu I-A)^{-1}}_2\norm{E}_2$. 由对角化,
\[
\norm{(\mu I-A)^{-1}}_2
\le\frac{\kappa_2(X)}{\min_i|\mu-\lambda_i|},
\]
合并即可. 正规矩阵有酉特征基, 条件数为 $1$.
\end{proof}
\begin{proposition}
设 $\lambda$ 为单特征值, $Ar=\lambda r$, $l^*A=\lambda l^*$, $l^*r\ne0$. 对光滑小扰动, 特征值的一阶变化为
\[
\Delta\lambda=\frac{l^*Er}{l^*r}+O(\norm{E}_2^2),
\qquad \kappa_{\lambda,\rm abs}=\frac{\norm{l}_2\norm{r}_2}{|l^*r|}.
\]
\end{proposition}
\begin{proof}
沿 $A(t)=A+tE$ 的局部单特征值分支对 $A(t)r(t)=\lambda(t)r(t)$ 求导, 左乘 $l^*$. 因 $l^*A=\lambda l^*$, 含 $r'(0)$ 的两项抵消, 得 $\lambda'(0)=l^*Er/(l^*r)$. 局部单根的光滑性来自特征多项式对 $\lambda$ 的非零导数与隐函数定理. Cauchy--Schwarz 给上界; 选择单位范数秩一扰动把 $r$ 方向映到 $l$ 方向, 可达到分子最大值, 得绝对条件数.
\end{proof}
对 Hermitian 矩阵可取 $l=r$ 为单位向量, 一阶绝对条件数为 $1$. 非正规矩阵的左右特征向量可能几乎正交, 特征值就会敏感. 缺陷点的一阶公式甚至不适用: $A=\begin{pmatrix}0&1\\0&0\end{pmatrix}$ 加 $E=\begin{pmatrix}0&0\\\eps&0\end{pmatrix}$ 后特征值为 $\pm\sqrt\eps$, 而 $\norm{E}_2=\eps$.

\begin{proposition}[近似特征对的后向解释]\label{prop:eigres}
对单位向量 $v$ 与标量 $\theta$, 记 $r=Av-\theta v$. 存在 $\norm{E}_2=\norm{r}_2$ 使 $(A+E)v=\theta v$, 具体可取 $E=-rv^*$. 若 $A$ Hermitian, 则还可直接得到 $\min_i|\theta-\lambda_i|\le\norm{r}_2$.
\end{proposition}
\begin{proof}
秩一构造立即给出 $(A-rv^*)v=\theta v$ 和范数等式. Hermitian 情形写 $v=\sum c_iq_i$, 则 $\norm{r}_2^2=\sum_i|\lambda_i-\theta|^2|c_i|^2\ge\min_i|\lambda_i-\theta|^2$.
\end{proof}
若取 $\theta=v^*Av$ 为实 Rayleigh 商, 则 $v^*r=0$, 还可选 Hermitian 扰动 $E=-rv^*-vr^*$, 在 $\Span(v,r)$ 上范数也是 $\norm{r}_2$. 这使近似对的解释保留 Hermitian 结构. 非正规矩阵中小残差仍是小后向误差, 但特征值前向距离还需要条件性控制.

\originalsection{Hessenberg 与三对角约化}
\begin{definition}
上 Hessenberg 矩阵满足 $h_{ij}=0$ 当 $i>j+1$, 即只允许主对角以下第一条副对角非零. Hermitian 的 Hessenberg 矩阵必为三对角.
\end{definition}
求 Schur 形式不能照搬 QR 分解的单边消元, 因为右乘恢复相似性时可能重新产生刚消掉的零. 更温和的目标是每列保留一个副对角元. 第 $k$ 步只在行 $k+1,\ldots,n$ 上作 Householder 反射, 把第 $k$ 列从第 $k+2$ 行以下消零; 同时右乘相同反射, 使变换为 $A\leftarrow H_k^*AH_k$.

\begin{proposition}
这个两侧反射不破坏前 $k-1$ 列已有的 Hessenberg 零. 最终得到 $A=QHQ^*$, $Q$ 酉, $H$ 上 Hessenberg. Hermitian 情形得到三对角矩阵.
\end{proposition}
\begin{proof}
左乘仅混合行 $k+1,\ldots,n$; 对前 $k-1$ 列, 这些位置已经为零, 混合后仍零. 右乘只混合列 $k+1,\ldots,n$, 不动前 $k$ 列, 因而不破坏本步和此前列的消零. 逐步归纳得 Hessenberg 结构. 两侧乘积为相似变换, Hermitian 性保持; 其下方零的共轭位置在上方也为零, 因而三对角.
\end{proof}
对实稠密矩阵, 一般 Hessenberg 约化主导约为 $\tfrac{10}{3}n^3$ 次操作; 利用对称结构的三对角约化约为 $\tfrac43n^3$. 成本只支付一次, 后续迭代显著便宜. 正交相似还允许像命题 \ref{prop:orthaccum} 那样累积误差; 标准实现的约化是后向稳定的, 但需要区分存储的近似 $Q$ 和分析中的精确酉邻近矩阵.

\originalsection{幂法与移位反迭代}
幂法从单位向量 $v_0$ 出发, 反复算 $w=Av_k$, $v_{k+1}=w/\norm{w}_2$, 并用 $\theta_k=v_k^*Av_k$ 估计特征值. 每一步只需一次矩阵向量乘法, 适合寻找主导特征方向. $w=0$ 时须停止或换初始向量, 不能继续除以零.

\begin{theorem}[Hermitian 幂法收敛]\label{thm:power}
设 $A$ Hermitian, $|\lambda_1|>|\lambda_2|\ge\cdots$, $\lambda_1\ne0$, 且 $q_1^*v_0\ne0$. 忽略单位复相位或实符号, $v_k$ 趋向 $q_1$, 其夹角正切为 $O((|\lambda_2/\lambda_1|)^k)$, Rayleigh 商误差为 $O((|\lambda_2/\lambda_1|)^{2k})$.
\end{theorem}
\begin{proof}
写 $v_0=\sum c_iq_i$, $c_1\ne0$. 归一化不改变方向, 故 $v_k$ 与
$q_1+\sum_{i>1}(c_i/c_1)(\lambda_i/\lambda_1)^kq_i$ 同方向. 正交分解给夹角正切不超过
$(\sum_{i>1}|c_i/c_1|^2)^{1/2}|\lambda_2/\lambda_1|^k$.
对任意单位 $v=\sum d_iq_i$, Rayleigh 商误差为
$\theta-\lambda_1=\sum_{i>1}(\lambda_i-\lambda_1)|d_i|^2$,
其绝对值不超过 $\max_i|\lambda_i-\lambda_1|\sin^2\angle(v,q_1)$, 得平方速度.
\end{proof}
若 $|\lambda_1|=|\lambda_2|$, 或初始向量对主导空间分量为零, 这个定理的结论不能保证. 例 $A=\diag(1,-1)$ 的一般初始向量会来回变化, 不趋一个方向. 缺陷和非正规矩阵需要另看 Jordan 结构和特征向量条件性.

移位反迭代反复解 $(A-\mu I)w=v_k$, 再归一化, 等价于对 $(A-\mu I)^{-1}$ 作幂法. 它优先放大离 $\mu$ 最近的特征值. 固定移位可复用 LU; 不应显式形成逆矩阵. 若最近特征值为唯一的 $\lambda_j$, 收敛因子为
$\max_{i\ne j}|(\lambda_j-\mu)/(\lambda_i-\mu)|$.
移位越接近目标, 理论速度越快, 但线性系统更病态, 求解误差也要监控.

\originalsection{Rayleigh 商迭代}
Rayleigh 商 $\rho(v)=v^*Av/(v^*v)$ 也是使 $\norm{Av-\alpha v}_2$ 最小的标量 $\alpha$: 把 $Av$ 投影到 $v$ 的张成空间即可. Hermitian 情形, 在单位球面上的驻点是特征向量; 局部特征值误差对方向误差为二阶.

Rayleigh 商迭代每次取 $\mu_k=\rho(v_k)$, 解 $(A-\mu_kI)w=v_k$, 归一化得到下一向量.
\begin{proposition}[局部三次收敛]
设 $A$ Hermitian, $\lambda_j$ 单, 与其余谱间隙为 $g=\min_{i\ne j}|\lambda_i-\lambda_j|>0$. 若 $v_k$ 已足够接近 $q_j$, 且移位系统可逆, 则
\[
\tan\angle(v_{k+1},q_j)\le C\tan^3\angle(v_k,q_j)
\]
对固定邻域内某个常数 $C$ 成立.
\end{proposition}
\begin{proof}
写 $v=c_jq_j+\sum_{i\ne j}c_iq_i$, 令 $t=(\sum_{i\ne j}|c_i|^2)^{1/2}/|c_j|$. Rayleigh 商满足
$|\mu-\lambda_j|\le M\sum_{i\ne j}|c_i|^2\le Mt^2$,
其中 $M=\max_i|\lambda_i-\lambda_j|$. 足够接近时该差小于 $g/2$, 所以 $|\lambda_i-\mu|\ge g/2$ ($i\ne j$). 反迭代后各非目标系数与目标系数的比例为
$ (c_i/c_j)(\lambda_j-\mu)/(\lambda_i-\mu)$.
由正交和得新夹角正切不超过 $(2M/g)t^3$, 即所述界.
\end{proof}
这个结论是局部的, 不保证任意初始向量收敛到预定特征值. 若已经得到精确特征向量, 移位系统奇异, 应以残差已为零而结束. 某些尚非特征向量的情形也可能恰好使 Rayleigh 商等于一个特征值, 简单求逆步骤无法执行; 实现需检测和处理, 而不是把证明中的“可逆”略去.

\originalsection{未移位 QR 与同时迭代}
QR 特征值算法从 $A_0=A$ 出发, 每次作 $A_{k-1}=Q_kR_k$, 再置 $A_k=R_kQ_k$. 因 $A_k=Q_k^*A_{k-1}Q_k$, 每一步保持特征值. 令 $\mathcal Q_k=Q_1\cdots Q_k$, 则 $A_k=\mathcal Q_k^*A\mathcal Q_k$.

\begin{proposition}\label{prop:qriter}
设 $A$ 可逆, 各 QR 分解按正实对角的相容约定完成. 则
\[
A^k=\mathcal Q_k\mathcal R_k,\qquad
\mathcal R_k=R_kR_{k-1}\cdots R_1.
\]
因此 QR 迭代等价于每步把 $A$ 作用到当前正交基后再正交化的同时迭代, 从单位基开始.
\end{proposition}
\begin{proof}
$k=1$ 成立. 若 $A^{k-1}=\mathcal Q_{k-1}\mathcal R_{k-1}$, 则
$A\mathcal Q_{k-1}=\mathcal Q_{k-1}A_{k-1}=\mathcal Q_{k-1}Q_kR_k=\mathcal Q_kR_k$.
相乘得 $A^k=\mathcal Q_kR_k\mathcal R_{k-1}$, 完成归纳. 这个等式也表明同一步的 $\mathcal Q_k$ 正是 $A\mathcal Q_{k-1}$ 的正交化结果.
\end{proof}
直接形成很高次 $A^k$ 会让主导方向淹没其他方向, 后续一次 QR 难以恢复它们. 同时迭代每步正交化保留不同方向, 以酉变换避免这个额外的数值损失.

为明确收敛条件, 继续假设 $A$ 可逆, 且 $A=V\Lambda V^*$ Hermitian, 特征值模严格降序. 对每个 $1\le j<n$, 取前 $j$ 个初始坐标向量组成 $E_j$, 写 $V^*E_j=\begin{pmatrix}C_1\\C_2\end{pmatrix}$, 假设 $C_1$ 可逆. 由
\[
V^*A^kE_j=\begin{pmatrix}\Lambda_1^kC_1\\\Lambda_2^kC_2\end{pmatrix}
\]
知其列空间是主导坐标空间上方的图, 下方图矩阵为
$F_k=\Lambda_2^kC_2C_1^{-1}\Lambda_1^{-k}$.
故 $\norm{F_k}_2\le\norm{C_2C_1^{-1}}_2|\lambda_{j+1}/\lambda_j|^k$.
对每个 $j$ 都满足该非退化条件时, 各级嵌套列空间趋向主导特征空间, 正交列依次趋向对应特征向量, 因而 $A_k$ 趋向对角形式. 这给出 QR 收敛的一个充分条件和证明; 不把未移位 QR 描述为对任意矩阵必收敛.

对可逆实对称 $A$, 由 $A^k=\mathcal Q_k\mathcal R_k$ 和对称性,
$A^{-k}=\mathcal Q_k\mathcal R_k^{-T}$.
把列顺序反转的置换记为 $P$, 则
$A^{-k}P=(\mathcal Q_kP)(P\mathcal R_k^{-T}P)$,
第二因子仍上三角. 因而最后一个 $\mathcal Q_k$ 列也可看作对反转初始基实施反迭代的首方向. 这一关系解释了移位通常从尾部目标特征值着手.

\originalsection{移位、退耦与实际 QR}
移位 QR 作
\[
A_{k-1}-\mu_kI=Q_kR_k,\qquad
A_k=R_kQ_k+\mu_kI.
\]
仍有 $A_k=Q_k^*A_{k-1}Q_k$. 相似关系与归纳还给
\[
\prod_{j=1}^k(A-\mu_jI)=\mathcal Q_k\mathcal R_k,
\]
各因子都是 $A$ 的多项式而可交换. 因此移位同时迭代用多项式滤波替代简单的 $A^k$. 一个合适移位可迅速压制尾部方向的耦合.

对实对称三对角矩阵, 常取尾对角元素 $a_{nn}$ 作 Rayleigh 移位. 但这个选择可能保留不利对称性. 例 $T=\begin{pmatrix}0&1\\1&0\end{pmatrix}$ 的尾对角为 $0$, 未移位 QR 不能消掉耦合. Wilkinson 移位使用末尾 $2\times2$ 块
$\begin{pmatrix}a&b\\b&d\end{pmatrix}$ 中离 $d$ 较近的特征值. 令 $\delta=(a-d)/2$, 可稳定写为
\[
\mu=d-\operatorname{sign}(\delta)
\frac{b^2}{|\delta|+\sqrt{\delta^2+b^2}},
\]
$\delta=0$ 时固定取一个符号以打破平局; $b=0$ 则直接退耦, 不计算 $0/0$. 这个写法避免 $\delta$ 与平方根的相近数相减. 在上例中可选 $\mu=1$ 或 $-1$, 对应已知特征值并使块迅速分离.

原手册对实对称三对角 QR 的 Wilkinson 移位给出全局收敛与至少二次的边界说明. 其严格全局证明和特殊退化情形超出本章局部迭代证明的范围; 这里采用这一专门算法结果, 不推广到一般非正规矩阵, 也不承诺每一步或每个输入必为三次速度. 本章已证明的三次结论是带可逆和谱间隙条件的 Rayleigh 商反迭代局部结论.

当副对角 $h_{i+1,i}$ 已足够小, 可把它设为零, 分离前后两个块分别迭代, 称为退耦或 deflation. 对称三对角中同时置对称位置为零, 此操作就是加一个只含该连接项的小 Hermitian 扰动, 范数为 $|h_{i+1,i}|$. 阈值应以附近对角尺度、机器精度和安全最小数设定, 不宜只用固定绝对 $10^{-12}$.

Hessenberg QR 每步可用 $O(n)$ 个 Givens 旋转, 每个更新 $O(n)$ 元素, 总计 $O(n^2)$. 对称三对角可用隐式移位和追赶凸起保持带状, 仅求特征值时每步约 $O(n)$; 若同时累计稠密特征向量矩阵, 成本更高. 显式构造稠密 $Q$ 再乘 $RQ$ 会失去这个复杂度优势. 实际库还采用多移位、分治等方法, 本章不展开所有工程细节.

\originalsection{SVD 的双对角约化与可靠输出}
对长方矩阵, 交替作左、右 Householder 变换可约化为上双对角矩阵 $B$, 即 $A=UBV^*$. 第一列左消元后, 在第一行剩余条目右消元, 再对尾块重复. 左变换处理列, 右变换处理行, 同时保留已有零的形状. 奇异值在这些酉左右变换下不变, 后续对双对角矩阵作专用 QR 或其他迭代. 这就是 Golub--Kahan 双对角化的稠密结构思想; 不能把它与直接形成 $A^*A$ 的算法混同.

对于 Hermitian 特征分解, 合理检查包含 $\norm{AQ-Q\Lambda}$ 与 $\norm{Q^*Q-I}$ 两项. 若算法后向误差为 $\norm{E}_2=O(u)\norm{A}_2$, 单个特征值的绝对误差也是这个量级; 很小的特征值相对误差却可能很大. 对接近的多个特征值, 单个特征向量可能不稳定, 但整个不变子空间仍可能准确. 正规矩阵的特征值绝对可靠性不自动等于每个特征向量相对可靠性.

\begin{exercise}
设 Hermitian 矩阵有单位特征向量 $q_j$, $v=\cos\theta\,q_j+\sin\theta\,w$, $w\perp q_j$, $\norm{w}_2=1$. 证明 $\rho(v)-\lambda_j=\sin^2\theta(w^*Aw-\lambda_j)$.
\end{exercise}
\begin{solution}
展开 $v^*Av$, 交叉项为零, 因 $q_j^*Aw=\lambda_jq_j^*w=0$. 主项为 $\lambda_j\cos^2\theta+(w^*Aw)\sin^2\theta$, 减去 $\lambda_j$ 即得.
\end{solution}
\begin{exercise}
若 $A$ Hermitian, 单位 $v$ 的 Rayleigh 商为 $\theta$, 残差为 $r$. 假设除目标特征值 $\lambda_j$ 外均满足 $|\lambda_i-\theta|\ge g>0$. 证明 $\sin\angle(v,q_j)\le\norm{r}_2/g$.
\end{exercise}
\begin{solution}
写 $v=\sum c_iq_i$, 有 $\norm{r}_2^2=\sum|\lambda_i-\theta|^2|c_i|^2\ge g^2\sum_{i\ne j}|c_i|^2$. 最后一和就是夹角正弦的平方. 谱间隙小的时候, 残差可以很小而单个方向误差仍较大.
\end{solution}
\begin{remark}
本章对应 L14--17, 核心资料为 Persson 的 Hessenberg/Tridiagonal Reduction 与 The QR Algorithm I/II. Schur 归纳、QR 同时迭代、反迭代、移位和历史手册的结构被保留并重新叙述. 特征值扰动证明、Rayleigh 局部三次推导、逐级子空间收敛与例题为自编补充. 全局 Wilkinson 收敛和完整库实现只在明确限定下作为引用结果, 不伪称本章已证明全部高性能特征值软件细节.
\end{remark}

% SOURCE: MIT 18.335J Spring 2019, Week 6 L17, Week 7 L18--19;
% Johnson, Why restarting Arnoldi/Lanczos is not trivial, March 21, 2019.
% ADAPTATION: self-authored supplemental explanation; not a translation of absent handouts.
\originalchapter{Krylov 子空间、Arnoldi 与 Lanczos}\label{ch:krylov}
\originalsection{从反复乘矩阵到保留整个子空间}
大型线性代数问题的困难常常不是一条公式不会写, 而是不能存储或分解矩阵. 设 $A\in\C^{n\times n}$. 若 $A$ 有 $s$ 个非零元, 一次矩阵向量积通常只需 $O(s)$ 次运算; 对某些积分算子或结构矩阵, 甚至可以在不显式构造 $A$ 的情况下计算 $Av$. 反之, 一个稠密的 $n\times n$ 分解可能需要 $O(n^2)$ 存储和 $O(n^3)$ 运算. 因此我们要把问题改写成一系列矩阵向量积和低维计算.

幂法从 $b$ 开始生成 $b,Ab,A^2b,\ldots$, 每一步通常只保留最新的方向. 这适合某些占优势的特征向量, 却丢掉了前面向量已经携带的信息. 若目标是求线性方程的解, 或多个特征方向, 比较自然的改进是保留这些向量的\emph{整个线性包}, 再在其中选最合适的近似.

\begin{definition}
对 $A\in\C^{n\times n}$ 和 $b\ne0$, 第 $k$ 个 Krylov 子空间定义为
\[
\mathcal K_k(A,b)=\Span\{b,Ab,\ldots,A^{k-1}b\},\qquad \mathcal K_0(A,b)=\{0\}.
\]
等价地, $v\in\mathcal K_k(A,b)$ 当且仅当 $v=q(A)b$, 其中 $q$ 是次数至多 $k-1$ 的多项式.
\end{definition}
这一等价表达非常重要. 它说明迭代方法并不是任意搜索 $n$ 维空间, 而是在不断增大的多项式族中选择系数. 后面的收敛分析因此会成为多项式逼近问题.

\begin{proposition}\label{prop:krylov-grade}
令 $d$ 是使 $b,Ab,\ldots,A^db$ 线性相关的最小正整数. 则 $\dim\mathcal K_k(A,b)=k$ 对 $1\le k\le d$ 成立, 且
\[
\mathcal K_{d+1}(A,b)=\mathcal K_d(A,b),\qquad
A\mathcal K_d(A,b)\subseteq\mathcal K_d(A,b).
\]
存在唯一首一、次数为 $d$ 的多项式 $m_b$ 使 $m_b(A)b=0$. 若 $A$ 可逆, 则 $m_b(0)\ne0$, 且 $A^{-1}b\in\mathcal K_d(A,b)$.
\end{proposition}
\begin{proof}
按 $d$ 的最小性, 前 $d$ 个向量线性无关. 第一次相关关系中 $A^db$ 的系数必须非零, 否则较早已经相关. 因此可唯一写成
\[
A^db=-\sum_{j=0}^{d-1}c_jA^jb,
\qquad m_b(z)=z^d+\sum_{j=0}^{d-1}c_jz^j.
\]
$A$ 把前 $d-1$ 个生成向量送到下一个, 把最后一个送到上述线性组合, 所以子空间不变. 唯一性也来自前 $d$ 个向量的线性无关.

若 $c_0=0$ 且 $A$ 可逆, 则把 $m_b(A)b=0$ 乘以 $A^{-1}$ 就得到次数 $d-1$ 的首一消去关系, 矛盾. 故 $c_0\ne0$. 再次乘以 $A^{-1}$ 并整理,
\[
A^{-1}b=-\frac1{c_0}\left(A^{d-1}b+\sum_{j=1}^{d-1}c_jA^{j-1}b\right)
\in\mathcal K_d(A,b).
\]
\end{proof}
$d$ 称为 $b$ 相对于 $A$ 的级数或 grade. 它可以远小于 $n$. 例如 $b$ 恰为特征向量时 $d=1$; 若 $A$ 为 Hermitian 矩阵, 且 $b$ 只在 $d$ 个不同特征值的特征子空间上有非零分量, 则 grade 正好是这个数目.

\begin{proposition}[移位不改变 Krylov 子空间]\label{prop:krylov-shift}
对任意 $\mu\in\C$ 和 $k\ge1$, 有
$\mathcal K_k(A+\mu I,b)=\mathcal K_k(A,b)$.
\end{proposition}
\begin{proof}
二项式展开表明 $(A+\mu I)^jb$ 是 $b,Ab,\ldots,A^jb$ 的线性组合, 因而一侧包含于另一侧. 把 $A$ 写成 $(A+\mu I)-\mu I$, 同样得到反向包含.
\end{proof}
这只是\emph{搜索空间}的等式. 用不同准则从同一空间中选出的近似解仍可能不同, 所以不能由此推断求解 $Ax=b$ 与 $(A+\mu I)x=b$ 同样容易.

\originalsection{Arnoldi 如何生成稳定的基}
直接把 $b,Ab,\ldots$ 放在矩阵中有两个问题: 各列长度可能相差巨大, 各列也可能几乎平行. 即使这些列在精确算术中独立, 浮点数中也可能严重病态. Arnoldi 的关键是每生成一个新方向就正交化, 而不是等全部幂次向量生成后再处理.

取 $v_1=b/\norm{b}_2$. 第 $j$ 步先计算 $w=Av_j$, 再用改进 Gram--Schmidt 依次执行
\begin{equation}\label{eq:arnoldi-step}
 h_{ij}=v_i^*w,\quad w\leftarrow w-h_{ij}v_i\quad(i=1,\ldots,j),\qquad
 h_{j+1,j}=\norm{w}_2,\quad v_{j+1}=\frac{w}{h_{j+1,j}}.
\end{equation}
最后的除法只在 $h_{j+1,j}\ne0$ 时进行. 在精确算术中, 每次更新后的 $w$ 都垂直于已经处理的基向量; 后面的减法不会破坏这一点, 因为基向量彼此正交.

记 $V_k=[v_1,\ldots,v_k]$, $H_k=(h_{ij})_{i,j=1}^k$, 并把 $H_k$ 下面补一行得到
\[
\overline H_k=
\begin{bmatrix}H_k\\ h_{k+1,k}e_k^*\end{bmatrix}\in\C^{(k+1)\times k},
\]
这里 $e_k$ 是 $\C^k$ 的最后一个坐标向量. 所有尚未定义且位于第一条次对角线以下的 $h_{ij}$ 取零.

\begin{theorem}[Arnoldi 关系]\label{thm:arnoldi}
若前 $k$ 步未发生 $h_{j+1,j}=0$, 则 $V_{k+1}^*V_{k+1}=I$, $\range(V_j)=\mathcal K_j(A,b)$, 且
\begin{equation}\label{eq:arnoldi-relation}
 AV_k=V_{k+1}\overline H_k
 =V_kH_k+h_{k+1,k}v_{k+1}e_k^*,\qquad H_k=V_k^*AV_k.
\end{equation}
$H_k$ 是上 Hessenberg 矩阵.
\end{theorem}
\begin{proof}
归纳假设 $V_j$ 是 $\mathcal K_j$ 的正交基. 正交化后得到
\[
Av_j=\sum_{i=1}^{j}h_{ij}v_i+h_{j+1,j}v_{j+1}.
\]
因此 $v_{j+1}\perp\range(V_j)$, 且 $v_{j+1}\in\mathcal K_{j+1}$. 因为它非零, 所以 $j+1$ 个向量独立. 另一方面 $\dim\mathcal K_{j+1}\le j+1$, 故这些向量恰为其基. 初始 $v_1$ 显然满足归纳起点.

把每一步的等式并列就得矩阵关系. 第 $j$ 列只出现 $v_1,\ldots,v_{j+1}$, 所以 $h_{ij}=0$ 对 $i>j+1$ 成立. 左乘 $V_k^*$ 后, 最后一项因 $V_k^*v_{k+1}=0$ 消失, 得到 $H_k=V_k^*AV_k$.
\end{proof}
Arnoldi 因而是\emph{部分 Hessenberg 化}: 只构造整个矩阵分解中当前需要的前几列. 它与先对所有幂次列做精确 Gram--Schmidt 得到相同的子空间, 但避免显式计算很高的幂次.

\begin{proposition}[精确中止]\label{prop:arnoldi-breakdown}
若第 $j$ 步 $h_{j+1,j}=0$, 则 $\mathcal K_j$ 是 $A$-不变子空间, 并且 $AV_j=V_jH_j$. 若 $A$ 可逆, 在此子空间中求解 $H_jy=\norm{b}_2e_1$ 就得到 $A^{-1}b=V_jy$.
\end{proposition}
\begin{proof}
前 $j-1$ 列的乘积已经落在 $\range(V_j)$; 中止的第 $j$ 列也完全落在该空间. 因而整个空间不变. $A$ 在该空间上的限制是单射, 有限维下也是满射, 所以 $H_j$ 可逆. 由 $V_j\norm{b}_2e_1=b$ 便得到所述解.
\end{proof}
这种中止常称为幸运中止. 数值实现中, 一个很小但非零的 $h_{j+1,j}$ 不能直接当作精确不变性证明; 应检查得到的残差与所需容差.

若一次 $Av$ 的代价为 $C_A$, 则 $k$ 步矩阵乘法代价为 $kC_A$. 第 $j$ 步需与 $j$ 个长度 $n$ 的向量作内积和更新, 所以总正交化代价为 $O(nk^2)$, 基的存储为 $O(nk)$. 这正是方法运行很久后不得不考虑重启的原因.

\originalsection{Rayleigh--Ritz 近似与可计算的残差}
有了一个正交基, 大型特征值问题可以投影成小问题. 我们搜索 $x=V_ky\ne0$ 和 $\theta\in\C$, 并要求特征残差 $Ax-\theta x$ 垂直于搜索空间. 此条件叫 Galerkin 条件:
\[
V_k^*(AV_ky-\theta V_ky)=0
\quad\Longleftrightarrow\quad H_ky=\theta y.
\]
小矩阵 $H_k$ 的特征值称为 Ritz 值, 相应的 $V_ky$ 称为 Ritz 向量. 这是一条选近似的规则, 并不是说每个 Ritz 值都已接近所需特征值.

\begin{proposition}[Ritz 残差]\label{prop:ritz-residual}
若 $H_ky=\theta y$, $\norm{y}_2=1$, 则 $x=V_ky$ 为单位向量, 且
\[
Ax-\theta x=h_{k+1,k}v_{k+1}e_k^*y,
\qquad \norm{Ax-\theta x}_2=|h_{k+1,k}|\,|e_k^*y|.
\]
若 $A$ 为正规矩阵, 则
$\min_{\lambda\in\operatorname{spec}(A)}|\theta-\lambda|\le\norm{Ax-\theta x}_2$.
\end{proposition}
\begin{proof}
把 Arnoldi 关系乘以 $y$ 得到残差式. 正交基保持长度, 因而 $\norm{x}_2=1$. 若 $A=U\Lambda U^*$ 是正规矩阵的酉对角化, 写 $x=Uc$, 则
\[
\norm{Ax-\theta x}_2^2
=\sum_i|\lambda_i-\theta|^2|c_i|^2
\ge\left(\min_i|\lambda_i-\theta|^2\right)\sum_i|c_i|^2.
\]
最后的和为 $1$, 开平方即得结论.
\end{proof}
残差小能保证一个\emph{近邻特征值}, 却未必能识别目标特征值, 也未必能保证特征向量误差小. 若 $A$ Hermitian, 目标是简单特征值 $\lambda_\ell$, 且 $|\lambda_i-\theta|\ge\delta>0$ 对所有 $i\ne\ell$ 成立, 则同一展开给出
\[
\sum_{i\ne\ell}|c_i|^2\le\frac{\norm{Ax-\theta x}_2^2}{\delta^2}.
\]
左侧正是 $x$ 与目标特征方向夹角的正弦平方. 这里特征值间隔是不可忽略的条件.

若 $A=A^*$, 则 $H_k=H_k^*$. 对任意单位 $y$, $y^*H_ky=(V_ky)^*A(V_ky)$, 所以
\[
\lambda_{\max}(H_k)=\max_{0\ne x\in\mathcal K_k}\frac{x^*Ax}{x^*x},\qquad
\lambda_{\min}(H_k)=\min_{0\ne x\in\mathcal K_k}\frac{x^*Ax}{x^*x}.
\]
由于搜索空间嵌套, 最大 Ritz 值随 $k$ 单调不减且不超过 $\lambda_{\max}(A)$, 最小 Ritz 值单调不增且不小于 $\lambda_{\min}(A)$. 这提供了对 Hermitian 极端特征值近似的清晰解释.

\originalsection{Lanczos 的三项递推为何成立}
Arnoldi 的内积数随迭代步数增加. 对 Hermitian 矩阵, 上 Hessenberg 结构与 Hermitian 对称性同时存在, 因而 $H_k$ 必为实对角、共轭对称的三对角矩阵. 在本算法取 $h_{j+1,j}\ge0$ 的约定下, 次对角线和超对角线都是同一个非负实数.

令 $\alpha_j=v_j^*Av_j\in\R$, $\beta_j=h_{j+1,j}\ge0$, 并设 $\beta_0=0$, $v_0=0$. 精确算术中的递推为
\begin{equation}\label{eq:lanczos-step}
 w=Av_j-\beta_{j-1}v_{j-1},\qquad
 \alpha_j=v_j^*w,\qquad
 w\leftarrow w-\alpha_jv_j,\qquad
 \beta_j=\norm{w}_2,\quad v_{j+1}=w/\beta_j.
\end{equation}
因此
\[
T_k=\begin{bmatrix}
\alpha_1&\beta_1&&\\
\beta_1&\alpha_2&\ddots&\\
&\ddots&\ddots&\beta_{k-1}\\
&&\beta_{k-1}&\alpha_k
\end{bmatrix},\qquad
AV_k=V_kT_k+\beta_kv_{k+1}e_k^*.
\]

\begin{proof}[三项递推的证明]
对 $i\le j-2$, Hermitian 性给出 $v_i^*Av_j=(Av_i)^*v_j$. 已完成的第 $i$ 步表明 $Av_i\in\Span\{v_1,\ldots,v_{i+1}\}$, 而 $i+1\le j-1$, 所以该内积为零. 对 $i=j-1$, 第 $j-1$ 步给出 $v_{j-1}^*Av_j=\beta_{j-1}$. 因而在第 $j$ 步, 只需消去 $v_{j-1}$ 和 $v_j$ 两个方向; 其他方向的系数在精确算术中本来就为零.
\end{proof}
三项递推来自对称性和已经建立的正交关系共同作用, 不能只凭矩阵是稀疏就使用它. 对一般非 Hermitian 矩阵, 删除较早方向的正交化会改变算法.

\begin{example}
对
\[
A=\begin{bmatrix}2&-1&0\\-1&2&-1\\0&-1&2\end{bmatrix},\qquad v_1=e_1,
\]
执行 Lanczos, 并解释两步 Ritz 值的意义.
\end{example}
\begin{solution}
第一步 $Av_1=2e_1-e_2$, 所以 $\alpha_1=2$, $\beta_1=1$, $v_2=-e_2$. 第二步
\[
Av_2=e_1-2e_2+e_3,\qquad
Av_2-\beta_1v_1=2v_2+e_3,
\]
故 $\alpha_2=2$, $\beta_2=1$, $v_3=e_3$. 第三步得到 $\alpha_3=2$, $\beta_3=0$. 因此
$T_2=\begin{bmatrix}2&1\\1&2\end{bmatrix}$ 的 Ritz 值为 $1,3$.

原矩阵的特征值为 $2-\sqrt2,2,2+\sqrt2$. 两步的最小值 $1$ 和最大值 $3$ 分别从内侧逼近两个极端特征值. 对 $T_2$ 的单位特征向量, 最后分量绝对值都是 $1/\sqrt2$, 所以两个 Ritz 残差均为 $1/\sqrt2$. 第三步空间已等于整个 $\R^3$, 残差项消失, 三个特征值全都精确出现.
\end{solution}
只求 $T_k$ 的特征值时, 精确 Lanczos 可只保留当前和前一向量, 加上 $\alpha_j,\beta_j$ 的标量序列, 存储为 $O(n+k)$. 若要恢复 Ritz 向量, 就要存储基, 或在第二遍重新生成基. 这些节省还会受到舍入误差的限制.

\originalsection{浮点数中的正交性与幽灵特征值}
在浮点运算中, 三项递推产生的 $\widehat v_i$ 不会永久满足 $\widehat v_i^*\widehat v_j=0$. 误差会被后续乘法传播, 而 Lanczos 没有主动消去所有旧方向. 当某个特征方向已经收敛时, 新向量中极小的该方向分量可能再次被放大. 结果是三对角矩阵的谱中出现同一特征值的额外近似副本, 通常叫\emph{幽灵特征值}. 它们不能直接解释成原矩阵特征值的真实重数.

这里有两个不同的检查对象. 一个是局部递推缺陷
\[
A\widehat V_k-\widehat V_k\widehat T_k
-\widehat\beta_k\widehat v_{k+1}e_k^*,
\]
另一个是全局正交缺陷 $\widehat V_k^*\widehat V_k-I$. 前者小不自动保证后者小; 没有接近酉的基, 也不能直接把精确 Ritz 残差公式当成完整误差证书. 最稳妥的验收仍是对最终 $\widehat x$ 计算 $A\widehat x-\widehat\theta\widehat x$.

常用的补救是全重正交化、选择性重正交化或部分重正交化. 全重正交化在每步把 $w$ 与全部旧基重新正交化, 可靠但失去理想的线性存储优势. 选择性方案主要防止向量重新进入已经收敛的特征方向. 部分方案则用便宜的估计监控正交缺陷, 超过阈值才重正交化. 阈值和误差传播需要专门分析, 不能用一条固定经验常数替代所有问题的判断.

Arnoldi 本身显式保存全部基并作正交化, 因而较少出现上述无控制的旧方向再进入, 但一次改进 Gram--Schmidt 也不是精确正交化. 若新向量几乎落在旧空间中, 投影减法会强烈消去有效数字. 实务中可对 $w$ 再做一次正交化, 并把第二遍内积加到原有 $h_{ij}$ 上, 从而保持同一个 Arnoldi 关系的数值版本. 不能只修改向量而忘记修改投影系数.

\originalsection{为什么重启不能只截取几个 Ritz 向量}
假设已运行到 $m$ 步, 记
\begin{equation}\label{eq:arnoldi-before-restart}
AV_m=V_mH_m+f_me_m^*,\qquad f_m\perp\range(V_m).
\end{equation}
为了缩减存储, 想保留一个 $k<m$ 维空间 $\widehat V_k=V_mY$, 其中 $Y^*Y=I$. 可以把 $Y$ 补成酉矩阵 $[Y,Y_\perp]$. 将原关系乘以 $Y$, 并把 $H_mY$ 分解为两个坐标块, 得
\begin{equation}\label{eq:restart-full}
A\widehat V_k
=\widehat V_k(Y^*H_mY)
+V_mY_\perp(Y_\perp^*H_mY)
+f_me_m^*Y.
\end{equation}
能继续普通 Arnoldi 的标准形式应是
\[
A\widehat V_k=\widehat V_k\widehat H_k+\widehat f_ke_k^*,
\qquad \widehat f_k\perp\range(\widehat V_k),
\]
其中 $\widehat H_k$ 上 Hessenberg, 且外部残差只出现在最后一列. 这比单纯保留一个好的近似子空间要求更多.

若选 $H_m$ 的 $k$ 个线性无关特征向量组成 $Z$, 作薄 QR 分解 $Z=YR$, 并设 $H_mZ=Z\Theta$, 则
\[
H_mY=YR\Theta R^{-1},\quad
Y_\perp^*H_mY=0,\quad
Y^*H_mY=R\Theta R^{-1}.
\]
$R\Theta R^{-1}$ 是上三角矩阵, 因而当然也是上 Hessenberg. 看似两大困难都消失了, 但最后一项仍是
\[
f_m(e_m^*Y).
\]
行向量 $e_m^*Y$ 一般在多个位置非零, 而不是某个标量乘以 $e_k^*$. 所以多个保留向量都带有外部残差, 不能把它们原样当作已经完成的普通 Arnoldi 前 $k$ 步.

\begin{example}
设保留两个正交 Ritz 向量后有 $e_m^*Y=(a,b)$, 其中 $a,b\ne0$. 即使 $Y^*H_mY=\diag(\theta_1,\theta_2)$, 新关系的第一列仍为
$A\widehat v_1=\theta_1\widehat v_1+a f_m$.
普通两步 Arnoldi 的第一列应完全落在 $\Span\{\widehat v_1,\widehat v_2\}$, 所以直接截取不符合该递推结构.
\end{example}
这个障碍并不是说 Ritz 向量不能用于重启. 它说明需要\emph{重新组织分解}: 隐式重启借助小 Hessenberg 矩阵上的移位 QR 步骤, 同时变换大空间基和残差结构; 厚重启也保留多个近似方向, 但使用相应的扩展递推形式. 二者都有比上述朴素截取更多的代数安排.

从多项式角度看, 隐式移位起到过滤不需要谱分量的作用. 若一次过滤对应 $p(A)$, 移位放在不需要的 Ritz 值附近, 则这些谱分量被削弱, 目标分量相对保留. 对非正规矩阵, 此图像还要考虑特征向量病态与多项式矩阵范数, 不能只看 $|p(\lambda)|$. 本章完成标准 Arnoldi 结构及朴素截取为何失效的证明; 隐式重启的完整实现还涉及移位 QR 的结构保持、截断和有限精度处理, 这里不把概述冒充其全部证明. 实际计算宜采用经过验证的实现, 并独立验算 Ritz 残差.

\originalsection{习题与解答}
\begin{exercise}
设 $A=\diag(1,2,2,4)$, $b=(1,0,1,0)^T$. 求 $b$ 相对于 $A$ 的 grade 和消去多项式, 并说明为什么 Arnoldi 不需要四步.
\end{exercise}
\begin{solution}
只有特征值 $1,2$ 的分量被激发. $b,Ab$ 独立, 而 $(A-I)(A-2I)b=0$, 因此 grade 为 $2$, $m_b(z)=(z-1)(z-2)$. Arnoldi 第二步产生不变子空间后中止; 未被 $b$ 激发的特征值 $4$ 不会凭空出现.
\end{solution}
\begin{exercise}
设 $A$ 可逆. 证明只用普通 Arnoldi 从一个固定 $b$ 开始, 无法保证发现 $A$ 的每个特征值. 给出一个实对称矩阵的例子, 并说明多个起始向量为何有助于补救.
\end{exercise}
\begin{solution}
取 $A=\diag(1,2,3)$, $b=e_1$, 则所有 Krylov 子空间都只有 $\Span\{e_1\}$, 唯一 Ritz 值为 $1$. 若用多个起始向量组成块 $B$, 则块空间 $\Span\{B,AB,\ldots\}$ 可包含更多特征方向. 但它能发现哪些方向仍取决于起始块对相应不变子空间的投影; 多起点不是无条件保证.
\end{solution}
\begin{exercise}
证明在 Arnoldi 未中止时, 第 $k$ 步产生的方向满足
\[
\min_{\substack{p\text{ 首一}\\\deg p=k}}\norm{p(A)v_1}_2
=\prod_{j=1}^k h_{j+1,j}.
\]
由此解释为何这个多项式问题与移位无关.
\end{exercise}
\begin{solution}
反复使用 $Av_j\in\Span\{v_1,\ldots,v_{j+1}\}$, 可知 $A^kv_1$ 在 $v_{k+1}$ 上的系数为 $\prod_{j=1}^k h_{j+1,j}$, 其余部分落在 $\mathcal K_k$. 首一多项式 $p$ 的低次项 $p(A)v_1-A^kv_1$ 可以遍历 $\mathcal K_k$, 所以最小化就是从 $A^kv_1$ 中消去它在 $\mathcal K_k$ 上的正交投影. 剩余长度正为右侧. 对 $A+\mu I$, 首一次数 $k$ 的 $p(z)$ 与首一次数 $k$ 的 $q(z)=p(z+\mu)$ 一一对应, 故最小值相同. 此结论是关于新 Krylov 方向的多项式最小化, 不能直接等同于每个 Ritz 对的同一误差界.
\end{solution}
\begin{exercise}
设 $A=A^*$, $x$ 为单位向量, $\theta=x^*Ax$. 若 $\theta$ 附近有两个非常接近的特征值, 为什么小残差通常只能证明 $x$ 接近相应二维不变子空间, 而不能证明它接近其中某一个固定特征向量?
\end{exercise}
\begin{solution}
在特征基中, 残差长度为 $\sum_i|\lambda_i-\theta|^2|c_i|^2$ 的平方根. 对远离 $\theta$ 的特征值, 系数受到间隔下界控制; 对两个接近 $\theta$ 的特征值, 两个系数之间如何分配可以几乎不改变残差. 因而只能控制该簇以外的投影. 即使重复特征值对应的空间已经精确得到, 其中的单个正交基向量本来也不唯一.
\end{solution}

% SOURCE: MIT 18.335J Spring 2019, Week 7, Lecture 20; Week 8, Lecture 21.
% ADAPTATION: self-authored supplemental exposition and examples.
\originalchapter{GMRES 与多项式收敛分析}\label{ch:gmres}
\originalsection{从大方程到小最小二乘}
设 $A\in\C^{n\times n}$ 可逆, 希望解 $Ax=b$. 选初值 $x_0$, 记 $r_0=b-Ax_0$. 若 $r_0=0$, 已经完成求解; 以下假设 $\beta=\norm{r_0}_2>0$. 与其在整个 $\C^n$ 中任意修正 $x_0$, 我们在不断增大的仿射空间
\[
x_0+\mathcal K_k(A,r_0)
\]
中搜索. 这一选择每多一维只需要一次新的矩阵向量积, 而且包含用次数越来越高的多项式作用于残差得到的修正.

\begin{definition}
第 $k$ 步 GMRES 近似 $x_k$ 是上述仿射空间中使真实残差二范数最小的向量:
\begin{equation}\label{eq:gmres-definition}
 x_k=\operatorname*{arg\,min}_{x\in x_0+\mathcal K_k(A,r_0)}\norm{b-Ax}_2.
\end{equation}
名称来自 generalized minimal residual, 即广义最小残差法.
\end{definition}
可逆性使最小化解唯一: 若写成某个满列秩基 $Z_k$ 的坐标, 则 $AZ_k$ 也满列秩, 二次目标严格凸. 本章主要讨论可逆方阵; 对奇异或不相容系统, 最小二乘解、空间中止与最小范数解之间还需额外条件, 不能直接沿用这里所有结论.

对 $(A,r_0)$ 做 Arnoldi, 令 $v_1=r_0/\beta$. 第 $k$ 步尚未中止时有
\[
AV_k=V_{k+1}\overline H_k,
\qquad V_{k+1}^*V_{k+1}=I.
\]
任何候选写成 $x=x_0+V_ky$, 因而
\begin{equation}\label{eq:gmres-small-residual}
 b-Ax=V_{k+1}(\beta e_1-\overline H_ky),\qquad
 \norm{b-Ax}_2=\norm{\beta e_1-\overline H_ky}_2.
\end{equation}
于是原来的 $n$ 维问题等价于一个 $(k+1)\times k$ 小最小二乘问题:
\begin{equation}\label{eq:gmres-small-ls}
 y_k=\operatorname*{arg\,min}_{y\in\C^k}\norm{\beta e_1-\overline H_ky}_2,
 \qquad x_k=x_0+V_ky_k.
\end{equation}
这里不能把 $\overline H_k$ 随意换成方阵 $H_k$. 最后一行 $h_{k+1,k}e_k^*$ 正是候选解在当前子空间外产生的残差, 忽略它会改变最小化目标.

\begin{proposition}[最小残差的正交条件]\label{prop:gmres-optimality}
GMRES 残差 $r_k=b-Ax_k$ 满足
\[
r_k\perp A\mathcal K_k(A,r_0).
\]
等价地, $\overline H_k^*(\beta e_1-\overline H_ky_k)=0$. 此条件在可逆情形下唯一确定 $x_k$.
\end{proposition}
\begin{proof}
对任意修正方向 $v\in\mathcal K_k$, 目标平方沿 $x_k+tv$ 为 $\norm{r_k-tAv}_2^2$. 分别对实参数 $t$ 和纯虚参数 $t$ 求一阶变化为零, 得 $(Av)^*r_k=0$. 反过来, 若此正交关系成立, 则
\[
\norm{r_k-Av}_2^2=\norm{r_k}_2^2+\norm{Av}_2^2,
\]
所以 $x_k$ 是唯一极小点. 在 Arnoldi 坐标中代入式 \eqref{eq:gmres-small-residual}, 即得小问题的正规方程.
\end{proof}
该正规方程用于理解最优性, 不是建议计算时形成 $\overline H_k^*\overline H_k$: 后者会平方条件数. 稳定实现应使用 QR 分解.

\originalsection{QR 更新、存储与停止检验}
因为 $\overline H_k$ 为上 Hessenberg, 可用相邻两行的 Givens 酉旋转逐个消去次对角元. 若 $Q_k$ 是这些旋转组成的小酉矩阵, 则
\[
Q_k^*\overline H_k=\begin{bmatrix}R_k\\0\end{bmatrix},\qquad
Q_k^*(\beta e_1)=\begin{bmatrix}g_k\\\gamma_{k+1}\end{bmatrix},
\]
其中 $R_k\in\C^{k\times k}$ 上三角且可逆. 酉变换保持二范数, 所以
\[
\norm{\beta e_1-\overline H_ky}_2^2
=\norm{g_k-R_ky}_2^2+|\gamma_{k+1}|^2.
\]
解三角方程 $R_ky_k=g_k$ 后, 最小残差恰为 $|\gamma_{k+1}|$.

从 $k-1$ 步扩展到 $k$ 步时, 先把旧旋转作用到新 Hessenberg 列, 再增加一个旋转消去新次对角元. 小 QR 更新每步只需 $O(k)$ 运算, $k$ 步共需 $O(k^2)$; 不必每一步重新分解整个小矩阵. 大空间正交化仍需 $O(nk^2)$, 所以 GMRES 的主要增长成本通常来自保存和处理所有基向量. 可以先只跟踪 $|\gamma_{k+1}|$, 在达到候选容差时再解三角方程并形成 $x_k$.

一个基本实现按下列顺序组织:
\begin{enumerate}
\item 计算 $r_0=b-Ax_0$, 归一化为 $v_1$, 初始化小 QR 的右端为 $\beta e_1$.
\item 用一次 $Av_j$ 和正交化扩展 Arnoldi, 必要时重正交化并同步累加 $h_{ij}$.
\item 更新小 QR 和变换后的右端, 用新末分量估计残差长度.
\item 达到候选容差、到达重启长度或遇到近中止时, 解三角系统、形成近似解, 再检查真实残差.
\end{enumerate}
最后一步很重要. 浮点数中, 基不完全正交, Arnoldi 关系也有缺陷, 小问题给出的递推残差可能低估 $b-A\widehat x_k$. 重启时重新计算真实残差还可防止递推误差一直累积.

若目标是验证线性方程被稳定地求解, 可以用尺度无关的残差指标
\begin{equation}\label{eq:linear-backward-indicator}
\eta(\widehat x)=\frac{\norm{b-A\widehat x}_2}
{\norm{A}_2\norm{\widehat x}_2+\norm{b}_2}.
\end{equation}
它控制范数意义下的相对后向扰动. 例如令 $r=b-A\widehat x$, 把 $r$ 分成作用于矩阵和右端的两个扰动, 可构造 $\Delta A,\Delta b$ 使
$(A+\Delta A)\widehat x=b+\Delta b$, 且
$\norm{\Delta A}_2\le\eta\norm{A}_2$, $\norm{\Delta b}_2\le\eta\norm{b}_2$.
当 $\widehat x\ne0$ 时, 可取
\[
\Delta A=\frac{\eta\norm{A}_2}{\norm{\widehat x}_2}
\frac{r}{\norm{r}_2}\widehat x^*,\qquad
\Delta b=-\eta\norm{b}_2\frac{r}{\norm{r}_2},
\]
对 $r=0$ 则取零扰动. $\widehat x=0$ 的情形只需 $\Delta b=-b$. 但即使后向误差小, 前向解误差仍可被 $\kappa_2(A)$ 放大. 对 $b\ne0$, 常用的简明上界是
\[
\frac{\norm{x_* -\widehat x}_2}{\norm{x_*}_2}
\le\kappa_2(A)\frac{\norm{b-A\widehat x}_2}{\norm{b}_2},
\qquad x_*=A^{-1}b.
\]
停止容差应结合用户需要的解精度、条件数和数据本身的误差制定.

\originalsection{残差多项式与有限步中止}
\begin{theorem}[GMRES 的多项式表述]\label{thm:gmres-polynomial}
在精确算术中,
\begin{equation}\label{eq:gmres-poly-min}
\norm{r_k}_2=
\min_{\substack{p\in\C[z],\ \deg p\le k\\p(0)=1}}
\norm{p(A)r_0}_2.
\end{equation}
因此未重启 GMRES 的残差范数单调不增. 若 $r_0$ 相对于可逆 $A$ 的 grade 为 $d$, 则至迟第 $d\le n$ 步得到 $r_d=0$.
\end{theorem}
\begin{proof}
任意修正 $v\in\mathcal K_k$ 可写为 $v=q(A)r_0$, $\deg q\le k-1$. 其残差为
\[
r_0-Av=(I-Aq(A))r_0=p(A)r_0,\qquad p(z)=1-zq(z).
\]
该多项式次数至多 $k$, 常数项为 $1$. 反过来, 任何这样的 $p$ 都使 $(1-p(z))/z$ 为次数至多 $k-1$ 的多项式, 因而产生一个合法修正. 两个最小化问题完全等价.

次数至多 $k$ 的可行集包含前一步的可行集, 所以最小值不增. 由命题 \ref{prop:krylov-grade}, $m_{r_0}(0)\ne0$. 取 $p=m_{r_0}/m_{r_0}(0)$, 就有 $p(0)=1$ 且 $p(A)r_0=0$, 于是第 $d$ 步的最小残差为零.
\end{proof}
有限步中止是代数上的完整性结论, 不是大型计算的运行时间保证. 对数百万维问题, 等到 $n$ 步既存不下基也承担不起正交化; 浮点数还会破坏精确消去关系. 实际关键是能否在远少于 $n$ 步时达到足够精度.

\begin{example}
设 $A=\diag(1,2)$, $x_0=0$, $b=(1,1)^T$. 求一步 GMRES 并用多项式复核.
\end{example}
\begin{solution}
一步候选为 $x=\alpha b$. 残差为 $(1-\alpha,1-2\alpha)^T$, 目标平方为
$(1-\alpha)^2+(1-2\alpha)^2$. 最小点为 $\alpha=3/5$, 因而
\[
x_1=(3/5,3/5)^T,\qquad r_1=(2/5,-1/5)^T,\qquad
\norm{r_1}_2=1/\sqrt5.
\]
这等于对 $p(z)=1-\alpha z$ 最小化 $|p(1)|^2+|p(2)|^2$. 第二步可取
$p(z)=(1-z)(1-z/2)$, 对两个特征值都为零, 因而精确求解. 一步结果一般不等于沿 $b$ 最小化能量范数所得的 CG 结果, 因为两者的最优性准则不同.
\end{solution}

\originalsection{正规矩阵、谱簇与异常值}
\begin{theorem}[可对角化情形的谱界]\label{thm:gmres-spectral}
若 $A=X\Lambda X^{-1}$, 则
\begin{equation}\label{eq:gmres-spectral-bound}
\frac{\norm{r_k}_2}{\norm{r_0}_2}
\le\kappa_2(X)
\min_{\substack{\deg p\le k\\p(0)=1}}
\max_i|p(\lambda_i)|.
\end{equation}
若 $A$ 正规, 可以取 $X$ 酉, 因而前面的因子为 $1$. 此时对给定 $r_0=Xc$, 还有精确表述
\[
\norm{r_k}_2^2=
\min_{\substack{\deg p\le k\\p(0)=1}}
\sum_i|c_i|^2|p(\lambda_i)|^2.
\]
\end{theorem}
\begin{proof}
多项式满足 $p(A)=Xp(\Lambda)X^{-1}$. 对任意合法 $p$,
\[
\norm{p(A)r_0}_2\le\norm{X}_2\norm{X^{-1}}_2
\max_i|p(\lambda_i)|\norm{r_0}_2.
\]
最小化即得第一式. 若 $X$ 酉, 则长度保持, 直接按坐标求和得到最后的等式.
\end{proof}
该界解释为什么\emph{谱的位置}比单一条件数包含更多收敛信息. 例如谱集中在不含原点的圆盘 $|z-c|\le\rho<|c|$ 中, 取
$p(z)=(1-z/c)^k$, 得
\[
\norm{r_k}_2\le\kappa_2(X)(\rho/|c|)^k\norm{r_0}_2.
\]
谱簇越小且越远离原点, 越容易构造满足 $p(0)=1$ 而在簇上很小的低次多项式.

少量离群值通常可以用少量多项式根处理. 若 $\lambda_1,\ldots,\lambda_s\ne0$ 是离群值, 可取
\[
p(z)=\prod_{i=1}^s(1-z/\lambda_i)\,(1-z/c)^{k-s},\qquad k\ge s.
\]
该多项式在离群值处精确为零, 在主谱簇上受后一个因子控制. 前面各因子在簇上的最大值会贡献常数, 所以不能简单断言消去离群值完全没有代价. 但它揭示了 Krylov 方法常比只按全局条件数预测的速度更快的原因.

\begin{example}\label{ex:gmres-cycle}
令 $A$ 为 $n$ 阶循环置换矩阵, 满足 $Ae_j=e_{j+1}$ 对 $j<n$ 成立, $Ae_n=e_1$. 取 $x_0=0$, $b=e_1$. 证明 GMRES 在前 $n-1$ 步均不改善残差, 第 $n$ 步才精确求解.
\end{example}
\begin{solution}
$A$ 酉, 因而 $\kappa_2(A)=1$. 对 $k<n$,
\[
\mathcal K_k(A,b)=\Span\{e_1,\ldots,e_k\},\qquad
A\mathcal K_k=\Span\{e_2,\ldots,e_{k+1}\}\perp b.
\]
所以对每个候选 $x$, $\norm{b-Ax}_2^2=1+\norm{Ax}_2^2$, 唯一最小点是 $x_k=0$, 残差长度保持 $1$. 第 $n$ 步空间等于整个空间, 解为 $x_*=e_n$. 谱由单位圆上的 $n$ 次单位根组成, 并未在某个远离原点的小区域集中. 这说明小条件数本身不能保证少步 GMRES.
\end{solution}

\originalsection{非正规矩阵为什么需要额外信息}
若矩阵很非正规, 特征向量可能几乎平行, 定理 \ref{thm:gmres-spectral} 中 $\kappa_2(X)$ 很大. 若矩阵根本不可对角化, 该定理不能使用. 例如 Jordan 块上的 $p(A)$ 还含有 $p'(\lambda),p''(\lambda)$ 等导数项; 只知道 $p$ 在特征值处很小远远不够.

\begin{example}
取
\[
A=\begin{bmatrix}1&M\\0&1\end{bmatrix},\qquad b=e_2,\qquad x_0=0,
\]
其中 $M>0$. 它只有特征值 $1$. 一步 GMRES 的残差却怎样随 $M$ 变化?
\end{example}
\begin{solution}
候选 $x=\alpha e_2$, 残差为 $(-M\alpha,1-\alpha)^T$. 最小二乘给出
$\alpha=1/(M^2+1)$, 所以
\[
\frac{\norm{r_1}_2}{\norm{r_0}_2}=\frac{M}{\sqrt{M^2+1}}.
\]
当 $M$ 很大时, 一步几乎没有进展. 多项式 $p(z)=1-z$ 虽在唯一特征值处为零, 但
$p(A)=\begin{bmatrix}0&-M\\0&0\end{bmatrix}$ 并非零矩阵. 第二步可以用 $(1-z)^2$ 消去整个 Jordan 块, 于是精确终止.
\end{solution}
即使保持可对角化, 类似问题仍存在. 对 $A=\begin{bmatrix}1&M\\0&2\end{bmatrix}$ 和同一 $b$, 一步相对残差为 $M/\sqrt{M^2+4}$, 尽管两个特征值始终固定. 其特征向量矩阵随 $M$ 变得病态, 正好补充了谱位置遗漏的信息.

并非所有非正规问题都只能放弃一般估计. 一个可直接证明的充分条件是矩阵 Hermitian 部分严格正定.
\begin{theorem}[正实部给出的收敛保证]\label{thm:gmres-positive-real}
若 $\operatorname{Re}(v^*Av)\ge\nu\norm{v}_2^2$ 对所有 $v$ 成立, 其中 $\nu>0$, 且 $L=\norm{A}_2$, 则
\[
\frac{\norm{r_k}_2}{\norm{r_0}_2}
\le\left(1-\frac{\nu^2}{L^2}\right)^{k/2}.
\]
\end{theorem}
\begin{proof}
对实数 $\alpha>0$,
\[
\norm{(I-\alpha A)v}_2^2
=\norm{v}_2^2-2\alpha\operatorname{Re}(v^*Av)+\alpha^2\norm{Av}_2^2
\le(1-2\alpha\nu+\alpha^2L^2)\norm{v}_2^2.
\]
取 $\alpha=\nu/L^2$, 得 $\norm{I-\alpha A}_2\le\sqrt{1-\nu^2/L^2}$. 将合法多项式 $p(z)=(1-\alpha z)^k$ 代入式 \eqref{eq:gmres-poly-min}, 并利用矩阵范数的次乘性, 即得所需界. 此条件也保证 $A$ 可逆, 因为 $Av=0$ 会迫使 $v=0$.
\end{proof}
这里的信息来自整个二次型, 不只是特征值. 更精细的非正规收敛理论会使用数值域、伪谱或矩阵多项式范数. 对实际问题, 观测残差历史、正交缺陷和预条件后的算子行为, 通常比盯着少数估计特征值更稳妥.

\originalsection{重启和预条件改变了什么}
GMRES($m$) 表示每个周期只运行最多 $m$ 步, 然后用当前解作新初值, 重新计算残差并生成新 Krylov 基. 它把基存储限制为 $O(nm)$, 每周期正交化限制为 $O(nm^2)$. 但新周期只在当前残差生成的空间中搜索, 不再对过去全部方向联合最小化.

若每个周期的残差多项式为 $p_j$, 则 $s$ 个周期后形式上有
\[
r^{(s)}=p_s(A)\cdots p_1(A)r_0,\qquad
\deg p_j\le m,\quad p_j(0)=1.
\]
总次数至多 $sm$, 但这些因子逐周期贪心选取, 不等于一次性求全局最优的次数 $sm$ 多项式. 这解释了全 GMRES 的有限中止结论为何不能直接移植到重启方法.

例 \ref{ex:gmres-cycle} 给出严格停滞的情形. 若 $m<n$, 每一个周期从 $x=0$, $r=e_1$ 开始, 周期最优解仍为 $0$. 因而 GMRES($m$) 可无限重复且残差完全不变. 在精确算术中, 可逆矩阵也不能保证任意固定重启长度的 GMRES 收敛.

预条件的目的就是在保持同一原问题解的同时, 改善迭代看到的算子. 本节假设 $M$ 可逆. 本书约定 $M$ 是容易求解的 $A$ 的近似, 实现只需要求 $Mz=v$, 不显式形成 $M^{-1}$. 左预条件系统是
\[
M^{-1}Ax=M^{-1}b.
\]
对它使用 GMRES, 最小化的是 $\norm{M^{-1}(b-Ax)}_2$. 若 $M^{-1}$ 对某些方向缩放很强, 该残差小可能不代表原始残差同样小, 所以应同时检查 $b-Ax$.

右预条件可写成修正形式
\[
AM^{-1}y=r_0,\qquad x=x_0+M^{-1}y.
\]
对该系统使用 GMRES, 最小化的就是原方程残差 $\norm{r_0-AM^{-1}y}_2$. 左右预条件算子相似:
$M^{-1}A=M^{-1}(AM^{-1})M$, 因而谱相同, 但欧氏内积下的非正规程度、起始残差和最小化范数可以不同, 实际残差历史不必相同.

若每次应用的预条件器变化, 不能再把整个过程当作一个固定矩阵 $AM^{-1}$ 的 Krylov 方法. 灵活 GMRES 在第 $j$ 步保存实际求得的 $z_j$, 用 $Az_j$ 正交化, 并令最终修正为 $\sum_j y_jz_j$. 它仍有小最小二乘结构, 但固定算子的残差多项式理论一般不再原样适用. 下一章会在 Hermitian 正定条件下推导更节省存储的 CG 以及对称预条件方式.

\originalsection{习题与解答}
\begin{exercise}
\label{nla:fom-exercise}
本题对应 MIT 18.335J Spring 2019 PS4 第 1 题 (见第\ref{ps4:q1}题), 并增加奇异压缩矩阵的比较. 在 GMRES 的同一个 Krylov 空间中, 另一种方法 FOM 要求 $r_k\perp\mathcal K_k$ 而非 $r_k\perp A\mathcal K_k$. 推导它的小方程, 并说明这个方程可能失败而 GMRES 最小二乘仍有解.
\end{exercise}
\begin{solution}
写 $x=x_0+V_ky$, 则 $V_k^*r=\beta e_1-H_ky$, 故 FOM 解 $H_ky=\beta e_1$. 即使 $A$ 可逆, 压缩矩阵 $H_k$ 也可能奇异. 例如 $A=\begin{bmatrix}0&1\\1&0\end{bmatrix}$, $r_0=e_1$, 一步有 $H_1=[0]$, FOM 方程无解; GMRES 的 $\overline H_1=(0,1)^T$ 满列秩, 最小点仍存在, 此时修正为零.
\end{solution}
\begin{exercise}
若 $A$ 正规, 谱包含在 $|z-3|\le1$ 内, 给出 $k$ 步 GMRES 的相对残差上界. 若只知道 $A$ 可对角化且同样的谱包含关系, 还缺少什么信息?
\end{exercise}
\begin{solution}
取 $p(z)=(1-z/3)^k$, 正规情形得 $\norm{r_k}_2/\norm{r_0}_2\le3^{-k}$. 可对角化情形还需控制特征向量矩阵的 $\kappa_2(X)$; 谱包含关系本身不能约束这个因子.
\end{solution}
\begin{exercise}
证明若 $\norm{I-\alpha A}_2\le q<1$ 对某个 $\alpha\in\C$ 成立, 则任意重启长度 $m\ge1$ 的 GMRES($m$) 每周期至少满足相对残差缩减因子 $q^m$.
\end{exercise}
\begin{solution}
对每个周期的当前残差 $r$, 合法的 $m$ 次多项式 $(1-\alpha z)^m$ 给出候选残差, 长度不超过 $q^m\norm{r}_2$. GMRES 周期最小值不大于此候选, 所以每周期都有同一个收缩保证. 若提前幸运中止, 残差已为零. 该充分条件解释了为什么定理 \ref{thm:gmres-positive-real} 的假设也能保证固定长度重启收敛, 而一般可逆性不能.
\end{solution}
\begin{exercise}
设左预条件 GMRES 输出 $\widehat x$ 且 $\norm{M^{-1}(b-A\widehat x)}_2\le\tau$. 证明原残差不超过 $\norm{M}_2\tau$, 并说明为什么实际仍宜直接计算原残差.
\end{exercise}
\begin{solution}
令 $\widetilde r=M^{-1}(b-A\widehat x)$, 则 $b-A\widehat x=M\widetilde r$, 长度至多 $\norm{M}_2\tau$. 若 $\norm{M}_2$ 很大, 这个界很弱; 若 $\tau$ 只是浮点递推估计, 还可能存在残差间隙. 直接计算 $b-A\widehat x$ 同时解决尺度解释和递推真实性两个问题.
\end{solution}

% SOURCE: MIT 18.335J Spring 2019, Week 8, Lectures 21--23.
% ADAPTATION: self-authored supplemental explanations, proofs, examples and exercises.
\originalchapter{共轭梯度、预条件与双共轭方法}\label{ch:cg}
\originalsection{正定性把解方程变成极小化}
现在假设 $A\in\C^{n\times n}$ Hermitian 正定, 即 $A=A^*$ 且 $x^*Ax>0$ 对所有 $x\ne0$ 成立. 定义能量内积与能量范数
\[
\inner{x}{y}_A=x^*Ay,\qquad \norm{x}_A=(x^*Ax)^{1/2},
\]
并令 $x_*=A^{-1}b$. 我们考虑实值二次函数
\[
f(x)=\frac12x^*Ax-\operatorname{Re}(b^*x).
\]
其变化由残差 $r=b-Ax$ 控制. 把 $x=x_*+h$ 展开, 交叉项因 $Ax_*=b$ 消去, 得到
\begin{equation}\label{eq:cg-quadratic-energy}
f(x)-f(x_*)=\frac12h^*Ah=\frac12\norm{x-x_*}_A^2.
\end{equation}
因此解方程、极小化 $f$, 以及极小化能量误差是同一个任务. 正定性同时保证极小点存在、唯一, 并给出一个真正的误差范数. 若 $A$ 只有 Hermitian 性而非正定性, $f$ 可能为鞍点函数, 后面分母的正性和极小化证明都将失效.

先看最简单的最速下降. 在当前残差方向 $r$ 上选择实步长 $\alpha$, 有
\[
f(x+\alpha r)=f(x)-\alpha r^*r+\frac{\alpha^2}{2}r^*Ar.
\]
唯一最优步长为 $\alpha=r^*r/(r^*Ar)$. 更新后 $r_{\mathrm{new}}=r-\alpha Ar$, 从而 $r^*r_{\mathrm{new}}=0$. 连续残差彼此正交看似有利, 但并不意味着同时保留了对过去各方向的最优性: 后续更新会重新改变以前已经调好的坐标.

\begin{example}
设 $A=\diag(1,100)$, $b=0$, $x_0=(1,1/100)^T$. 解释最速下降的折返和慢收敛.
\end{example}
\begin{solution}
$r_0=(-1,-1)^T$, $\alpha_0=2/101$, 故
\[
x_1=(99/101,-99/10100)^T,\qquad
r_1=(-99/101,99/101)^T.
\]
两个残差方向正交. 由于 $r_1$ 的两个分量绝对值仍相等, 下一步仍用 $\alpha_1=2/101$, 得
$x_2=(99/101)^2x_0$. 之后这一模式重复, 能量误差每步乘以 $99/101$. 二次函数的等高线很扁, 欧氏最陡方向在长轴两侧折返, 因而仅沿当前梯度精确搜索仍很慢.
\end{solution}
若 $a=\lambda_{\min}(A)>0$, $c=\lambda_{\max}(A)$, 则选择固定步长 $2/(a+c)$ 已保证能量误差每步不超过 $(c-a)/(c+a)$. 最速下降的精确线搜索不会更差, 因而其一般保证需要约 $O(\kappa_2(A)\log(1/\tau))$ 步达到相对误差 $\tau$. 共轭梯度的改进点不是再找一个更精细的单步长度, 而是使新的搜索不破坏所有旧方向的极小化.

\originalsection{共轭方向与整个子空间上的最优性}
\begin{definition}
非零向量 $p_0,\ldots,p_{k-1}$ 若满足 $p_i^*Ap_j=0$ 对 $i\ne j$ 成立, 称为 $A$-共轭或 $A$-正交. 因 $A$ 正定, 它们线性无关.
\end{definition}
这里的共轭并不是复共轭运算的同义词, 而是通过 $A$ 定义的正交性. 如果用这些方向作坐标, 二次函数中的混合项全部消失, 各坐标的一维极小化就不会互相干扰.

\begin{proposition}[共轭方向的联合极小化]\label{prop:conjugate-directions}
设 $p_0,\ldots,p_{k-1}$ 两两 $A$-共轭. 从 $x_0$ 开始依次令
\[
\alpha_j=\frac{p_j^*r_j}{p_j^*Ap_j},\qquad
x_{j+1}=x_j+\alpha_jp_j,\qquad
r_{j+1}=r_j-\alpha_jAp_j.
\]
则 $r_k\perp\Span\{p_0,\ldots,p_{k-1}\}$, 且 $x_k$ 是该仿射子空间上 $f$ 的唯一极小点, 也使能量误差最小. 在复数情形, 上述系数允许复数; 这是沿复线的极小化公式.
\end{proposition}
\begin{proof}
更新消去 $p_j^*r_j$ 的分量, 所以 $p_j^*r_{j+1}=0$. 对 $i<j$, 已有 $p_i^*r_j=0$, 新变化为 $-\alpha_jp_i^*Ap_j=0$, 所以旧正交性保留. 最终残差垂直于全部方向.

对任意 $h$ 在这些方向的线性包中, 展开二次函数得
\[
f(x_k+h)-f(x_k)=-\operatorname{Re}(h^*r_k)+\frac12h^*Ah
=\frac12h^*Ah\ge0.
\]
等号仅在 $h=0$ 时成立. 再用式 \eqref{eq:cg-quadratic-energy} 即得能量误差最小化.
\end{proof}
任意给定的一组方向都可以用能量内积做 Gram--Schmidt, 但保存所有旧方向并作全部投影依然昂贵. CG 的特殊之处是从残差生成方向时, 绝大多数能量投影自动为零, 最终只需前一个方向.

\originalsection{从 Krylov 极小化推导三项算法}
令 $x_k$ 是 $x_0+\mathcal K_k(A,r_0)$ 上 $f$ 的唯一极小点, $r_k=b-Ax_k$. 这里先把 $x_k$ 用最优性定义, 再导出怎样低成本地得到它, 以免把待证明的递推性质当成假设.

\begin{lemma}\label{lem:cg-residual-space}
在精确算术中,
\[
r_k\perp\mathcal K_k(A,r_0),\qquad r_k\in\mathcal K_{k+1}(A,r_0).
\]
若 $r_k\ne0$, 则 $\mathcal K_{k+1}=\mathcal K_k\oplus\Span\{r_k\}$. 不同非零残差彼此正交.
\end{lemma}
\begin{proof}
对任意 $h\in\mathcal K_k$, 沿 $x_k+th$ 和 $x_k+\mathrm it h$ 的实参数一阶条件给出 $h^*r_k=0$. 又 $x_k-x_0\in\mathcal K_k$, 所以
$r_k=r_0-A(x_k-x_0)\in\mathcal K_{k+1}$. 若 $r_k\ne0$, 它正交于 $\mathcal K_k$ 且不属于其中, 故维数至少增加一. Krylov 空间每一步最多增加一维, 所以所述直和成立. 对 $i<k$, $r_i\in\mathcal K_{i+1}\subseteq\mathcal K_k$, 故 $r_i^*r_k=0$.
\end{proof}

现在令 $p_0=r_0$. 假设 $p_0,\ldots,p_{k-1}$ 已是 $\mathcal K_k$ 的 $A$-共轭基. 对 $r_k$ 做能量正交化, 得
\[
p_k=r_k-\sum_{i=0}^{k-1}
\frac{p_i^*Ar_k}{p_i^*Ap_i}\,p_i.
\]
对 $i\le k-2$, $p_i\in\mathcal K_{i+1}$, 因而 $Ap_i\in\mathcal K_{i+2}\subseteq\mathcal K_k$. Hermitian 性及引理给出
\[
p_i^*Ar_k=(Ap_i)^*r_k=0.
\]
只剩最后一个投影. 因此 $p_k=r_k+\beta_{k-1}p_{k-1}$, 其中
\[
\beta_{k-1}=-\frac{p_{k-1}^*Ar_k}{p_{k-1}^*Ap_{k-1}}.
\]
新方向不为零, 因为其欧氏正交于旧空间的分量为非零的 $r_k$. 它补齐 $\mathcal K_{k+1}$, 并与所有旧方向 $A$-共轭.

命题 \ref{prop:conjugate-directions} 表明沿 $p_k$ 更新就得到新的子空间极小点. 因 $r_k$ 垂直旧方向,
$p_k^*r_k=r_k^*r_k$, 故
\[
\alpha_k=\frac{r_k^*r_k}{p_k^*Ap_k},\qquad
r_{k+1}=r_k-\alpha_kAp_k.
\]
为了把 $\beta$ 改写成只含残差长度的公式, 用前一步关系
$Ap_{k-1}=(r_{k-1}-r_k)/\alpha_{k-1}$, 得
\[
p_{k-1}^*Ar_k
=\frac{(r_{k-1}-r_k)^*r_k}{\alpha_{k-1}}
=-\frac{r_k^*r_k}{\alpha_{k-1}}.
\]
再代入 $\alpha_{k-1}=r_{k-1}^*r_{k-1}/(p_{k-1}^*Ap_{k-1})$, 最终得到
$\beta_{k-1}=r_k^*r_k/(r_{k-1}^*r_{k-1})$.

\begin{theorem}[CG 算法及其精确性质]\label{thm:cg-algorithm}
设 $A=A^*>0$, $r_0=b-Ax_0\ne0$, $p_0=r_0$. 在残差非零时反复计算
\begin{equation}\label{eq:cg-algorithm}
\begin{aligned}
q_k&=Ap_k,&\alpha_k&=\frac{r_k^*r_k}{p_k^*q_k},\\
x_{k+1}&=x_k+\alpha_kp_k,&r_{k+1}&=r_k-\alpha_kq_k,\\
\beta_k&=\frac{r_{k+1}^*r_{k+1}}{r_k^*r_k},&
p_{k+1}&=r_{k+1}+\beta_kp_k.
\end{aligned}
\end{equation}
若 $r_{k+1}=0$, 在计算下一方向前停止. 每个非零搜索方向的分母 $p_k^*Ap_k$ 为正. 由算法得到的 $x_k$ 恰为 $x_0+\mathcal K_k(A,r_0)$ 中能量误差最小的向量; 残差彼此欧氏正交, 搜索方向彼此 $A$-共轭. 若 grade 为 $d$, 则至迟第 $d\le n$ 步精确终止.
\end{theorem}
\begin{proof}
前面的构造从唯一的子空间极小点逐步推出了式 \eqref{eq:cg-algorithm}, 故它同时证明递推与该极小点一致. 正分母来自正定性与方向非零. 当维数不再增长时, 引理 \ref{lem:cg-residual-space} 表明残差既属于 $\mathcal K_d$ 又垂直它, 因而只能为零. 也可由 $A^{-1}r_0\in\mathcal K_d$ 得知精确解已经属于候选空间.
\end{proof}
每步只需一次矩阵向量积、若干内积和向量更新, 不必保存全部 Krylov 基. 总向量存储为 $O(n)$, 是 CG 相比全 GMRES 的重要优势. 这并不是把 GMRES 的基随意删除, 而是正定极小化和 Hermitian 三项结构带来的合法压缩.

\begin{example}
对 $A=\diag(1,2)$, $b=(1,1)^T$, $x_0=0$, 计算两步 CG 并与上一章的一步 GMRES 比较.
\end{example}
\begin{solution}
$r_0=p_0=(1,1)^T$, $\alpha_0=2/3$, 因而
\[
x_1=(2/3,2/3)^T,\qquad r_1=(1/3,-1/3)^T,
\quad\beta_0=1/9,\quad p_1=(4/9,-2/9)^T.
\]
$Ap_1=(4/9,-4/9)^T$, $p_1^*Ap_1=8/27$, $r_1^*r_1=2/9$, 故 $\alpha_1=3/4$. 更新得 $x_2=(1,1/2)^T$, $r_2=0$. 可直接核对 $p_0^*Ap_1=0$.

一步 GMRES 为 $(3/5,3/5)^T$, 欧氏残差长度 $1/\sqrt5$; 一步 CG 的残差长度为 $\sqrt2/3$, 略大. 但 CG 的能量误差为 $1/\sqrt6$, 小于一步 GMRES 的 $\sqrt{9/50}$. 两者都在同一个一维空间中取最优点, 区别来自所最小化的范数.
\end{solution}

\originalsection{Chebyshev 界与谱簇带来的进一步加速}
CG 的多项式观点与 GMRES 相似, 只是优化对象换成能量误差. 记 $e_k=x_*-x_k$. 因 $r_0=Ae_0$, 有
\[
e_k=e_0-q(A)r_0=(I-Aq(A))e_0=p(A)e_0,\qquad p(0)=1,\quad\deg p\le k.
\]
子空间极小化立即给出
\begin{equation}\label{eq:cg-poly}
\norm{e_k}_A=
\min_{\substack{\deg p\le k\\p(0)=1}}\norm{p(A)e_0}_A
\le\min_{\substack{\deg p\le k\\p(0)=1}}
\max_{\lambda\in\operatorname{spec}(A)}|p(\lambda)|\,\norm{e_0}_A.
\end{equation}
最后一步把 $A$ 在酉特征基中对角化, 使两边平方成为带权和 $\sum_i\lambda_i|p(\lambda_i)|^2|c_i|^2$ 的比较. 正定性使所有权重非负.

为从特征值所在区间构造多项式, 定义 Chebyshev 多项式 $T_0(t)=1$, $T_1(t)=t$, $T_{k+1}(t)=2tT_k(t)-T_{k-1}(t)$. 余弦加法公式归纳给出
$T_k(\cos\theta)=\cos(k\theta)$, 所以 $|T_k(t)|\le1$ 在 $[-1,1]$ 上成立. 对 $t>1$, 令 $\rho=t+\sqrt{t^2-1}>1$, 递推同样给出
\[
T_k(t)=\frac12(\rho^k+\rho^{-k})\ge\frac12\rho^k.
\]
不需要先掌握最佳一致逼近定理, 这两条性质就足以构造我们需要的收敛界.

\begin{theorem}[CG 的条件数界]\label{thm:cg-chebyshev}
设 $0<a=\lambda_{\min}(A)<c=\lambda_{\max}(A)$, $\kappa=c/a$. 则精确 CG 满足
\begin{equation}\label{eq:cg-chebyshev}
\frac{\norm{e_k}_A}{\norm{e_0}_A}
\le2\left(\frac{\sqrt\kappa-1}{\sqrt\kappa+1}\right)^k.
\end{equation}
若 $a=c$, 则 $A=aI$, 一步即得到精确解.
\end{theorem}
\begin{proof}
线性函数 $t(z)=(c+a-2z)/(c-a)$ 把 $[a,c]$ 映到 $[-1,1]$, 且 $t(0)=(c+a)/(c-a)>1$. 取合法多项式
\[
p_k(z)=\frac{T_k(t(z))}{T_k(t(0))}.
\]
分子在谱区间上的绝对值至多 $1$, 分母至少为 $\rho^k/2$, 其中
\[
\rho=t(0)+\sqrt{t(0)^2-1}
=\frac{\sqrt c+\sqrt a}{\sqrt c-\sqrt a}.
\]
于是 $\max_{z\in[a,c]}|p_k(z)|\le2\rho^{-k}$. 代入式 \eqref{eq:cg-poly} 得结论.
\end{proof}
由 $\log((\sqrt\kappa+1)/(\sqrt\kappa-1))\ge2/\sqrt\kappa$, 达到能量相对误差 $\tau$ 的充分步数约为
$\tfrac12\sqrt\kappa\log(2/\tau)$. 与最速下降的 $\kappa$ 量级相比, 条件数依赖改善为平方根. 但该界只使用谱的两个端点, 所以常常偏保守.

例如有 $s$ 个小特征值远离主簇 $[a,c]$, 可以在多项式中先放入 $s$ 个对应根, 再对主簇使用 Chebyshev 因子. 那么端点比主要由较窄的主簇控制, 乘上有限个离群值因子在主簇上的界. 迭代中逐渐抑制离群方向后, 观测收敛常会加快, 有时称为超线性阶段. 这是一种依赖谱分布和起始误差分量的现象, 不是所有 SPD 矩阵都必然出现的统一速度定理.

\originalsection{对称预条件与 PCG 的正确内积}
把左预条件系统写成 $M^{-1}Ax=M^{-1}b$ 时, 即使 $A,M$ 都 Hermitian 正定, $M^{-1}A$ 在欧氏内积中通常也不 Hermitian. 直接把它当作普通 SPD 矩阵运行标准 CG, 没有正确性依据. 正确推导要通过一个对称变换, 或等价地改变内积.

设 $M=CC^*$ 是一个可逆 Cholesky 因子分解. 令 $y=C^*x$, 则原方程等价于
\begin{equation}\label{eq:pcg-symmetric}
\widetilde A y=\widetilde b,\qquad
\widetilde A=C^{-1}AC^{-*},\qquad\widetilde b=C^{-1}b.
\end{equation}
$\widetilde A$ 显然 Hermitian, 且 $v^*\widetilde Av=(C^{-*}v)^*A(C^{-*}v)>0$, 所以它适合普通 CG. 算法中不必真的存储 $C$; 它只是证明对称性的工具.

将变换后的残差记为 $\widetilde r=C^{-1}r$. 定义 $z=M^{-1}r$, 则
$\widetilde r^*\widetilde r=r^*M^{-1}r=r^*z$. 把变换后搜索方向映回原坐标, 就得到以下递推.

\begin{theorem}[预条件共轭梯度]\label{thm:pcg}
若 $A=A^*>0$, 固定预条件矩阵 $M=M^*>0$, 令 $r_0=b-Ax_0$; 若 $r_0=0$ 则直接停止. 否则解 $Mz_0=r_0$, 取 $p_0=z_0$, $\rho_0=r_0^*z_0$. 对非零残差计算
\begin{equation}\label{eq:pcg-algorithm}
\begin{aligned}
q_k&=Ap_k,&\alpha_k&=\frac{\rho_k}{p_k^*q_k},\\
x_{k+1}&=x_k+\alpha_kp_k,&r_{k+1}&=r_k-\alpha_kq_k,\\
Mz_{k+1}&=r_{k+1},&\rho_{k+1}&=r_{k+1}^*z_{k+1},\\
\beta_k&=\frac{\rho_{k+1}}{\rho_k},&p_{k+1}&=z_{k+1}+\beta_kp_k.
\end{aligned}
\end{equation}
若 $r_{k+1}=0$ 则停止. 这是对 \eqref{eq:pcg-symmetric} 运行普通 CG 后的等价算法. 它仍使原坐标中的能量误差最小, 但候选空间为
\[
x_0+\mathcal K_k(M^{-1}A,M^{-1}r_0).
\]
误差界 \eqref{eq:cg-chebyshev} 中的条件数换成 $\kappa_2(C^{-1}AC^{-*})$.
\end{theorem}
\begin{proof}
变换后初始残差为 $C^{-1}r_0$, 初始搜索方向相同. 映回原坐标时, 搜索方向成为 $C^{-*}C^{-1}r_0=M^{-1}r_0$. 每次 $\widetilde p^*\widetilde A\widetilde p$ 等于 $p^*Ap$, 每次 $\widetilde r^*\widetilde r$ 等于 $r^*M^{-1}r$, 所以 CG 的 $\alpha,\beta$ 完全变成上述公式. 更新同样一一对应.

映回的子空间可写为
$C^{-*}\mathcal K_k(\widetilde A,C^{-1}r_0)=\mathcal K_k(M^{-1}A,M^{-1}r_0)$. 能量满足
$\norm{C^*e}_{\widetilde A}^2=e^*Ae$, 因而变换后能量最小化和误差界恰为原坐标中的所述结论.
\end{proof}
另一种看法是 $M^{-1}A$ 在内积 $\inner{u}{v}_M=u^*Mv$ 中自伴且正定. 其特征值是广义特征值问题 $Av=\lambda Mv$ 的正实数, 与 $\widetilde A$ 的谱相同. 预条件成功通常要求这些特征值集中, 同时应用 $M^{-1}$ 足够便宜.

\begin{proposition}[谱等价的含义]\label{prop:pcg-spectral-equivalence}
若存在 $0<c_1\le c_2$ 使
\[
c_1v^*Mv\le v^*Av\le c_2v^*Mv\qquad\text{对所有 }v,
\]
则 $\operatorname{spec}(\widetilde A)\subseteq[c_1,c_2]$, 从而预条件条件数不超过 $c_2/c_1$.
\end{proposition}
\begin{proof}
令 $v=C^{-*}w$, 则 $v^*Mv=w^*w$, $v^*Av=w^*\widetilde Aw$. 不等式成为 $c_1\norm{w}_2^2\le w^*\widetilde Aw\le c_2\norm{w}_2^2$. 取单位特征向量就得到谱界.
\end{proof}
这个准则能把预条件器的设计转化为可证明的二次型比较. 只要求 $M$ 的每个元素看起来接近 $A$ 通常不够; 要关心难收敛方向在两个二次型中受到怎样的缩放.

\originalsection{常用预条件器、有限精度与验收}
Jacobi 预条件取 $M=\diag(A)$, 对 SPD 矩阵其对角元均为正. 它便宜, 能修正变量尺度差异, 但只看局部独立坐标, 不一定处理好网格方程的低频误差. 块 Jacobi 把强耦合的变量成组, 在每个小块中求解, 因而能捕捉比单点更强的局部结构.

不完全 Cholesky 在消元时舍弃部分填充, 得到易解的近似因子 $M=LL^*$. 但舍弃可能使未修正的分解遇到非正主元; 必须用保证正定的构造、移位或其他受控修改, 不能因原矩阵 SPD 就断言任意丢弃策略也 SPD. 对一般非对称矩阵, 不完全 LU 常与 GMRES 或稳定化双共轭方法配合; 非对称 ILU 不可无条件用于 PCG.

多重网格通过细网格平滑处理高频误差, 再在粗网格校正低频误差. 若作为 PCG 预条件器使用, 整个循环的线性作用还应满足所需的对称性和正定性, 例如使用配对的前后平滑与一致的转移算子. 多重网格并非只把一个 Jacobi 步重复很多次, 它的核心是跨尺度校正.

预条件器的总成本应比较
\[
\text{构造成本}+k\bigl(\text{一次 }Av\text{ 的成本}+\text{一次 }M^{-1}v\text{ 的成本}\bigr).
\]
更少的迭代不一定更快. 若有许多右端, 较昂贵的预条件构造可被摊销; 若矩阵频繁变化, 则构造成本更敏感. 在理想但昂贵的极限 $M=A$ 下, PCG 一步收敛, 可是应用预条件器本身已经等于解原问题.

浮点 CG 中, 残差正交性、方向共轭性和 $n$ 步精确终止都不再严格成立. 递推残差还会与真实残差分离. 每隔适当阶段或临近收敛时重算 $b-A\widehat x$ 有助于检查; 若用真实残差替换递推残差, 后续方向也可能需要重新启动, 不能任意替换后仍声称精确共轭性质保持不变.

CG 保证单调下降的是能量误差和二次目标, 而非欧氏残差范数. 正确的停止检验宜同时考虑真实残差、相对尺度和应用误差. 对 SPD 矩阵, $\norm{e}_A^2=r^*A^{-1}r$, 并且
\[
\frac{\norm{r}_2}{\sqrt{\lambda_{\max}(A)}}
\le\norm{e}_A\le
\frac{\norm{r}_2}{\sqrt{\lambda_{\min}(A)}}.
\]
这些界需要特征值端点或可靠估计; 不能把一个小残差百分比自动当作同样大小的能量误差百分比.

若预条件器随迭代改变, 标准 PCG 的固定变换证明失效. 只有当变动方式符合某个灵活 CG 方法的假设并采用相应递推, 才能继续给出理论保证. 对变动、非线性或不精确内层求解, 灵活 GMRES 往往是更直接的替代, 仍需检查实际算子和收敛.

\originalsection{BiCG 的动机与不可忽略的中断}
非 Hermitian 矩阵通常不能既保留 GMRES 的欧氏残差最小化, 又维持简单的三项递推. 双共轭梯度 BiCG 采用两个互为对偶的 Krylov 序列, 用双正交性换取短递推. 它的内存较低, 但失去了 SPD 能量极小化的保护.

以下给出实矩阵版本, 以避免复共轭系数遮住基本结构. 取初始残差 $r_0$ 和影子残差 $\widetilde r_0$, 要求 $\rho_0=\widetilde r_0^Tr_0\ne0$, 并置 $p_0=r_0$, $\widetilde p_0=\widetilde r_0$. 形式上的递推为
\begin{equation}\label{eq:bicg}
\begin{aligned}
q_k&=Ap_k,&\widetilde q_k&=A^T\widetilde p_k,
&\alpha_k&=\frac{\rho_k}{\widetilde p_k^Tq_k},\\
x_{k+1}&=x_k+\alpha_kp_k,&
r_{k+1}&=r_k-\alpha_kq_k,&
\widetilde r_{k+1}&=\widetilde r_k-\alpha_k\widetilde q_k,\\
\rho_{k+1}&=\widetilde r_{k+1}^Tr_{k+1},&
\beta_k&=\frac{\rho_{k+1}}{\rho_k},\\
p_{k+1}&=r_{k+1}+\beta_kp_k,&
\widetilde p_{k+1}&=\widetilde r_{k+1}+\beta_k\widetilde p_k.
\end{aligned}
\end{equation}
在没有中断的精确双 Lanczos 构造中, 残差满足双正交关系 $\widetilde r_i^Tr_j=0$ 对 $i\ne j$, 方向满足 $\widetilde p_i^TAp_j=0$ 对 $i\ne j$. 这里没有正定性保证 $\rho_k$ 或 $\widetilde p_k^TAp_k$ 非零. 任一关键分母为零时, 上述公式就不能继续, 即使原矩阵可逆也是如此.

\begin{example}
取 $A=\begin{bmatrix}0&1\\1&0\end{bmatrix}$, $r_0=\widetilde r_0=e_1$. 则 $\rho_0=1$, 但 $\widetilde p_0^TAp_0=e_1^TAe_1=0$. BiCG 第一个步长即无定义, 而 $A$ 可逆, 方程当然有唯一解.
\end{example}
从块矩阵的角度,
\[
B=\begin{bmatrix}0&A\\A^*&0\end{bmatrix}
\]
是 Hermitian, 但一般不正定. 若 $\sigma_i(A)>0$, 则其特征值含 $\pm\sigma_i(A)$: 对一个奇异向量对 $Av=\sigma u$, $A^*u=\sigma v$, 向量 $(u,\pm v)^T$ 分别对应 $\pm\sigma$. 因而借助此块矩阵得到的共轭构造面向鞍点, 不能套用正定极小化证明. 这解释了课程中将 BiCG 与 Hermitian 不定块系统联系起来时, 为什么中断成为中心问题.

QMR 通过准最小残差处理双 Lanczos 序列; BiCGSTAB 在双共轭多项式之外增加局部残差平滑因子, 常可减少剧烈振荡并避免显式乘以 $A^T$. 这些方法是有价值的低存储选择, 但名称中的稳定化不表示绝无中断或对所有问题保证单调收敛. 若影子配对近于零、局部最小化分母退化或有限精度损坏双正交性, 仍可能失败. 应报告中断原因, 检查真实残差, 并考虑改变预条件器、使用经过验证的稳定化实现或改用 GMRES. 对实对称不定系统, MINRES 利用 Lanczos 和最小残差条件, 通常比硬套 CG 更合适.

\originalsection{习题与解答}
\begin{exercise}
证明非零的两两 $A$-共轭向量线性无关. 指出如果 $A$ 只是 Hermitian 不定, 证明的哪一步可能失效.
\end{exercise}
\begin{solution}
若 $\sum_jc_jp_j=0$, 左乘 $p_i^*A$ 得 $c_ip_i^*Ap_i=0$. 正定性保证 $p_i^*Ap_i>0$, 故每个 $c_i=0$. 不定情形可能有非零 $p_i$ 满足 $p_i^*Ap_i=0$, 于是无法除去该因子; 一个非零方向甚至可能与自己 $A$-正交.
\end{solution}
\begin{exercise}
设 $A$ SPD 且只有 $s$ 个不同特征值. 证明无论初值如何, 精确 CG 至迟 $s$ 步终止, 并说明为何这不等于浮点数中固定 $s$ 步达到任意精度.
\end{exercise}
\begin{solution}
取 $p(z)=\prod_{j=1}^s(1-z/\lambda_j)$, 则 $p(0)=1$, $p(A)=0$. 能量多项式最小化给出第 $s$ 步误差为零. 浮点迭代不严格保持这些多项式关系、残差正交性与共轭性; 数据和运算还会扰动谱, 所以实际需以真实残差和应用误差验收.
\end{solution}
\begin{exercise}
取 $A=\diag(1,100)$, $M=\diag(1,10)$. 求对称预条件矩阵的特征值和条件数. 比较 $M=I$ 与 $M=A$ 的情况, 并解释为什么仅比较迭代数不足以决定预条件器优劣.
\end{exercise}
\begin{solution}
$M^{-1/2}AM^{-1/2}=\diag(1,10)$, 条件数为 $10$. $M=I$ 时为 $100$, $M=A$ 时为 $1$, 后者一步终止. 但一般矩阵下精确应用 $A^{-1}$ 可能比许多便宜迭代还昂贵, 因此要计入预条件器构造和每次求解成本. 这个对角例子中三者都很便宜, 不能据此推广到复杂稀疏问题.
\end{solution}
\begin{exercise}
证明 $M^{-1}A$ 与 $M^{-1/2}AM^{-1/2}$ 相似, 并说明为什么前者的欧氏二范数条件数不一定等于后者的特征值比.
\end{exercise}
\begin{solution}
$M^{-1}A=M^{-1/2}(M^{-1/2}AM^{-1/2})M^{1/2}$. 相似保持特征值, 却只有酉相似才保持奇异值和欧氏二范数. 因此 $M^{-1}A$ 的欧氏条件数可能受 $M^{1/2}$ 的尺度影响; PCG 理论中的条件数是对称变换矩阵的正特征值比.
\end{solution}

% SOURCE: MIT 18.335J Spring 2019, Lecture 24 sparse matrix handout
% (Per-Olof Persson slides dated November 28, 2006), and Week 8 L23 summary.
% ADAPTATION: self-authored connected explanation, graph arguments, cost derivations and examples.
\originalchapter{稀疏存储、填充与直接法的复杂度}\label{ch:sparse}
稀疏矩阵并不只是一个有很多零的稠密矩阵. 我们必须同时改变存储、运算和消元顺序, 才能真正从零结构中获益.

\originalsection{稀疏结构与存储方式}
若一个 $n\times n$ 矩阵只有 $s$ 个非零元, 用稠密数组仍需 $n^2$ 个数的位置. 稀疏格式记录非零数值及其索引, 使存储约为 $O(s+n)$, 矩阵向量积约为 $O(s+n)$. 这里的 $n$ 包含输出向量初始化和行列指针, 不应在矩阵可能有大量空行时被省略.

结构矩阵与稀疏矩阵是不同概念. Toeplitz 矩阵可能几乎没有零, 但可用少量参数描述并快速相乘; 有些算子甚至不显式存储任何矩阵. Krylov 方法只需要快速的矩阵向量积, 因而适用范围比显式稀疏矩阵更广.

坐标格式 COO 存三个数组 $(i_\ell,j_\ell,a_\ell)$, 表示各项所在行、列和数值. 它便于组装, 特别是有限元计算中不同单元会向同一个矩阵位置贡献数值. 组装完成后应合并重复索引, 并按实际策略删除精确零; 否则存储中的条目数与数学上的非零元数不一致.

\begin{definition}
压缩行格式 CSR 用三个数组表示矩阵: 数值数组、列索引数组, 以及长度 $n+1$ 的行指针数组. 第 $i$ 行的所有条目位于第 $i$ 和第 $i+1$ 个行指针之间. 压缩列格式 CSC 交换行列角色, 按列保存数值与行索引.
\end{definition}
为避免编程语言下标差异, 下面例子用从 $1$ 开始的数值与索引数组, 并规定最后一个指针等于 $s+1$.

\begin{example}
把
\[
A=\begin{bmatrix}
4&0&-1&0\\
0&3&0&2\\
-1&0&5&0\\
0&2&0&6
\end{bmatrix}
\]
写成 CSR, 并计算其一次矩阵向量积的组织方式.
\end{example}
\begin{solution}
按每行列号递增排列, 数值数组为 $(4,-1,3,2,-1,5,2,6)$, 列索引为 $(1,3,2,4,1,3,2,4)$, 行指针为 $(1,3,5,7,9)$. 给定 $x$, 第 $i$ 个输出为该行对应条目的 $a_\ell x_{j_\ell}$ 之和. 此例共用 $8$ 次乘法, 每行两项合并需一次加法. 这些条目逐行连续读取, 不用扫描 $16$ 个稠密位置.
\end{solution}
CSR 适合逐行矩阵向量积, CSC 适合按列访问和许多稀疏分解. 这不意味着另一格式无法完成同一个数学运算, 而是局部访问、输出写入和库实现的代价不同. 若行内索引已经排序, 查找一个指定位置可二分搜索, 但插入新条目可能需要搬移后面的整个数组. 因此应区分\emph{组装格式}和\emph{计算格式}, 不宜在每次更新时都向压缩数组中插入一个元素.

稀疏矩阵相加或相乘也可能失去稀疏性. 若 $A_{ik}\ne0$ 且 $B_{kj}\ne0$, 则乘积的 $(i,j)$ 项有潜在贡献. 许多不同的 $k$ 会汇聚到同一个位置, 也可能恰好相消. 运算复杂度依赖这些结构交集和输出非零数, 不能只用输入各自的非零数简单相加来界定矩阵乘法.

\originalsection{消元为何产生填充}
对 Hermitian 正定矩阵, 先考虑消去一个主元的分块形式
\[
A=\begin{bmatrix}a&c^*\\c&B\end{bmatrix},\qquad a>0.
\]
Cholesky 或对称消元后, 其余变量面对的 Schur 补为
\begin{equation}\label{eq:sparse-schur}
S=B-\frac{cc^*}{a},\qquad S_{ij}=B_{ij}-\frac{c_i\overline{c_j}}a.
\end{equation}
若 $c_i,c_j\ne0$ 而 $B_{ij}=0$, 新值就是非零的 $-c_i\overline{c_j}/a$. 原来没有直接耦合的两个变量, 因为共同连接被消去的变量而出现了耦合. 这叫\emph{填充}.

\begin{definition}
对对称非零结构, 建立无向图 $G(A)$: 每个变量为一个顶点, 不同顶点 $i,j$ 间有边当且仅当结构上允许 $A_{ij}\ne0$. 消去一个顶点时, 把其尚未消去的邻居连接成团, 再删除该顶点. 新增边称为符号填充.
\end{definition}
式 \eqref{eq:sparse-schur} 就是这个图规则的代数依据. 若某条边原来已有数值, Schur 补更新可能出现精确相消; 符号分析通常按潜在非零处理, 忽略特殊数值相消, 因而给出预计结构或上界. 对一般数值, 这些潜在填充通常真的出现. 不能把每一条符号边无条件说成浮点输出中的严格非零.

\begin{example}\label{ex:sparse-star}
取
\[
A=\begin{bmatrix}
4&-1&-1&-1\\
-1&2&0&0\\
-1&0&2&0\\
-1&0&0&2
\end{bmatrix}.
\]
比较先消去中心变量 $1$, 与先消去叶子 $2,3,4$ 的填充.
\end{example}
\begin{solution}
叶子块为 $2I$, 中心对叶子块的 Schur 补为 $4-3/2=5/2>0$, 所以 $A$ SPD. 若先消去中心, 剩余矩阵为
\[
2I-\frac14\begin{bmatrix}1&1&1\\1&1&1\\1&1&1\end{bmatrix},
\]
三个叶子之间都出现 $-1/4$, 形成完整三角形. Cholesky 下三角因子一般有 $4+3+3=10$ 个非零条目, 即完整下三角.

若先消去任一叶子, 它只有一个剩余邻居, 不会产生邻居之间的新边. 三个叶子依次消去, 只把中心对角依次减少 $1/2$, 最后为 $5/2$. 下三角因子只需 $4$ 个对角条目和 $3$ 条星形边, 共 $7$ 个非零条目. 矩阵的数值和特征值没有改变, 只是变量顺序改变, 存储和后续更新就不同.
\end{solution}
若 $P$ 是置换矩阵, 对 $PAP^*$ 求分解并用相应顺序解方程, 得到的原解仍相同. 对 SPD 矩阵, 这样的对称置换保持正定性, 可以在数值分解前专门选择顺序以减少填充.

\originalsection{带宽、最小度与分隔子排序}
半带宽 $w$ 意味着 $A_{ij}=0$ 在 $|i-j|>w$ 时成立. 对 SPD 带状矩阵, 按自然顺序做 Cholesky 不会产生带外填充. 原因是消去第 $j$ 个变量时, 尚存邻居都在 $j+1,\ldots,j+w$ 中, 它们之间的距离不超过 $w$. 归纳即可保持带宽.

每列最多有 $w$ 个下三角非零元. 一次主元消去的稠密外积更新最多需要 $O(w^2)$ 工作, 所以带状 Cholesky 的上界为
\[
\text{存储 }O(nw),\qquad\text{分解 }O(nw^2),\qquad
\text{每个右端回代 }O(nw).
\]
三对角矩阵 $w=1$ 因此能在线性时间和线性存储中直接求解. 这正是前面 Lanczos 小三对角问题很便宜的一个原因.

带宽压缩方法如反向 Cuthill--McKee 试图把相邻顶点排到接近位置, 对长而窄的网格往往有效. 但较小带宽不等价于最小填充: 填充受局部邻居团结构和整个消元树影响, 不是只看最远的一条边.

最小度方法每一步消去当前剩余图中度数较小的顶点. 一个度数 $d$ 的顶点最多可能在其邻居间引入 $d(d-1)/2$ 条边, 因而低度数是合理的局部启发. 但当前度数会随填充改变, 而局部最小并不保证全局最佳. 实际的近似最小度方法通过压缩表示和度数估计降低维护图的成本.

\begin{definition}
图的分隔子是一个顶点集合 $S$, 删除 $S$ 后图分成彼此不相连的多个部分. 嵌套分割排序先递归排列这些部分, 再把分隔子 $S$ 的变量排在最后.
\end{definition}
把分隔子放在最后, 就能先分别消去子区域, 不让一个区域的内部消元直接把另一个区域的全部内部变量连接起来. 消元最终汇集到较小的界面. 对二维规则网格, 一条网格线可把区域近乎平分, 分隔子大小约为未知数总量的平方根; 对三维网格, 分隔面大小约为总量的 $2/3$ 次方. 二维和三维稀疏直接法的规模差异, 正是由这些界面大小引起.

求最优填充顺序本身是困难的组合优化问题. 实际排序常把嵌套分割与近似最小度混合使用. 不同图、机器内存、并行要求和数值主元策略下, 没有一个无需检验的万能顺序.

\originalsection{二维与三维网格的复杂度推导}
现在把 $N$ 用作未知数总数, 把 $m$ 用作每个方向的网格点数, 避免把二者混在一起. $d$ 维 $m\times\cdots\times m$ 网格有 $N=m^d$ 个未知数, 最近邻差分矩阵有 $O(N)$ 个非零元.

二维自然逐行编号的半带宽约为 $m=\sqrt N$. 带状估计给出存储 $O(N^{3/2})$, 分解 $O(N^2)$. 即使输入只有 $O(N)$ 个非零元, 填充仍可能比输入大得多. 对规则二维 Poisson 网格, 这些量级确实描述自然顺序的主要成本. 嵌套分割能显著改善它们.

为了精确说明估计的适用范围, 下述结论采用规则网格或具有相同比例分隔子的网格模型. 每次近乎平分区域, 在含 $N_s$ 个未知数的区域中, 假设其分隔子以及消元时的相关前沿大小均为 $O(N_s^{(d-1)/d})$. 最长边依次分割的规则网格满足这一量级. 消元前沿按稠密矩阵存储和更新, 一个大小 $s$ 的前沿需要 $O(s^2)$ 存储和 $O(s^3)$ 工作. 这些是本次分析的结构假设, 不是任意稀疏图自动拥有的性质.

\begin{proposition}[网格模型下的嵌套分割成本]\label{prop:nested-dissection-cost}
在上述模型中, 嵌套分割稀疏 Cholesky 的分解成本和因子存储量满足
\[
\begin{array}{c|ccc}
&d=1&d=2&d=3\\\hline
\text{因子存储}&O(N)&O(N\log N)&O(N^{4/3})\\
\text{分解工作}&O(N)&O(N^{3/2})&O(N^2)
\end{array}
\]
每个右端的前代、回代成本与因子非零数同阶.
\end{proposition}
\begin{proof}
在二叉递归的第 $\ell$ 层, 有 $O(2^\ell)$ 个区域, 每个含 $O(N/2^\ell)$ 个变量. 记分隔子大小
$s_\ell=O((N/2^\ell)^{(d-1)/d})$. 层存储上界为 $O(2^\ell s_\ell^2)$, 层工作为 $O(2^\ell s_\ell^3)$, 直到区域大小降为常数, 共 $O(\log N)$ 层.

当 $d=2$, 每层存储为 $O(N)$, 所以总量 $O(N\log N)$. 层工作为 $O(N^{3/2}2^{-\ell/2})$, 对几何级数求和得到 $O(N^{3/2})$.

当 $d=3$, 层存储为 $O(N^{4/3}2^{-\ell/3})$, 层工作为 $O(N^2 2^{-\ell})$, 两个级数分别为 $O(N^{4/3})$ 和 $O(N^2)$.

当 $d=1$, 分隔子和前沿大小为常数, 每层成本为 $O(2^\ell)$, 直到 $2^\ell=O(N)$; 总和为 $O(N)$. 最后还需保存和处理叶子区域, 同样为 $O(N)$. 三角求解沿因子条目作乘加, 因而成本为 $O(\operatorname{nnz}(L))$.
\end{proof}
对于二维模型, 最顶层界面虽然只有 $O(\sqrt N)$ 个未知数, 它的稠密分解已经需 $O(N^{3/2})$; 对三维, 顶层 $O(N^{2/3})$ 的界面需 $O(N^2)$. 这解释了为什么三维问题往往在较小的总未知数规模就遭遇直接法内存和时间瓶颈.

这些是模型量级, 不能用来预测每个具体常数. 多重前沿、超节点、消元树上的并行和稠密块 BLAS 可以显著改善实际速度, 却不消除几何界面导致的主要幂次. 非规则、细长、强连接或近乎稠密的图, 需要重新分析分隔子与前沿大小.

\originalsection{数值主元、符号分析与多右端复用}
稀疏直接求解一般分为三阶段. 第一阶段分析图结构, 选择排列, 预测因子结构、消元树和所需工作区; 第二阶段按当前数值执行分解; 第三阶段对一个或多个右端作三角求解. 若数值改变但非零结构不变, 符号阶段通常可以复用. 若矩阵不变、右端改变, 数值因子也可以复用, 这是直接法的重要优势.

SPD Cholesky 不需要为了避免小或负主元而作普通部分选主元, 所以符号结构较容易预先确定. 但浮点下接近奇异、缩放很差的 SPD 输入仍可能暴露数值问题, 不能把结构分析当成准确性证明.

一般非对称稀疏 LU 还要协调两个目标: 结构排列减少填充, 数值主元保证稳定. 可能先采用列排序减少预期填充, 再在分解中进行行主元选择或延迟消元. 数值主元变化会影响实际填充和前沿大小, 所以仅凭符号排序不能完全固定全部成本. 为了更稀疏而接受过小主元, 会增大消元乘子、元素增长和舍入误差; 为了更稳定而改变顺序, 可能增加填充. 二者必须共同判断.

若需要高精度残差或修正解, 可以利用已经得到的因子作迭代改进: 计算真实残差, 用因子近似求解校正方程, 再更新解. 它是否有效仍取决于因子误差、条件数、残差计算精度和舍入范围, 不是任意不稳定因子的自动补救.

\originalsection{直接法和迭代法怎样比较}
一个有用的比较不应只数原矩阵的非零元. 对直接法要估计\emph{填充后的因子}, 对迭代法要估计\emph{达到指定精度的迭代数及预条件成本}. 相同的 $N$ 和 $\operatorname{nnz}(A)$, 可以对应截然不同的计算难度.

以 $d$ 维 Poisson 方程的标准最近邻差分为例, 假设等距网格并采用齐次 Dirichlet 边界. 对每个方向 $m$ 个内部点, 不带 $h^{-2}$ 缩放的离散负 Laplace 特征值为
\[
\lambda_{k_1,\ldots,k_d}
=4\sum_{j=1}^d\sin^2\left(\frac{\pi k_j}{2(m+1)}\right),
\qquad 1\le k_j\le m.
\]
其谱公式可直接推导. 一维取向量 $v_i=\sin(\pi ki/(m+1))$, 边界 $v_0=v_{m+1}=0$. 余弦加法给 $2v_i-v_{i-1}-v_{i+1}=2(1-\cos(\pi k/(m+1)))v_i=4\sin^2(\pi k/(2(m+1)))v_i$. 不同 $k$ 给出一维全部正交特征方向. $d$ 维算子是各坐标的一维算子的 Kronecker 和, 各方向正弦向量的张量积因此是特征向量, 特征值为各一维特征值之和, 得上述公式.

最小值为 $\Theta(m^{-2})$, 最大值为 $\Theta(1)$, 所以 $\kappa_2(A)=\Theta(m^2)=\Theta(N^{2/d})$. 若乘以 $h^{-2}$, 两端同时缩放, 条件数不变. 因而未预条件 CG 的 Chebyshev 保证约需
$O(N^{1/d}\log(1/\tau))$ 步; 每步矩阵向量积为 $O(N)$, 总工作上界约为
\[
O(N^{1+1/d}\log(1/\tau)),\qquad\text{向量存储 }O(N).
\]
二维中这个幂次与嵌套分割分解的 $N^{3/2}$ 相近, 但精度、常数、存储和多右端复用不同. 三维未预条件 CG 的该界为 $N^{4/3}$ 量级, 小于直接分解的 $N^2$ 幂次, 可仍然随网格加密增加迭代数. 若有合适的网格无关多重网格预条件, 迭代数还可能保持有界, 总工作接近线性; 这依赖实际方程和预条件器的有效性, 不能只由稀疏性推出.

选择方法时, 可以按以下问题逐项判断:
\begin{enumerate}
\item 是否只有快速 $Av$, 却没有可访问的矩阵条目? 若是, 普通稀疏直接分解无从实施, Krylov 方法更自然.
\item 是否 SPD、Hermitian 不定或一般非对称? 这决定 CG、MINRES、GMRES 等方法的适用条件.
\item 因子填充能否装入内存? 三维界面、图中高度连接区域和主元变动可能使输入很小而因子很大.
\item 有多少个右端, 矩阵多久改变一次? 反复复用分解可能摊薄直接法的构造成本; 有效预条件器同样可以复用.
\item 所需精度多高, 系统有多病态? 迭代法可能后期停滞, 直接法也可能需要缩放、改进和残差验算.
\item 在并行机器上, 三角求解依赖和全局内积通信各自有多昂贵? 运算量相同不代表实际时间相同.
\end{enumerate}
较小稠密问题优先采用稳定的成熟稠密分解; 因子仍可容纳的稀疏问题常适合稀疏直接法; 最大规模或仅有矩阵向量积的问题常需要预条件迭代. 这些是由结构和资源支持的选择原则, 不应把某个固定维数阈值当作所有硬件和问题通用的定理. 对边界情形, 使用同一真实残差标准做小规模试算最有说服力.

\originalsection{习题与解答}
\begin{exercise}
对长度 $N$ 的路径图, 比较沿路径自然顺序消去, 与先消去一个内部顶点的第一步. 为什么三对角直接法本来就有线性复杂度, 嵌套分割不会改变其主要幂次?
\end{exercise}
\begin{solution}
自然顺序每次消去路径端点, 只剩一个邻居, 没有新边. 先消去内部顶点会连接其两个邻居, 产生一条填充边, 但剩余图仍是路径. 无论适当的递归顺序如何组织, 一维分隔子和前沿都是常数大小, 总成本为 $O(N)$. 排序可能改变常数和并行性, 不能把已是线性的工作降为次线性, 因为仅读写输入输出就需线性成本.
\end{solution}
\begin{exercise}
一个 $N$ 阶 SPD 矩阵有半带宽 $w=N^{1/3}$. 给出带状 Cholesky 的存储和分解上界. 仅凭这个半带宽, 能否断言嵌套分割一定获得二维模型的成本?
\end{exercise}
\begin{solution}
存储为 $O(N^{4/3})$, 分解为 $O(N^{5/3})$. 半带宽只说明当前编号下边的跨度, 不保证存在二维网格式 $O(\sqrt N)$ 分隔子, 也不控制递归区域的前沿. 所以不能直接使用二维嵌套分割公式.
\end{solution}
\begin{exercise}
设两个互不相连的子图通过分隔子 $S$ 连接. 证明先消去两个子图内部时, 一个子图的消元不会在另一个子图的内部顶点之间直接引入由前者产生的新边.
\end{exercise}
\begin{solution}
两个子图的内部之间原来没有边. 前者内部某顶点的剩余邻居只能位于同一内部或 $S$. 把这些邻居连接成团后, 仍没有后者内部顶点参与. 归纳所有内部消元步骤即可. 但 $S$ 上会积累更新, 最终处理 $S$ 时又会产生稠密界面耦合, 这正是分隔子成本必须计入的原因.
\end{solution}
\begin{exercise}
一个稀疏矩阵固定不变, 要求解 $1000$ 个不同右端. 为什么一次直接分解可能优于逐个独立运行 Krylov 方法? 又为什么这个结论不能仅凭右端个数就确定?
\end{exercise}
\begin{solution}
直接法只分解一次, 后续每个右端只做两次三角求解. 独立 Krylov 求解重复生成空间和执行内积, 成本可能很大. 但若因子填充过多装不下, 直接法不可行; 若预条件迭代很快, 或能复用子空间、使用块方法, 迭代总成本也可能更低. 需要比较构造成本、每右端成本、内存与目标精度.
\end{solution}
\begin{exercise}
二维规则网格的每层嵌套分割存储为何都是 $O(N)$, 而每层分解工作却按几何级数递减? 用此解释 $O(N\log N)$ 与 $O(N^{3/2})$ 的不同来源.
\end{exercise}
\begin{solution}
第 $\ell$ 层有 $2^\ell$ 个区域, 分隔子大小约 $(N/2^\ell)^{1/2}$. 平方后乘区域数为 $N$, 每层都贡献同阶存储, 层数 $O(\log N)$. 立方后乘区域数为 $N^{3/2}2^{-\ell/2}$, 越深层越便宜, 所以工作求和由顶层主导, 不再多一个对数因子.
\end{solution}

% SOURCE: course use of Julia; self-authored experiments. Official Julia documentation API checked 2026-10-08, no Julia executable in preparation environment.
\originalchapter{数值实验与综合训练}\label{ch:experiment}
\originalsection{先决定实验要检验什么}
数值实验最有价值的作用是区分几个经常混淆的命题: 分解残差小不等于基完全正交, 原方程残差小不等于解误差小, 特征值近似准确不等于特征向量稳定, 迭代步数少不等于总运行时间短. 一个好的实验围绕某个明确问题设计, 再用独立指标检验它.

本章给出可复现的数学实验与自编 Julia 示例. 代码采用实数 Float64 数据, 只调用 LinearAlgebra、SparseArrays 与 Random 标准库. API 名称按 Julia 官方文档核对; 编写环境未安装 Julia, 所以本章不把这些代码或性能称为本环境已运行的 Julia 实测. 随附脚本保留完整实现, 下文解释应观测哪些量、哪些现象有数学保证. 不复制原课程 Julia 简介中的第三方权利保留图像.

实验记录至少包括数据生成方式、浮点类型、矩阵尺寸、容差、软件与 BLAS 版本、线程设置. 若比较时间, 先运行一次使编译和初始化完成, 将数据生成排除在计时之外, 报告多次重复. 若比较准确度, 先固定同一个原始问题和同一个验收指标.

\originalsection{正交化实验}
取第 \ref{ch:qr} 章的近线性相关矩阵
\[
A=\begin{pmatrix}1&1&1\\\eps&0&0\\0&\eps&0\\0&0&\eps\end{pmatrix},\qquad \eps=10^{-8}.
\]
分别做 CGS、MGS 和 Householder QR, 记录
\[
\eta_{\rm fact}=\frac{\norm{A-QR}_F}{\norm{A}_F},\qquad
\eta_{\rm orth}=\norm{Q^*Q-I}_2.
\]
第一个量检验因子乘积是否近似原矩阵, 第二个检验基是否正交. 二者必须同时看. 某个 $Q$ 的列几乎平行时, 乘积 $QR$ 仍可能很接近 $A$; 若随后把 $Q^*$ 当作正交逆使用, 错误才会显著暴露.

用于比较 CGS 与 MGS 的教学实现如下. 它没有重正交化, 也不是供生产环境替代成熟库的代码.
\begin{verbatim}
function gs_demo(A; modified=true)
    m, n = size(A)
    Q = zeros(Float64, m, n)
    R = zeros(Float64, n, n)
    for j in 1:n
        v = copy(A[:, j])
        for i in 1:j-1
            R[i,j] = dot(Q[:,i], modified ? v : A[:,j])
            v -= R[i,j] * Q[:,i]
        end
        R[j,j] = norm(v)
        R[j,j] > 0 || error("zero remainder")
        Q[:,j] = v / R[j,j]
    end
    return Q, R
end
\end{verbatim}
变动 $\eps$ 可观察从清晰独立到接近数值秩亏的过程. 当 $\eps$ 不太小时, 两种 Gram--Schmidt 都可能表现良好; 当 $\eps$ 小到有效余量接近舍入尺度时, CGS 的正交性可显著丢失. MGS 往往改善, 但不能据几次测试证明它在所有输入上保持机器精度正交. 再正交化与 Householder 的差异, 应用同样两个指标比较.

还可解同一 $A$ 的最小二乘, 比较用正规方程、QR 与 SVD 产生的 $x$. 检验不仅是 $\norm{Ax-b}$, 还包括 $\norm{A^*(b-Ax)}$ 和相对于高精度或已知真解的误差. 对不相容问题, 残差本来可能很大; 正交最优性才是拟合是否达到最小值的检验.

\originalsection{残差与病态方向实验}
令 $A=\diag(1,10^{-12})$, $b=(1,0)^T$, $\widehat x=(1,1)^T$. 这个实验无需任何求解器, 就能严格算出相对残差为 $10^{-12}$, 相对解误差为 $1$. 它是理解停止准则的基线. 然后在不同奇异方向上给 $b$ 加同样大小的噪声, 比较解的变化, 验证第 \ref{ch:condition} 章的最坏方向解释.

第二个实验可用指定奇异值构造矩阵. 先取两个正交矩阵 $U,V$, 设置 $A=U\diag(1,10^{-2},10^{-4},10^{-6})V^T$, 选已知 $x_*$, 令 $b=Ax_*$. 给 $b$ 沿最小左奇异向量加扰动 $\eps u_4$, 则解恰好变化 $10^6\eps v_4$. 这一步是精确数学结论, 与所用求解器无关. 后向稳定算法应在这个敏感性之外只引入小量级后向误差.

对数据有噪声的版本, 可比较截断 SVD 和 Tikhonov. 在绘图时把横轴设为正则化强度或截断秩, 同时画数据残差与真解误差. 两条曲线的最优位置不一定相同. 若没有真解, 不能把数据拟合最好直接当成去噪最好.

\begin{example}
设 $A=\diag(1,\sigma)$, 真实 $x_*=(1,1)^T$, 数据为 $b=(1,\sigma+\eta)^T$. 精确逆解的第二分量误差是多少? Tikhonov 第二分量的偏差与噪声项分别是什么?
\end{example}
\begin{solution}
逆解第二分量为 $1+\eta/\sigma$, 误差 $\eta/\sigma$. Tikhonov 给
\[
(x_\lambda)_2=\frac{\sigma(\sigma+\eta)}{\sigma^2+\lambda},\qquad
(x_\lambda)_2-1=-\frac{\lambda}{\sigma^2+\lambda}
+\frac{\sigma\eta}{\sigma^2+\lambda}.
\]
第一项是正则化偏差, 第二项是噪声传播. 增大 $\lambda$ 抑制第二项, 却通常增大第一项, 不能只看其中一项选择参数.
\end{solution}

\originalsection{幂法与 Rayleigh 商速度实验}
取 $A=\diag(4,2,1)$, $v_0=(1,1,1)^T/\sqrt3$. 精确幂法方向与 $(4^k,2^k,1)^T$ 相同. 因而
\[
\tan\theta_k=\frac{\sqrt{4^k+1}}{4^k},\qquad
\rho(v_k)=\frac{4\cdot16^k+2\cdot4^k+1}{16^k+4^k+1}.
\]
对足够大的 $k$, 夹角约以 $1/2$ 的因子收缩, Rayleigh 商误差约以 $1/4$ 的因子收缩. 在半对数图上, 后者斜率约为前者两倍. 浮点误差底到达后会停滞, 不能把最后几步的比值当作新的收敛阶.

实验记录三个量: 与已知主特征向量的夹角、Rayleigh 商误差、真实特征残差. 若只记录相邻特征值估计的差, 一个停滞在错误位置的算法也可能看起来“变化很小”. 再把矩阵改为 $\diag(4,3.99,1)$, 观察谱间隙变小怎样降低幂法速度. 对移位反迭代, 应记录每个移位线性系统的求解残差, 而不仅记录外层向量变化.

\originalsection{稀疏 Poisson、CG 与分解复用}
一维 Dirichlet Poisson 矩阵为 $A=\operatorname{tridiag}(-1,2,-1)$. 取 $x_*=(1,\ldots,1)^T$, $b=Ax_*$. 稀疏存储有 $3n-2$ 个非零元, 按自然顺序的 Cholesky 下三角因子有 $2n-1$ 个非零元. 这说明一个规模很大的矩阵也可能有线性成本的直接解法; 不能只凭维数就断言必须迭代.

可用如下初始化和库分解:
\begin{verbatim}
using LinearAlgebra, SparseArrays
n = 100
A = spdiagm(-1 => -ones(n-1), 0 => 2.0 .* ones(n),
             1 => -ones(n-1))
xtrue = ones(n)
b = A*xtrue
F = cholesky(Symmetric(A))
xdirect = F \ b
residual = norm(b-A*xdirect)
\end{verbatim}
Symmetric 只是声明程序应该使用相应的对称视图, 并不证明任意输入真的是对称或正定. 本例结构本身已经保证这些条件. 库也可能先作减少填充的置换, 所以实际因子非零数要按其采用的顺序记录, 不自动等同于自然顺序的计数. 对真实测量矩阵, 必须先检查数据; 把一个非对称矩阵包成 Symmetric 可能意味着选择某一半条目来定义了不同矩阵.

用自编 CG 和同一问题比较时, 记录每步真实残差、能量误差和总矩阵向量积次数. 为了教学核验每步重新计算真实残差, 会额外增加一次乘法; 计时应明确这个测量开销. 向量相对误差只有已知真解时可直接计算. 对能量误差, 可用 $\sqrt{(x_*-x)^TA(x_*-x)}$, 不把欧氏误差自动当作能量误差.

然后把一个右端改为许多右端, 复用同一 $F$ 解多个系统. 直接法的构造成本只支付一次, 迭代法的预条件器也应复用后公平比较. 对二维和三维网格, 还要记录因子非零数、排序和峰值内存; 输入 $\operatorname{nnz}(A)=O(N)$ 并不保证因子也线性.

\originalsection{方法选择的综合练习}
\begin{exercise}
一个满列秩 $m\times n$ 矩阵 ($m\gg n$) 条件数约为 $10^7$, 需要解不相容最小二乘. 某程序形成 $A^TA$, 用 Cholesky 得到很小的正规方程残差. 这是否足以验收? 你会怎样改进计算与检查?
\end{exercise}
\begin{solution}
不足. 正规矩阵条件数约为 $10^{14}$, 它在形成时已经引入误差; 很小的正规系统残差只检验了这个形成后的问题. 宜用 Householder QR 或必要时 SVD 直接处理原 $A$, 检查原数据的最小二乘残差、$A^T(b-Ax)$、分解后向误差以及相应条件性. 不相容问题的残差可以本来就大, 不能把它必须接近零作为验收条件.
\end{solution}
\begin{exercise}
两个可逆矩阵条件数都为 $1$, 一个是 $2I$, 另一个是第 \ref{ch:gmres} 章的循环置换矩阵. 对初始残差 $e_1$ 的 GMRES, 为什么一个一步收敛而另一个可能前 $n-1$ 步完全停滞?
\end{exercise}
\begin{solution}
$2I$ 的最小消去多项式为 $1-z/2$, 所以一步能消去残差. 循环置换的残差方向依次走到相互正交的坐标轴, 在 $k<n$ 时 $A\mathcal K_k\perp e_1$, 最小修正为零. 条件数描述输入解映射的最坏放大, 不描述满足 $p(0)=1$ 的低次多项式能否在整个谱和当前残差方向上变小.
\end{solution}
\begin{exercise}
某 SPD 问题使用每次不同的非对称 ILU 近似作为“PCG 预条件器”, 测试中偶尔收敛. 能否据此使用第 \ref{ch:cg} 章的 Chebyshev 界?
\end{exercise}
\begin{solution}
不能. 标准 PCG 界依赖固定 Hermitian 正定的 $M$ 和等价对称变换. 非对称且变化的 ILU 不满足这些条件, 成功算例不补足证明缺口. 应改用满足条件的对称正定预条件器及 PCG, 或使用适用于变化预条件的灵活 GMRES 等算法, 并重新检查真实残差与成本.
\end{solution}
\begin{exercise}
一个稀疏非对称系统, 全 GMRES 残差持续下降, 但 GMRES(20) 每个周期几乎无改进. 你会检查哪些信息, 可在不改变原方程的前提下考虑哪些改进?
\end{exercise}
\begin{solution}
先验证小问题递推残差与真实残差一致, 检查正交性、预条件应用和周期末更新. 若不是实现问题, 可能是短周期丢失了关键谱方向, 或算子很非正规. 可在内存允许下增大重启长度, 改善预条件, 保留有价值的近似不变子空间或采用成熟的增广重启实现. 单靠更紧的容差通常不恢复丢失的搜索空间. 对一般矩阵, 固定长度重启本来就没有无条件收敛保证.
\end{solution}
\begin{remark}
本章实验题与示例代码为自编补充. 课程原有 Julia 资料提供工具背景, 官方 Julia 文档只用于核对调用名称与基本对象. 本章严格区分可证明的精确算术预测、允许读者执行的代码和本环境已经核算的数学例题, 不虚构运行时间或代码执行记录.
\end{remark}

\appendix
\originalchapter{来源与覆盖}\label{app:sources}
\originalsection{核心手册与章节对应}
以下页数是官方 PDF 的实际总页数, 含原有题名和 MIT OCW 许可末页; 它们不是中文正文页数. 原文件题名里的 Lecture 编号可能来自 2006 年等旧版本, 本册课次按 2019 年资源索引对应, 不擅自改写原文件日期. 中文讲义按数学主题重组, 下表表示内容来源与改编位置, 不表示逐字对照翻译.

\begin{longtable}{p{.8cm}p{8.2cm}p{.8cm}p{3cm}}
\toprule
编号 & 原手册题名 & 页数 & 本册对应\\
\midrule
\endhead
N01 & Floating-Point Arithmetic, the IEEE Standard (Per-Olof Persson) & 10 & 第 \ref{ch:float} 章\\
N02 & Some myths about floating-point arithmetic (Steven G. Johnson) & 2 & 第 \ref{ch:float} 章\\
N03 & Notes on the accuracy of naive summation (S. G. Johnson, 2019) & 3 & 第 \ref{ch:float} 章\\
N04 & Backwards stability of recursive summation (Steven G. Johnson, 2012/2015) & 3 & 第 \ref{ch:float} 章\\
N05 & Notes on the equivalence of norms (Steven G. Johnson, 2012) & 3 & 第 \ref{ch:float} 章\\
N06 & Gram-Schmidt Orthogonalization (Per-Olof Persson, 2006) & 3 & 第 \ref{ch:qr} 章\\
N07 & Householder Reflectors and Givens Rotations (Per-Olof Persson, 2006) & 13 & 第 \ref{ch:qr} 章\\
N08 & Hessenberg/Tridiagonal Reduction (Per-Olof Persson, 2006) & 2 & 第 \ref{ch:eigen} 章\\
N09 & The QR Algorithm I (Per-Olof Persson, 2006) & 4 & 第 \ref{ch:eigen} 章\\
N10 & The QR Algorithm II (Per-Olof Persson, 2006) & 9 & 第 \ref{ch:eigen} 章\\
N11 & Why restarting Arnoldi/Lanczos is not trivial (Steven G. Johnson, 2019) & 4 & 第 \ref{ch:krylov} 章\\
N12 & Sparse Matrix Algorithms (Per-Olof Persson, 2006) & 3 & 第 \ref{ch:sparse} 章\\
\midrule
合计 & 直接相关核心手册 & 59 & 12 份 PDF\\
\bottomrule
\end{longtable}
N01--N12 的官方链接在\href{https://ocw.mit.edu/courses/18-335j-introduction-to-numerical-methods-spring-2019/resources/lecture-notes/}{课程讲义资源栏}集中列出. 原件校验清单逐项记录了直链、文件名、页数、字节数与 SHA-256. 本册数学公式重新写成可编辑的原生 LaTeX, 没有用整页 PDF 截图代替正文.

\originalsection{性能资料与课堂主题}
第 \ref{ch:performance} 章另采用三份 Johnson 性能资料: Ideal-Cache Terminology (2 页), Performance Experiments with Matrix Multiplication (7 页), Experiments with Cache-Oblivious Matrix Multiplication (7 页), 共 16 页. 两份实验手册的硬件和编译器为历史环境, 本册保留其背景, 未把曲线当作新的实测结果.

本册的主课次映射为:
\begin{enumerate}
\item L2--4: 浮点、求和稳定性与范数, 对应第 \ref{ch:float} 章.
\item L5、L7: 条件数与扰动, 对应第 \ref{ch:condition} 章.
\item L7--8: SVD、回归、GSVD, 由 Alan Edelman 客座讲授, 对应第 \ref{ch:svd} 章.
\item L9--10: 正交化、Householder/Givens QR, 对应第 \ref{ch:qr} 章.
\item L11--12: 缓存与数据移动, 对应第 \ref{ch:performance} 章.
\item L13--14: LU、Cholesky、结构求解器, 对应第 \ref{ch:direct} 章.
\item L14--17: Schur、Hessenberg、幂法与 QR, 对应第 \ref{ch:eigen} 章.
\item L18--19: Arnoldi、Lanczos、Ritz 与重启, 对应第 \ref{ch:krylov} 章.
\item L20: GMRES 与收敛, 对应第 \ref{ch:gmres} 章.
\item L21--23: 预条件、CG、BiCG, 对应第 \ref{ch:cg} 章.
\item L24: 稀疏直接法, 对应第 \ref{ch:sparse} 章.
\end{enumerate}
上述主题均有官方逐周 HTML 说明, 但 L5、L7、L9、L13--14、L16--18、L20、L22--23 等课次没有覆盖整讲的独立手册. 一些课次列出补充阅读或往期 Persson 手册, 仍不等于完整逐讲讲稿. 本册因此补写了条件数定义与可达后向扰动、SVD 存在性与低秩逼近、LU/Cholesky 证明、Krylov 结构、GMRES 多项式论证、CG/PCG 推导和网格稀疏复杂度等内容. 这些证明与例题属于本册自编补充, 不标为 MIT 原文逐字译稿.

第 \ref{ch:experiment} 章的综合实验与 Julia 示例也是自编补充. 原课程 Julia 简介和速查表属于工具背景, 不是数值代数的数学证明来源. API 核对链接列入书目, 但未在编写环境运行 Julia, 因而没有声称其代码已做实机性能测试.

\originalsection{范围与许可边界}
官方讲义栏的 27 份 PDF 共 261 页, 涵盖整个 Introduction to Numerical Methods 的多个主题. 本册未纳入 L1 的求根、L6 的 ODE 以及 L25 以后优化、积分、Chebyshev 逼近和 FFT 的非数值代数全文, 以免与其他学科讲义重复. 原课提供的过往考试、问题集解答及外部版权教材也没有被混算为本册已经完整翻译的讲义体量.

主要改编材料依据 MIT OCW 默认 CC BY-NC-SA 4.0 许可, 仍需保留教师及手册作者、课程、题名、特殊署名、来源和改编说明. 课程页面明确提醒外部网站有独立条款, 不能因 MIT 链接了它们便默认继承 OCW 许可. 本册不下载或复制未获许可的 Trefethen--Bau 教材全文. Julia 简介中的 Project Jupyter 权利保留图像与优化概览中的 Springer 权利保留图像未进入本册改编或默认原件上传包.

本册的理论证明以所写明的精确算术和浮点模型为准. 成熟库的完整逐指令稳定性常数、Wilkinson 移位的专门全局收敛证明、多重网格与稳定化双共轭算法的完整实现, 只在明确条件和范围内作说明或引用, 不冒称已在这一本基础讲义中全部证明. 数值实验与机械排版检查也不替代数学证明.

\begin{thebibliography}{9}
\bibitem{ocw} Steven G. Johnson. Introduction to Numerical Methods, 18.335J, Spring 2019. MIT OpenCourseWare. \url{https://ocw.mit.edu/courses/18-335j-introduction-to-numerical-methods-spring-2019/}.
\bibitem{index} MIT OpenCourseWare. Resource Index and Lecture Notes for 18.335J, Spring 2019. \url{https://ocw.mit.edu/courses/18-335j-introduction-to-numerical-methods-spring-2019/pages/resource-index/}; \url{https://ocw.mit.edu/courses/18-335j-introduction-to-numerical-methods-spring-2019/resources/lecture-notes/}.
\bibitem{persson} Per-Olof Persson. Floating-Point Arithmetic, the IEEE Standard; Gram-Schmidt Orthogonalization; Householder Reflectors and Givens Rotations; The QR Algorithm I and II; Sparse Matrix Algorithms. MIT 18.335J/6.337J handouts, 2006--2007, republished in the Spring 2019 course resources. 各份实际题名、署名和页数见项目来源清单.
\bibitem{tb} Lloyd N. Trefethen and David Bau III. Numerical Linear Algebra. SIAM, 1997. 仅列为原课程指定教材与进一步阅读, 本册未复制其全文.
\bibitem{julia} The Julia Project. Linear Algebra and Sparse Arrays standard-library documentation. \url{https://docs.julialang.org/en/v1/stdlib/LinearAlgebra/}; \url{https://docs.julialang.org/en/v1/stdlib/SparseArrays/}. API核对日期: 2026-10-08.
\bibitem{terms} MIT OpenCourseWare. Privacy and Terms of Use. \url{https://ocw.mit.edu/pages/privacy-and-terms-of-use/}. Creative Commons BY-NC-SA 4.0: \url{https://creativecommons.org/licenses/by-nc-sa/4.0/}.
\end{thebibliography}
\end{document}
